% Tusculanae Disputationes (tusculanae-disputationes) --- English translation
% Translated by Claude (Anthropic) from the Latin (Perseus).
% Style guide: STYLE.md.

\ciceroBook{1}

\ciceroSection{1} When at last I had been freed, wholly or in great part, from the labors of legal defense and the duties of the Senate, I turned myself back, Brutus — at your urging above all — to those studies which I had kept alive in my mind, set aside as the times required, and, after a long interval of neglect, now took up again; and since the system and discipline of all the arts that bear on the right way of living is contained in the pursuit of wisdom, which is called philosophy, I thought this ought to be illuminated by me in Latin letters — not because philosophy could not be grasped through Greek writers and Greek teachers, but because it has always been my judgment that our countrymen have either discovered everything more wisely on their own than the Greeks, or improved upon what they took from them, in those matters at least which they had judged worth the effort.

\ciceroSection{2} For in our manners and the institutions of our life, in our households and our domestic affairs, we surely both maintain order better and live more handsomely; and our public commonwealth our ancestors regulated, beyond doubt, with better institutions and laws. What shall I say of the art of war? In this our people prevailed greatly by their valor, and even more by their discipline. As for what they achieved by nature rather than by letters, that is to be compared neither with Greece nor with any nation whatever. For what gravity so great, what such steadiness, what greatness of spirit, integrity, good faith, what virtue so outstanding in every kind, has existed in any people that it could stand comparison with our ancestors?

\ciceroSection{3} In learning, and in every kind of letters, Greece surpassed us; and there it was easy to win against men who put up no fight. For while among the Greeks the oldest class of the learned is that of the poets — since Homer and Hesiod lived before Rome was founded, and Archilochus in the reign of Romulus — we took up poetry later. About five hundred and ten years after the founding of Rome, Livius produced a play, in the consulship of Gaius Claudius, son of Caecus, and Marcus Tuditanus, a year before the birth of Ennius. He was older than Plautus and Naevius. Late, then, were poets either recognized or received among us. There is, to be sure, a passage in the \textit{Origines} which says that guests at banquets used once to sing to the pipe of the virtues of famous men; yet that no honor belonged to this calling is shown by a speech of Cato's, in which he reproached Marcus Nobilior as if with a disgrace, that he had taken poets with him into his province — for that consul had taken Ennius, as we know, into Aetolia. The less honor, then, there was for poets, the smaller was the zeal for poetry; and yet, where men of great genius arose in that field,

\ciceroSection{4} they did not fall short of the Greeks' glory. Or do we suppose that, had it been counted to the credit of Fabius, a man of the highest nobility, that he painted, we too would not have had many a Polyclitus and Parrhasius? Honor nourishes the arts, and all men are fired to their pursuits by glory, while whatever is held in low esteem among any people always lies neglected. The Greeks held the highest learning to lie in the music of strings and voices; and so Epaminondas, in my judgment the foremost man of Greece, is said to have played the lyre beautifully, while Themistocles, some years earlier, when he declined the lyre at a banquet, was thought the less educated for it. Music, therefore, flourished in Greece; all men learned it, and one who did not know it was not reckoned fully accomplished in learning.

\ciceroSection{5} Geometry was held in the highest honor among them, and so nothing was more illustrious than the mathematicians; but we have fixed the limit of this art at its usefulness for measuring and reckoning. The orator, by contrast, we embraced at once — not learned at first, yet apt for speaking, and afterward learned as well. For Galba, Africanus, and Laelius are recorded to have been learned, and Cato, who preceded them in age, to have been studious; then Lepidus, Carbo, the Gracchi, and from there men so great down to our own day that little or nothing was conceded to the Greeks. Philosophy has lain neglected up to this present age and has had no light from Latin letters; it is for us to illuminate and rouse it, so that, if while busy we did our fellow citizens some good, we may do them good as well, if we can, in our leisure.

\ciceroSection{6} In this we must labor all the harder, because there are said to be many Latin books already, written without due care by those excellent men, to be sure, but men not learned enough. Now it can happen that a man judges rightly yet cannot express what he judges with polish; but for anyone to commit his thoughts to writing who can neither order them nor make them clear nor draw the reader on with any delight, is the mark of a man squandering both his leisure and his letters without restraint. And so they read their own books with their own circle, and no one touches them except those who want the same license of writing granted to themselves. Therefore, if by my industry I have brought any credit to oratory, with far greater zeal I shall open up the springs of philosophy, from which those streams too were flowing.

\ciceroSection{7} But just as Aristotle — a man of the highest genius, knowledge, and abundance — when he was stirred by the fame of the rhetorician Isocrates, began to teach the young men to speak as well, and to join wisdom with eloquence, so it is my resolve neither to lay down my old pursuit of speaking and yet to occupy myself in this greater and richer art. For I have always judged that philosophy perfect which could speak fully and elegantly on the greatest questions; and to this exercise I have given my efforts so zealously that I have even ventured to hold formal discussions in the manner of the Greeks. Just so, recently at my place at Tusculum, after your departure, when a number of my friends were with me, I tried what I could do in that line. For as I once used to declaim cases — which no one did longer than I — so this is now my old man's declamation. I would have someone propose a topic he wished to hear discussed; and on that, either sitting or walking about, I would argue.

\ciceroSection{8} And so I have gathered the discussions of five days, the "schools" as the Greeks call them, into the same number of books. It went in this way: when the man who wished to listen had stated his own view, I would then argue against it. For this, as you know, is the old Socratic method of discussing against another's opinion — since Socrates believed that in this way what was nearest the truth could most easily be found. But, that our discussions may be set out more conveniently, I shall present them as though the matter were being acted out, not as though it were being narrated. The opening, then, will arise in this way. — Death seems to me to be an evil.

\ciceroSection{9} — For those who have died, or for those who must die? — For both. — It is wretched, then, since it is an evil. — Certainly. — Therefore both those for whom it has already come about that they should die, and those for whom it is yet to come, are wretched. — So it seems to me. — No one, then, is not wretched. — No one at all. — And indeed, if you wish to be consistent, all who have been born or ever will be are not only wretched but always wretched. For if you called only those wretched who must die, you would exclude no one of the living — for all must die — yet there would still be an end of wretchedness in death. But since the dead too are wretched, we are born into everlasting wretchedness. For those who perished a hundred thousand years ago must be wretched — or rather all, whoever have been born. — So I judge, entirely. — Tell me, I beg you:

\ciceroSection{10} do those things frighten you — three-headed Cerberus down among the dead, the roar of Cocytus, the crossing of Acheron, Tantalus worn out with thirst, his chin touching the surface of the water? Or that which Sisyphus undergoes — \textit{he heaves the rock, sweating in the strain, and gains not a hair's breadth}? Perhaps also the inexorable judges, Minos and Rhadamanthus? Before them no Lucius Crassus will defend you, no Marcus Antonius; nor, since the case will be tried before Greek judges, will you be able to bring in Demosthenes; you will have to plead your own cause before a vast assembly, on your own behalf. These things, perhaps, you fear, and for that reason you think death an everlasting evil. — Do you think me so far gone in folly as to believe such things? — Do you not believe them, then? — Not in the least. — A bad business, by Hercules, you tell me. — Why so, I beg? — Because I could be eloquent if I were to argue against all that. For who could not, in a case of that kind? Or what trouble is it to refute these monstrosities of poets and painters?

\ciceroSection{11} — And yet whole books are full of philosophers arguing against those very things. — Absurdly, to be sure. For who is so witless as to be moved by them? — If, then, men are not wretched among the dead, there are not in fact any men at all among the dead. — So I judge, entirely. — Where, then, are those whom you call wretched, or what place do they inhabit? For if they exist, they cannot be nowhere. — But I think they are nowhere. — They do not exist, then, either? — Just so, in that way; and yet wretched for that very reason — because they do not exist.

\ciceroSection{12} — Now I would rather you feared Cerberus than spoke so carelessly. — How do you mean? — The same man whom you say does not exist, you say exists. Where is your sharpness? For when you say he is wretched, you then say that one who does not exist exists. — I am not so dull as to say that. — What do you say, then? — That Marcus Crassus, for instance, is wretched, who let go those fortunes of his at death; that Gnaeus Pompeius is wretched, who was robbed of such great glory; in short, that all are wretched who are deprived of this light. — You roll back to the same point. For they must exist, if they are wretched; yet you were just now denying that those who have died exist. If, then, they do not exist, they can be nothing; and so they are not wretched either. — Perhaps I do not even say what I feel; for that very thing — not to exist, when you have existed — I count the most wretched of all. — What?

\ciceroSection{13} more wretched than never to have existed at all? Then those not yet born are wretched already, because they do not exist; and we, if we are to be wretched after death, were wretched before we were born. But I do not remember being wretched before I was born; you, if you have a better memory, I should like to know whether you recall anything of yourself. — You jest, as if I said that those are wretched who have not been born, and not those who have died. — You say, then, that they exist. — On the contrary: it is because they do not exist, after they have existed, that they are wretched. — Do you not see that you are contradicting yourself? For what could be more contradictory than that one who does not exist should be not only wretched but anything at all? When you have gone out by the Porta Capena and see the tombs of Calatinus, of the Scipios, the Servilii, the Metelli, do you think those men wretched? — Since you press me on a word, hereafter I shall not say that they are wretched, but only "wretched," for that very reason — because they do not exist. — You do not say, then, "Marcus Crassus is wretched," but only "Marcus Crassus wretched"? — Just so, plainly. — As if it were not necessary that whatever you pronounce in that manner must either exist or not exist! Have you not so much as a smattering of logic?

\ciceroSection{14} For the first thing handed down is this: every proposition (for so it occurs to me at present to render \textit{[Greek: axioma]} — I shall use another term later, if I find a better one) — that, then, is a proposition, which is either true or false. So when you say "Marcus Crassus wretched," you are either saying "Crassus is wretched," so that it can be judged whether this is true or false, or you are saying nothing at all. — Come, I now grant that those who have died are not wretched, since you have wrung from me the admission that those who do not exist at all cannot even be wretched. But what of this: we who are alive, since we must die, are we not wretched? For what pleasantness can there be in life, when day and night one must reflect that at any moment one must die? — Do you grasp, then, how much evil you have cast off from the human condition?

\ciceroSection{15} — In what way? — Because, if death were wretched even for the dead, we would have in life a kind of infinite and everlasting evil; but now I see the finish line, and once we have run down to it, there is nothing further to dread. But you seem to me to be following the maxim of Epicharmus, a man as sharp and far from dull as a Sicilian would be. — Which maxim? For I do not know it. — I shall say it, if I can, in Latin. For you know that I am no more in the habit of speaking Greek in Latin discourse than Latin in Greek. — And rightly so. But what, after all, is that maxim of Epicharmus? — \textit{I have no wish to die, but being dead I count as nothing.} — Now I recognize the Greek. But since you have compelled me to grant that those who have died are not wretched, finish the work, if you can, and make me think it not even wretched that one must die.

\ciceroSection{16} "That, surely, is no trouble at all — but I am attempting something greater." How is it no trouble? And what, pray, is this greater thing? "Because, since after death there is nothing evil, not even death itself is an evil — death, to which the very next moment is the time after death, in which you grant there is nothing evil. So not even the necessity of dying is an evil; for it is only the necessity of arriving at that which we admit to be no evil." Develop that more fully, please. These thornier points compel me to concede before I am persuaded. But what is this greater thing you say you are attempting? "To teach you, if I can, not only that death is no evil, but that it is even a good." That I do not demand — yet I long to hear it. For even if you do not accomplish what you intend, you will still accomplish this: that death is no evil. But I shall not interrupt you; I would rather hear a continuous discourse. What if I put a question to you?

\ciceroSection{17} Will you not answer? "That would be arrogant of me — yet unless something requires it, I would rather you did not ask." I shall humor you, and explain what you wish as best I can — not, however, like the Pythian Apollo, so that what I say must stand as certain and fixed, but as one ordinary man among many, following what is probable by conjecture. For I have no way of going further than to see what resembles the truth. Those will speak with certainty who claim both that these things can be grasped and that they themselves are wise. As you see fit; we are ready to listen.

\ciceroSection{18} Death itself, then, which seems the most familiar of things — what it is, must first be examined. For there are those who think death is the departure of the soul from the body; there are those who hold that no departure takes place, but that soul and body perish together, the soul being extinguished within the body. Of those who think the soul departs, some hold that it is scattered at once, others that it endures for a long while, others that it endures forever. And as to what the soul itself is, or where, or whence — there is great disagreement. To some the heart itself seems to be the soul, from which men are called "heartless," "deranged," and "of one heart"; and that shrewd Nasica, twice consul, was called "Corculum" — Little Heart — and Aelius Sextus, "an outstandingly hearted man, the keen Aelius Sextus." Empedocles holds the soul to be blood suffused around the heart; to others a certain part of the brain seemed to hold the soul's command; to others neither the heart itself satisfies, nor the notion that some part of the brain is the soul,

\ciceroSection{19} but some have placed the soul's seat and dwelling in the heart, others in the brain. As for the soul, some call it breath — as most of our own people do; the very word declares it, for we speak of "drawing breath" and "breathing out," and of the "spirited" and the "well-spirited," and of acting "from the soul's mind"; and the soul itself \textit{(animus)} is named from breath \textit{(anima)}. To Zeno the Stoic the soul seems to be fire. But these views I have named — heart, brain, breath, fire — are commonly held; the rest are each held by single men. As, long before, the ancients held, and most recently Aristoxenus, a musician and likewise a philosopher: that the soul is a certain tension of the body itself, like what in song and on the lyre is called harmony \textit{[Greek: harmonia]} — so that from the nature and shape of the whole body various motions are stirred, as sounds are in a song.

\ciceroSection{20} Here he did not depart from his own craft, and yet he said something — which, whatever its worth, had long before been both stated and explained by Plato. Xenocrates denied that the soul has any shape or, so to speak, body, and said it was number, whose force, as had already seemed to Pythagoras, is supreme in nature. His teacher Plato fashioned a threefold soul, whose ruling part — that is, reason — he placed in the head as in a citadel; and he wished the two other parts to obey it, anger and desire, which he set apart in their own regions: he located anger in the breast, desire below the midriff.

\ciceroSection{21} Dicaearchus, however, in that dialogue which he sets out in three books as held at Corinth among learned men in debate, in the first book makes many speak; in the other two he brings on a certain Pherecrates, an old man of Phthia, whom he says was descended from Deucalion, arguing that there is no soul at all, that this is a wholly empty word, that creatures are called "animate" and "living" to no purpose, that there is no soul or breath in man or beast, and that the entire force by which we either act or perceive is spread evenly through all living bodies and is inseparable from the body — being in fact nothing at all — and that there is nothing but body, one and simple, so shaped that by the tempering of its nature it lives and perceives.

\ciceroSection{22} Aristotle — far surpassing all others, Plato I always except, in genius and in diligence — having taken in those four well-known kinds of first principles from which all things arise, holds that there is a certain fifth nature, from which mind comes. For to think, to foresee, to learn, to teach, to discover something, and to remember so many other things, to love and hate, to desire and fear, to feel anguish and joy — these and their like he thinks reside in none of those four kinds. He brings in a fifth kind that has no name, and so calls the soul itself, by a new name, \textit{[Greek: endelecheia]} — a kind of continuous and perpetual motion. Unless some perhaps escape me, these are roughly the views about the soul. For Democritus — a great man, to be sure, but one who makes the soul out of smooth and round little bodies by some chance collision — let us leave aside; for there is nothing among those people that the swarm of atoms does not produce.

\ciceroSection{23} Which of these views is true, some god may see; which is likeliest to the truth is a great question. Do we prefer, then, to decide between these views, or to return to our subject? I should wish for both, if it were possible, but it is hard to combine them. So if we can be freed from the fear of death without discussing those questions, let us do that; but if it cannot be done unless this question of the soul is worked out, then let us take this now, if you think best, and the other another time. What I understand you to prefer — that, I think, is the more convenient course; for reason will bring it about that, whichever of the views I have set out is true, death is either no evil, or rather a good.

\ciceroSection{24} For if the soul is heart, or blood, or brain, then surely, since it is body, it will perish along with the rest of the body; if it is breath, perhaps it will be scattered; if fire, it will be extinguished; if it is Aristoxenus's harmony, it will be dissolved. What am I to say of Dicaearchus, who says the soul is nothing at all? On all these views nothing after death can concern anyone; for sensation is lost together with life, and to one who does not perceive there is nothing that can matter in any direction. The remaining views bring hope — if this perhaps delights you — that souls, when they have left their bodies, may reach heaven as their own dwelling. It does indeed delight me; and I should wish first that it be so, and then, even if it should not be so, that I might still be persuaded of it. What need, then, have you of our help? Can we surpass Plato in eloquence? Unroll carefully that book of his which is about the soul: you will find nothing more to desire. By Hercules, I have done so — and indeed more than once; but somehow, while I read, I assent; when I have laid down the book and begin to reflect by myself on the immortality of souls, all that assent slips away. What?

\ciceroSection{25} Do you grant this — that souls either remain after death or perish in death itself? I grant it. What if they remain? I concede that they are blessed. But if they perish? That they are not wretched, since they no longer even exist; for just now, under your compulsion, we conceded that very point. How then, or why, do you say death seems to you an evil? It will make us either blessed, if our souls remain, or not wretched, if we are without sensation.

\ciceroSection{26} Set it out, then, if it is no trouble — first, if you can, that souls remain after death; then, if you cannot establish that — for it is hard — you will teach that death is free of every evil. For this is the very thing I fear may be an evil: not, I say, to be without sensation, but to have to be without it. For the view you wish to maintain we can call on the best authorities — which in all cases both ought to count, and usually does count, for most — and first of all the whole of antiquity, which, the nearer it stood to its origin and to its divine descent, the better, perhaps, it discerned what was true. And so this one belief was rooted in those men of old, whom Ennius calls the "ancients": that in death there is sensation, and that a man is not so blotted out by his departure from life as to perish utterly;

\ciceroSection{27} and this may be understood from many other things, and above all from pontifical law and from the rites of the tomb, which men endowed with the greatest minds would not have observed with such care, nor sanctioned the violation of with so inexpiable a religious dread, had there not clung in their minds the conviction that death is not a destruction wiping out and effacing all things, but a kind of migration and change of life, which in famous men and women was wont to be a guide to heaven, while in the rest it was held to the ground and yet endured.

\ciceroSection{28} From this, and from the belief of our own people, "Romulus passes his days in heaven among the gods," as Ennius said, assenting to common report; and among the Greeks, and thence handed down to us and as far as the Ocean, Hercules is held a god so great and so present; from this Liber, born of Semele, and by the same celebrity of report the Tyndarid brothers, who are said to have been not only helpers of the Roman people's victory in battle but also messengers of it. What of Ino, daughter of Cadmus — is she not named "White Goddess" \textit{[Greek: Leukothea]} by the Greeks and held to be Matuta by us? What of this — to pursue no more — is not nearly the whole of heaven filled with the human race?

\ciceroSection{29} And if I were to search the ancient records and dig out from them what the writers of Greece have handed down, those very gods who are counted among the greater families would be found to have set out for heaven from here. Ask whose tombs are shown in Greece; recall, since you have been initiated, what is handed down in the mysteries: then at last you will understand how widely this extends. But those who had not yet learned the natural philosophy that men began to handle many years afterward had persuaded themselves of only as much as nature, prompting them, had taught; they did not grasp the reasons and causes of things, and were often moved by certain visions, those above all that came by night, so that men who had departed from life seemed to be alive.

\ciceroSection{30} Now, just as this seems the strongest proof to be brought forward for believing that the gods exist — that there is no race so savage, no one among all men so brutish, whose mind the belief in gods has not imbued (many hold corrupt views about the gods — for that is usually the work of vicious habit — yet all judge that there is a divine power and nature; and this is not produced by men's conversation or agreement, nor confirmed as an opinion by institutions or by laws; but in every matter the consensus of all nations is to be reckoned a law of nature) — who, then, does not first grieve at the death of his own precisely because he thinks them bereft of life's goods? Take away that belief, and you will have taken away grief. For no one mourns over his own misfortune; men feel pain, perhaps, and anguish, but that mournful lamentation and grieving weeping arises from this: that we judge the one we loved to be deprived of life's goods, and to feel it. And these things we feel by nature's guidance, with no reasoning and no instruction. But the strongest proof is that nature herself judges, in silence, concerning the immortality of souls — namely, that all men care, and care greatly, about what will be after death. "A man plants trees to benefit another age," as Statius says in the \textit{Synephebi} — with what in view, unless he thinks that later ages too concern him? Will the careful farmer, then, plant trees whose berry he himself will never lay eyes on, and a great man not plant laws, institutions, a commonwealth? What do the begetting of children, the propagation of one's name, the adoption of sons, the careful drawing of wills, and the very monuments and inscriptions of tombs signify, except that we think of the future as well? What?

\ciceroSection{32} As for this, surely you do not doubt that the standard of nature is rightly drawn from each best nature? What nature, then, is better in the human kind than that of those who hold themselves born to help, to protect, to preserve their fellow men? Hercules passed over to the gods: he would never have passed over, had he not, while he was among men, paved that road for himself. These are old examples by now, hallowed by the religion of all. But what are we to suppose was in the minds of so many great men in this republic of ours who gave their lives for the republic? That their name should be bounded by the same limits as their life? No one would ever offer himself to death for his country without great hope of immortality.

\ciceroSection{33} Themistocles might have lived at his ease; so might Epaminondas; so — not to seek old examples and foreign ones — might I. But somehow there clings in our minds a kind of foreboding of the ages to come, and this both arises most powerfully and shows itself most readily in the greatest natures and the loftiest spirits. Take that away, and who would be so mad as to live forever amid labors and dangers?

\ciceroSection{34} I speak of leading men. And what of the poets? Do they not wish to be honored after death? Whence, then, this: \textit{Behold, O citizens, the features of the likeness of old Ennius. He set in verse the greatest deeds of your fathers.} He demands the reward of glory from those whose fathers he had crowned with glory. And the same poet: \textit{Let no one honor me with tears.} Why? \textit{Alive, I fly upon the lips of men.} But why speak of poets? Even craftsmen wish to be honored after death. Why else did Phidias enclose a likeness of himself in the shield of Minerva, when he was not permitted to inscribe his name? And what of our philosophers? Do they not inscribe their own names in the very books they write on the contempt of glory?

\ciceroSection{35} But if the agreement of all is the voice of nature, and all men everywhere agree that there is something which concerns those who have departed from life, then we too must hold the same; and if we shall judge that those whose spirit excels in genius or in virtue, being of the best nature, discern most keenly the force of nature, then it is likely — since each best man serves posterity most — that there is something of which he will have awareness after death.

\ciceroSection{36} But just as we hold by nature that the gods exist, and learn by reason what they are like, so we judge by the agreement of all nations that souls endure, while it is by reason that we must learn in what abode they remain and what they are like. Ignorance of this invented the underworld and those terrors which you seemed, not without cause, to despise. For since bodies fall to the earth and are covered with soil — whence the word "to inhume" — men supposed that the remaining life of the dead is lived beneath the earth; and great errors followed upon that supposition, errors which the poets magnified.

\ciceroSection{37} For the crowded gathering of the theater, in which there are women and children, is stirred to hear so grand a strain: \textit{Here I am, and I come from Acheron, by a way scarcely passable, steep and hard, through caverns built of rough hanging rocks, the greatest, where the stiff, thick gloom of the dead stands fixed.} And the error had such force — though to me at least it now seems lifted — that, although they knew the bodies were burned, they still imagined things happening in the underworld which without bodies could neither happen nor be conceived. For they could not grasp with the mind souls living of themselves; they sought some shape and figure for them. From this comes Homer's whole \textit{Summoning of the Dead [Greek: nekyia]}; from this what my friend Appius used to perform, his \textit{Consultation of the Dead [Greek: nekyomanteia]}; from this, in our own neighborhood, Lake Avernus, \textit{whence the souls are roused, shrouded in dark shade, the images of the dead, from the deep mouth of Acheron, with salt blood.} And yet they will have these images speak — which can be done neither without a tongue nor without a palate nor without the force and form of throat, sides, and lungs. For they could see nothing with the mind; they referred everything to the eyes.

\ciceroSection{38} But it is the mark of a great mind to call thought away from the senses and to lead reflection away from habit. And so I do indeed believe that others said it too, in all those ages; but the first who said, as far as the record stands in writing, that the souls of men are everlasting was Pherecydes of Syros, an ancient man, to be sure — for he lived in the reign of my own clansman. His pupil Pythagoras confirmed this view above all; he came into Italy in the reign of Tarquin the Proud, and held all that Magna Graecia in his grip, both by the honor of his teaching and by his authority; and for many ages afterward the name of the Pythagoreans so flourished that no others seemed learned at all. But I return to the ancients. They scarcely ever gave a reason for their view, unless something had to be explained by numbers or by diagrams.

\ciceroSection{39} They say that Plato came into Italy to learn from the Pythagoreans, and learned all their teachings; and that on the eternity of souls he not only held the same view as Pythagoras but also brought forward a reason for it. That reason — unless you have something to say — let us pass over, and let go this whole hope of immortality. — What? After raising me to the highest expectation, you abandon me? By Hercules, I would rather err with Plato, whom I know how highly you value and whom I admire on your own report, than hold true views with those others.

\ciceroSection{40} Bravely said! For I myself would not be unwilling to err with that same man. Do we, then, doubt — ? Or is it as with most questions? Though here at least not in the slightest. For the mathematicians persuade us that the earth, set in the middle of the universe, occupies in relation to the embrace of the whole heaven something like the size of a point, which they call its \textit{center [Greek: kentron]}; and further, that such is the nature of the four body-producing elements that, as if they held their impulses portioned and divided among themselves, the earthy and the moist are borne by their own inclination and their own weight at right angles down into the earth and into the sea, while the remaining two parts — one fiery, the other airy — just as those former two are borne by heaviness and weight into the middle place of the world, so these in turn fly up by straight lines into the celestial region, whether because nature itself seeks the higher, or because the lighter are by nature repelled by the heavier. Since these things stand firm, it must be clear that souls, when they have departed the body — whether they be airy, that is, breathable, or fiery — are borne aloft.

\ciceroSection{41} But if the soul is some kind of number — which is said with more subtlety than clarity — or that fifth nature, unnamed rather than ununderstood, then souls are even far more unmixed and pure, so that they carry themselves up the farthest from the earth. The soul, then, is some one of these — lest a mind so vigorous lie sunk in the heart, or in the brain, or in the blood of Empedocles. Dicaearchus, however, along with his contemporary and fellow pupil Aristoxenus, both learned men to be sure, let us leave aside; of whom the one seems never even to have felt a pang, since he does not perceive that he has a soul, while the other so delights in his own melodies that he tries to carry them over into these questions as well. Now harmony we can recognize from the intervals of sounds, whose varied combination produces yet more harmonies; but how the arrangement of the limbs and the figure of a body devoid of soul can produce a harmony, I do not see. But this man, learned though he is, as indeed he is, may well leave these matters to his master Aristotle and teach singing himself; for that Greek proverb gives sound counsel: let each man practice the craft he knows.

\ciceroSection{42} But that chance concurrence of smooth round indivisible bodies let us cast out root and branch — which Democritus nonetheless will have to be warmed and breathable, that is, animate. The soul, however, which, if it belongs to these four kinds out of which all things are said to be composed, consists of inflamed air — as I see Panaetius is most inclined to think — must seek the upper regions. For these two kinds have nothing that tends downward, and always seek the heights. So whether they are dispersed, this happens far from the earth, or whether they endure and keep their condition, all the more must it be that they are borne to the sky, and that by them is broken through and divided this thick and condensed air which is nearest the earth. For the soul is warmer, or rather more burning, than this air which I just called thick and condensed; and this can be known from the fact that our bodies, made up of the earthy kind of elements, grow warm with the heat of the soul.

\ciceroSection{43} It follows that the soul escapes the more easily from this air, which I now invoke so often, and breaks through it, because nothing is swifter than the soul: there is no speed that can contend with the speed of the soul. And if the soul remains uncorrupted and true to itself, it must be borne in such a way that it penetrates and divides this whole sky, in which clouds and rains and winds are gathered, which is both moist and dark on account of the exhalations of the earth. When the soul has risen above this region and reached and recognized a nature like its own, it comes to rest among fires joined of thin air and the tempered burning of the sun, and ceases to lift itself higher. For when it has attained a lightness and warmth like its own, balanced as though by equal weights, it moves in no direction, and that at last is its natural seat, when it has penetrated to what is like itself; there, lacking nothing, it will be nourished and sustained by the same things by which the stars are sustained and nourished. And since we are accustomed to be set aflame by the torches of the body toward almost every desire, and to be the more inflamed because we strive against those who have what we long to have, surely we shall be blessed when, the body left behind, we shall be free of both desires and rivalries;

\ciceroSection{44} and what we do now, when our cares are loosened — that we wish to look at something and behold it — that we shall do then far more freely, and shall set ourselves wholly to contemplating and examining things, because there is by nature in our minds a certain insatiable desire to see the truth, and the very shores of those regions, once we have arrived, will give us an easier knowledge of the celestial things, and with it a greater desire of knowing.

\ciceroSection{45} For this beauty even on earth roused that ancestral and forefathers' philosophy, as Theophrastus says, kindled by the desire of knowledge. But those above all will enjoy it who, even then, while they dwelt on these lands wrapped in gloom, still longed to discern with the mind's edge. For if those think they have attained something who have seen the mouth of the Pontus and that strait through which passed the ship named \textit{Argo}, \textit{because in her the chosen men of the Argives were carried in their quest for the golden fleece of the ram}; or those who have seen those straits of Ocean, \textit{where the ravenous wave divides Europe from Libya} — what spectacle, after all, do we think it will be, when it will be granted to behold the whole earth, and its position, its form, its compass, and then both the habitable regions and again those left empty of all cultivation through the violence of cold or heat?

\ciceroSection{46} For not even now do we discern with the eyes the things we see; for there is no sense in the body, but — as not only the natural philosophers teach but the physicians too, who have seen these things laid open and exposed — there are, as it were, certain channels bored from the seat of the soul to the eyes, to the ears, to the nostrils. And so often, hindered either by thought or by some force of illness, we neither see nor hear, though the eyes and ears are open and unimpaired — so that it can easily be understood that it is the soul that both sees and hears, not those parts which are, as it were, the windows of the soul, by which nonetheless the mind can perceive nothing unless it is at work and present. What of the fact that with the same mind we grasp things utterly unlike — color, taste, warmth, smell, sound? These the soul would never recognize through its five messengers, were not all things referred to it, and were it not the sole judge of all. And surely those things will then be discerned far more pure and clear, when the soul, set free, has arrived where nature carries it.

\ciceroSection{47} For as things stand now, although nature has fashioned with the most cunning craftsmanship those channels that open from the body to the soul, still they are in a manner walled off by these earthen, compacted bodies of ours. But when there is nothing left except the soul, no obstacle thrown in its way will hinder it from perceiving the true nature of each thing. However copiously we might speak of this, if the matter required it — how many, how various, how vast the spectacles the soul would have before it in the regions of heaven.

\ciceroSection{48} Thinking on this, I am often given to wonder at the insolence of certain philosophers, who marvel at the knowledge of nature and pour out exultant thanks to its discoverer and founder, and revere him as a god; for they say they have been freed by him from the harshest of masters — from everlasting terror and from fear by day and by night. What terror? What fear? What old woman is so out of her wits as to dread the things you would no doubt have dreaded yourselves, had you not studied natural philosophy — \textit{the deep Acherontian halls of Orcus, the pale realms of death, the places clouded over with darkness}? Is a philosopher not ashamed to boast of this, that he does not fear such things and has learned them to be false? From which one may gather how sharp these men are by nature, since without instruction they would have believed it all.

\ciceroSection{49} A splendid prize, surely, they have won — to have learned that, when the time of death comes, they will perish entire. Suppose it is so — for I put up no fight — what is there in the matter either to rejoice in or to glory over? And yet nothing whatever occurs to me to show that the doctrine of Pythagoras and of Plato is not the true one. For even if Plato offered no reasoning at all — see how much I grant the man — by his very authority he would overcome me; but he has brought forward so many reasons that he plainly wished to convince the rest and certainly convinced himself.

\ciceroSection{50} But most men strive against this, and condemn souls to death as if to a capital sentence; and there is no other reason why the eternity of souls should seem incredible to them, except that they cannot conceive what the soul is like when free of the body, or grasp it in thought. As though they understood what it is like within the body itself — its conformation, its magnitude, its place; so that, if everything now hidden in a living man could already be seen, the question would be whether the soul would fall within view at all, or whether it is of such fineness that it escapes the keenest sight.

\ciceroSection{51} Let those who deny they can conceive the soul without a body weigh this: they will see how much they conceive of it within the body itself. For my own part, when I contemplate the nature of the soul, a far more difficult thought confronts me, far more obscure — what the soul is like within the body, as in another's house — than what it is like when it has gone out and come, as it were, into its own home, the free expanse of heaven. For if we cannot understand the nature of a thing we have never seen, surely we can embrace in thought both god himself and the divine soul released from the body. Dicaearchus, indeed, and Aristoxenus, because it was difficult to understand what the soul was, or of what kind, denied that there was any soul at all.

\ciceroSection{52} It is the greatest thing of all to see the soul by the soul itself; and surely this is the force of Apollo's precept, by which he counsels each man to know himself. For I do not believe he prescribes that we should know our limbs, our stature, or our shape; nor are we our bodies, nor, in saying this to you, am I speaking to your body. When therefore he says "know thyself," he says this: know your soul. For the body is, as it were, a vessel or some receptacle of the soul; whatever is done by your soul is done by you. To know it, then — unless this were a divine thing, the precept of so penetrating a mind would not have been ascribed to a god: namely, that one can come to know oneself.

\ciceroSection{53} But if the soul itself does not know what kind of thing the soul is, tell me, I ask — will it not even know that it exists, will it not even know that it moves? From this was born that argument of Plato's, which Socrates sets out in the \textit{Phaedrus} and which I placed in the sixth book of the \textit{Republic}: "What is forever in motion is eternal; but what imparts motion to another, and is itself driven from elsewhere, must, when its motion has an end, have an end of living. That alone, then, which moves itself, because it is never abandoned by itself, never even ceases to move; indeed, for all other things that move, it is the source, it is the first principle of motion.

\ciceroSection{54} But of a first principle there is no origin; for from a first principle all things arise, while the principle itself can be born of nothing else; for it would not be a first principle if it were begotten from elsewhere. And if it never arises, it never perishes either; for a first principle, once extinguished, will neither be reborn from another nor create another out of itself, since it is necessary that all things arise from a first principle. So it comes about that the beginning of motion is from that which moves itself; and this can neither be born nor die — or else the whole heaven must collapse and all nature halt, finding no force by which it might be set in motion from the first. Since it is therefore plain that whatever moves itself is eternal, who is there to deny that this nature has been granted to souls? For everything that is driven by an external impulse is without soul; but what is a living being is roused by motion inward and its own; for this is the proper nature and force of the soul. And if it is the one thing among all that forever

\ciceroSection{55} moves itself, then surely it was not born, and it is eternal." Let all the common philosophers come running together — for so they seem fit to be called, those who dissent from Plato and Socrates and that family — not only will they never explain anything so elegantly, but they will not even understand how subtly this very conclusion has been drawn. The soul, then, perceives that it moves; and when it perceives this, it perceives at the same time that it moves by its own force and not another's, and that it can never come to pass that it is abandoned by itself. From this its eternity follows — unless you have something to say against it. As for me, I would gladly let nothing even come into my mind against it: so much do I favor that view. — What of this?

\ciceroSection{56} Do you, after all, judge those things lighter which declare that there is something divine present in the souls of men? If I could discern how these could come to be, I should also see how they might perish. For blood, bile, phlegm, bones, sinews, veins — the whole shape, in short, of the limbs and of the entire body — I think I can say out of what they were compounded and how they were made. But the soul itself — if there were nothing in it but the means by which we live, I should suppose a man's life sustained by nature just as a vine's is, or a tree's; for of these too we say that they live. Likewise, if a man's soul had nothing but the power to desire or to flee, that too it would share with the beasts.

\ciceroSection{57} It has, first of all, memory — and that an infinite memory of innumerable things. This Plato will have it to be the recollection of an earlier life. For in that book entitled the \textit{Meno}, Socrates puts to a certain young boy some geometrical questions about the measurement of a square. To these the boy answers as a boy would, and yet the questions are so easy that, answering step by step, he arrives at the same point he would have reached had he learned geometry. From this Socrates will have it follow that learning is nothing other than recollection. This passage he explains far more carefully still in that discourse he held on the very day he departed from life; for he shows that anyone, however untaught in all things he may appear, by his answers to one who questions well makes it plain that he is not then learning those things but recognizing them by remembering — and that it could in no way come to pass that from boyhood we should hold so many notions of so great matters, implanted and as it were sealed upon our souls, which they call \textit{[Greek: ennoiai]}, unless the soul, before it entered the body, had been vigorous in the knowledge of things.

\ciceroSection{58} And since nothing truly is — as is argued by Plato in every passage — for he thinks that nothing is which comes to be and perishes, and that that alone is which is forever such as it is (which he calls \textit{[Greek: idea]}, and we a "form") — the soul, shut up in the body, could not recognize these things, but brought the knowledge of them already gained; and from this the wonder at the knowledge of so many things is taken away. Nor does the soul see them plainly when it has suddenly migrated into so unfamiliar and so disordered a dwelling; but when it has collected and refreshed itself, then it recognizes them by remembering.

\ciceroSection{59} So learning is nothing other than recollecting. As for me, I marvel at memory in a still greater way. For what is that by which we remember? What is its force, and what its origin? I am not asking how great a memory Simonides is said to have had, or Theodectes, or Cineas, who was sent as envoy by Pyrrhus to the Senate, or, lately, Charmadas, or, only recently, Metrodorus of Scepsis, or our own Hortensius. I speak of the common memory of mankind, and most of all of those engaged in some greater pursuit and art, the measure of whose minds it is hard to estimate, so much do they remember.

\ciceroSection{60} To what, then, does this discussion point? I think we must understand what that force is and from where it comes. Certainly it is not of the heart, nor of the blood, nor of the brain, nor of atoms; whether it is of air or of fire I do not know, and I am not ashamed, like those men, to confess that I do not know what I do not know. This much, if I could affirm anything at all about an obscure matter, I would swear to: whether the soul is air or fire, it is divine. For tell me, I beg you — does so great a force of memory seem to you to have been sown or compounded out of earth, under this misty and murky sky? If you do not see what this thing is, still you see what kind of thing it is; if not even that, still you surely see how great it is.

\ciceroSection{61} What then? Do we suppose there is some capacity in the soul, into which, as into some vessel, the things we remember are poured? That, surely, is absurd; for what bottom, or what such shape of the soul, can be conceived, or what capacity so vast at all? Or do we suppose the soul to be stamped, as wax is, and that memory consists in the traces of things imprinted upon the mind? But what traces can there be of words, what of the things themselves; and what magnitude, again, could be so immense as to figure forth so many things? What of this?

\ciceroSection{62} And that power, what in the end is it, which searches out hidden things — the power we call discovery and invention? Does it seem to you compounded of this earthly, mortal, perishable nature? Or the man who first imposed names on all things — which Pythagoras counted the mark of the highest wisdom? Or the man who gathered scattered human beings together and summoned them into a shared life? Or the man who bounded the sounds of the voice, which seemed infinite, within the few marks of letters? Or the man who charted the courses of the wandering stars, their advances and their standings still? All these were great men; greater still those before them — who discovered grain, who clothing, who shelter, who the ordered ways of living, who defenses against wild beasts: men by whom we were tamed and refined, and so passed on from the necessary crafts to the more elegant ones. For great delight has been provided to the ears as well, once the variety and the nature of sounds was discovered and tuned; and we have looked up at the stars — both those fixed in their settled places and those that wander, not in fact but in name — and the man who in his mind saw their revolutions and all their motions taught that his own mind was like the mind of the one who had fashioned them in the heavens.

\ciceroSection{63} For when Archimedes bound the motions of the moon, the sun, and the five wandering stars into a sphere, he did the same thing as that god of Plato who in the \textit{Timaeus} built the world — making a single revolution govern motions utterly unlike one another in slowness and in speed. And if this cannot happen in our world without a god, then not even in the sphere could Archimedes have imitated those same motions without divine genius.

\ciceroSection{64} To me, indeed, not even these more familiar and more obvious powers seem to lack a divine force — so much so that I cannot believe a poet pours out a weighty and full song without some heavenly prompting of the mind, or that eloquence flows in abundance, with resounding words and rich thoughts, without some greater power. And philosophy, the mother of all the arts — what is it but, as Plato says, a gift, or, as I say, a discovery of the gods? It was philosophy that first schooled us in the worship of the gods, then in the law of men, which rests in the fellowship of the human race, and then in restraint and greatness of soul; and it was philosophy that drove the fog away from the soul, as from the eyes, so that we might see all things — above and below, first and last and in between.

\ciceroSection{65} Surely this power that accomplishes things so many and so great seems to me divine. For what is memory of things and of words? What, further, is invention? Surely it is something greater than which nothing can be conceived even in a god. For I do not suppose the gods take their joy in ambrosia or nectar or in Youth pouring their cups, nor do I listen to Homer, who says that Ganymede was carried off by the gods for his beauty, to pour out drink for Jupiter — no just cause for so great a wrong to be done to Laomedon. Homer made this up, and transferred human things to the gods: I would rather he had transferred divine things to us. And what are the divine things? To be vigorous, to be wise, to discover, to remember. Therefore the soul that …, as I say, is divine, or, as Euripides dares to say, is a god. And indeed, if a god is either breath or fire, then so too is the soul of man; for as that heavenly nature is free of earth and free of moisture, so the human soul is without share in either of these. But if there is a certain fifth nature, first brought in by Aristotle, this belongs both to the gods and to souls. Following this view, I expressed it in these very words in my \textit{Consolation}:

\ciceroSection{66} ``No origin for souls can be found on earth; for there is nothing in souls that is mixed or compounded, nothing that seems born or formed out of earth, nothing even moist or airy or fiery. For in these natures there is nothing that possesses the power of memory, of mind, of thought — nothing that holds the past and foresees the future and can grasp the present. These powers alone are divine, and there will never be found any source from which they could come to man except from a god. Singular, then, is this nature and power of the soul, set apart from these common and familiar natures. And so, whatever it is that perceives, that knows, that lives, that thrives, must be heavenly and divine, and for that reason eternal. Nor indeed can the god himself, as we understand him, be understood in any other way than as a mind unbound and free, severed from all mortal compounding, perceiving and moving all things and itself endowed with everlasting motion. Of this kind, and from this same nature, is the human mind.'' Where, then, is that mind, and of what sort is it?

\ciceroSection{67} —Where is yours, and of what sort? Can you say? Or, if I do not have everything I would wish to have for understanding, am I not to be allowed by you to use even what I do have? The soul has not the power to see itself; but, like the eye, the soul, though it does not see itself, discerns other things. It does not see — and this is the least of it — its own form (though perhaps it sees that too; but let us set it aside); its force, its keenness, its memory, its motion, its speed it certainly does see. These are great, these are divine, these are everlasting; what its face may be, or where it dwells, is not even worth asking.

\ciceroSection{68} It is as when we behold first the beauty and brightness of the sky, then the speed of its revolution, greater than we can conceive, then the alternations of days and nights and the fourfold changes of the seasons, fitted to the ripening of crops and the tempering of bodies, and the sun, the moderator and leader of all these things; and the moon, marking and signaling the days, as in a calendar, by the waxing and waning of its light; then, in that same circle divided into twelve parts, the five stars borne along, holding most steadily to those same courses with motions unlike one another; and the form of the night sky, adorned on every side with constellations; then the globe of the earth rising up out of the sea, fixed in the central place of the whole universe, made habitable and cultivated in two regions far apart from each other — of which the one that we inhabit lies \textit{beneath the pole, toward the seven stars, from where the shuddering howl of the North Wind drives down the icy snows}, while the other, the southern, is unknown to us, and the Greeks call it

\ciceroSection{69} the \textit{[Greek: antichthona]}, the counter-earth — and the remaining regions lie untilled, because they are either stiff with cold or scorched with heat; whereas here, where we dwell, in its own season there is no failing of how \textit{the sky grows bright, the trees put forth their leaves, the gladdening vines swell with their tendrils, the branches bend low with the richness of their berries, the standing corn yields its grain, all things flower, the springs gush, the meadows are clothed with grass}; then the multitude of animals, some for food, some for the working of the fields, some for carriage, some for the clothing of our bodies; and man himself, set as it were to contemplate the heavens and to worship the gods,

\ciceroSection{70} with all the fields and seas obedient to man's use — when, then, we discern these things and others beyond number, can we doubt that some being presides over them: either a maker, if these things were born, as Plato holds, or, if they have always been, as Aristotle is pleased to maintain, a moderator of so great a work and gift? So too with the mind of man: though you do not see it, just as you do not see god, yet, as you recognize god from his works, so from memory of things and from invention and from the speed of its motion and all the beauty of virtue, recognize the divine power of the mind. In what place, then, is it? In the head, I believe — and I can bring reasons why I believe so. But where the soul is, another time; that it is in you is certain. What is its nature? Its own, I think, and proper to it. But grant it to be fiery, grant it to be breath: that is nothing to the matter at hand. Only see this: that as you know god, even though you are ignorant both of his place and of his face, so your soul ought to be known to you, even if you are ignorant both of its place and of its form.

\ciceroSection{71} But in the knowledge of the soul we cannot doubt — unless we are utterly leaden in natural philosophy — that there is nothing mixed into souls, nothing compounded, nothing coupled, nothing cemented together, nothing twofold; and since this is so, surely the soul can neither be separated nor divided nor torn apart nor pulled asunder, and therefore cannot perish either. For perishing is, as it were, the departure and parting and sundering of those parts that, before the perishing, were held together by some joining. Moved by these and like considerations, Socrates neither sought an advocate for his capital trial nor begged before his judges, but brought to it a free defiance drawn from greatness of soul, not from arrogance; and on the last day of his life he discoursed at length on this very subject; and a few days before, when he could easily have been led out of prison, he refused; and then, when he was already all but holding in his hand that deadly cup, he spoke in such a way that he seemed not to be driven to death but to be climbing up into heaven.

\ciceroSection{72} For so he held, and so he discoursed: that there are two paths and a twofold course for souls departing from the body. Those who had defiled themselves with human vices and given themselves wholly to lusts — blinded by which they had either stained themselves with private vices and outrages or, by violating the commonwealth, had conceived crimes beyond all atonement — for these there is a certain devious road, shut off from the council of the gods. But those who had kept themselves whole and pure, who had had the least contact with their bodies and had always called themselves back from them, and who, in human bodies, had imitated the life of the gods — for these the return lies easy and open to those from whom they had set out.

\ciceroSection{73} And so he recalls how, as the swans — dedicated to Apollo not without reason, but because they seem to have from him the gift of prophecy, by which, foreseeing what good there is in death, they die with song and with joy — so all good and learned men should do the same. (And indeed no one could doubt this, if only the same thing did not happen to us when we think carefully about the soul as happens often to those who, gazing keenly with their eyes at the failing sun, lose their sight altogether: so the mind's edge, gazing at itself, sometimes grows dull, and for that reason we lose our diligence in contemplation. And so, doubting, looking about, hesitating, fearing many adverse things, our discourse is carried along as if on a raft on a measureless sea.

\ciceroSection{74} ) But these things are both old and from the Greeks; while Cato departed from life rejoicing that he had found a cause for dying. For that god who rules within us forbids us to depart from here without his command; but when god himself has given a just cause — as then to Socrates, now to Cato, often to many others — then surely, by the God of Faith, the wise man will pass gladly out of this darkness into that light; and yet he will not break the chains of his prison — for the laws forbid it — but, as if summoned and released by a magistrate or by some lawful authority, so he will go out, called forth and discharged by god. For the whole life of philosophers, as the same man says, is a rehearsal of death.

\ciceroSection{75} For what else are we doing when we call the soul away from pleasure — that is, from the body — from property, which is the body's servant and handmaid, from the commonwealth, from all business: what, I say, are we then doing but summoning the soul to itself, compelling it to be with itself, and drawing it as far as possible away from the body? And to separate the soul from the body is nothing else than to learn how to die. Therefore let us rehearse this, believe me, and let us part ourselves from our bodies — that is, let us grow used to dying. This, while we are still on earth, will be like that heavenly life; and when, released from these chains, we are carried up there, the course of our souls will be the less hindered. For those who have always been in the fetters of the body, even when set loose, advance more slowly, like men who have been bound in irons for many years. And when we have come there, then at last we shall live. For this life, at least, is death — which I could lament, if I cared to. —You have lamented it enough in your \textit{Consolation};

\ciceroSection{76} and when I read it, I want nothing more than to leave these things behind — and now that I have heard all this, far more so. The time will come, and quickly, whether you draw back or hasten on; for life flies. And so far is death from being an evil, as it seemed to you a while ago, that I fear there may be nothing else that is not an evil to man, nothing else more to be preferred as a good — if indeed we are to be either gods ourselves or in the company of the gods. —What does it matter? There are those present who do not approve of this. —But I shall never let you go from this discussion in such a way that death could by any reasoning seem to you an evil. —How can it, now that I have learned these things? —You ask how it can?

\ciceroSection{77} Crowds come out against me on the other side — and not only Epicureans, whom for my part I do not look down on, though somehow it is precisely the most learned who hold them in contempt. The keenest of all, my own darling Dicaearchus, argued against this immortality. He wrote three books, called the \textit{Lesbiaca}, because the conversation is set at Mytilene; in them he means to prove that souls are mortal. The Stoics, for their part, grant us a lease, as if we were crows: they say souls will last a long while — they deny they last forever. Well then, do you not want to hear why, even if it should be so, death is still no evil? — As you please; but no one will dislodge me from immortality. — I commend that, though one ought never to be too confident of anything;

\ciceroSection{78} for we are often shaken by some sharply turned conclusion; we waver and change our minds even in matters that are clearer than this — and in this there is a certain obscurity. So if that should happen, let us be armed. — By all means; but I will see to it that it does not happen. — Is there any reason, then, why we should not dismiss our friends the Stoics? I mean those who say that souls remain after they have left the body, but not forever. — Those, surely, who take on what is the hardest thing of all in this whole question — that the soul can remain when it is free of the body — but will not grant the other thing, which is not only easy to believe, but, once their own point is conceded, follows from it: that, having remained a long while, it should not perish.

\ciceroSection{79} — A fair reproach, and the matter does stand so. Are we then to believe Panaetius, dissenting from his own Plato? For the man whom he everywhere calls divine, calls the wisest, calls the most holy, calls the Homer of philosophers — this single doctrine of his, about the immortality of souls, he does not accept. He holds — what no one denies — that whatever has been born perishes; and that souls are born, as is shown by the resemblance of those who are begotten to their parents, a resemblance that appears in temperament too, not in bodies alone. He brings forward a second argument as well: that there is nothing capable of feeling pain that is not also capable of falling ill; and that whatever is liable to sickness will also perish; but souls feel pain — therefore they also perish.

\ciceroSection{80} These points can be refuted. They come of failing to understand that, when we speak of the eternity of souls, we are speaking of the mind, which is always free of every turbulent motion — not of those parts in which distress and anger and the appetites have their seat, parts which the very man against whom this is said holds to be set apart from the mind and walled off from it. Now resemblance shows itself more plainly in beasts, whose souls are devoid of reason; among human beings the resemblance stands out rather in the figure of the body, and it matters greatly to the souls themselves in what kind of body they are lodged. For many things arise from the body that sharpen the mind, and many that dull it. Aristotle, indeed, says that all men of genius are melancholic — so that I need not take it amiss to be of the slower sort. He lists many such men, and offers, as though it were settled, a reason why this should happen. But if there is so great a power to shape the disposition of the mind in the things engendered in the body — and these, whatever they are, are what produce the resemblance — then resemblance brings no necessity at all that souls should be born. I pass over the cases of dissimilarity.

\ciceroSection{81} I could wish Panaetius were here — he lived with Africanus — and I would ask him which of his own kin the grandson of Africanus's brother resembled: in face, perhaps his father's; in life, so like all the ruined men there are that he was easily the worst of them. I would ask the same about a man who was likewise the grandson of Publius Crassus, a man wise and eloquent and of the first rank, and about the grandsons and sons of many other distinguished men, whom there is no point in naming. But what are we doing? Have we forgotten that what we set before ourselves now — once we had said enough about eternity — was to show that, even if souls do perish, there is nothing evil in death? — I had not forgotten; but while you were speaking of eternity I was content to let you wander from the point.

\ciceroSection{82} I see that you look high, and wish to migrate to heaven. I hope that lot may fall to us. But suppose, as those men will have it, that souls do not remain after death: I see that, if it is so, we are deprived of the hope of a happier life; but what evil does that view bring with it? Grant that the soul perishes just as the body does: is there then any pain, or any sensation at all, in the body after death? No one says so — and though Epicurus charges Democritus with it, the Democriteans deny it. Neither, then, does sensation remain in the soul; for the soul itself is nowhere. Where, then, is the evil, since there is no third thing? Is it that the soul's very departure from the body does not happen without pain? Even granting that it is so — how slight a thing that is! But I judge it false, and that for the most part it happens without sensation, sometimes even with pleasure, and that the whole business is trivial, whatever it is;

\ciceroSection{83} for it happens in an instant of time. What torments — or rather, racks — us is the departure from all those things that are good in life. See whether it might not more truly be called a departure from evils. Why should I now lament the life of men? I could do it, truly and with justice; but what need is there, when my whole aim is that we should not think ourselves wretched after death, to make life itself more wretched still by deploring it? I did that in the book in which I consoled my own self as well as I could. Death, then, leads us away from evils, not from goods — if we want the truth. And this point, indeed, is argued so copiously by Hegesias the Cyrenaic that he is said to have been forbidden by King Ptolemy to say such things in his lectures, because many who heard them took their own lives.

\ciceroSection{84} There is, too, the epigram of Callimachus on Theombrotus of Ambracia, who — though no misfortune had befallen him — threw himself, the poet says, from the wall into the sea, having read Plato's book. And of that Hegesias I mentioned there is a book entitled \textit{The Man Who Starves Himself to Death} \textit{[Greek: Apokarteron]}, in which a man taking his leave of life by starvation is called back by his friends; answering them, he reckons up the troubles of human life. I could do the same — though less than he, who thinks it profits no one at all to live. I leave the rest aside: does it profit even us? We who, robbed of the consolations and adornments both of home and of the Forum — surely, if we had died sooner, death would have drawn us away from evils, not from goods.

\ciceroSection{85} Suppose, then, there is some man who has nothing evil, who has taken no wound from fortune: the famous Metellus, with his four honored sons, or Priam with his fifty, seventeen of them born of a lawful wife. Over both of them fortune held the same power, but in the one she used it: for Metellus many sons and daughters, grandsons and granddaughters, laid upon the pyre; Priam, bereft of so great a progeny, was killed by an enemy's hand when he had fled to the altar. If he had died while his sons lived and his kingdom stood unharmed — \textit{while barbarian wealth stood about him, beneath fretted and paneled ceilings} — would he then, after all, have departed from goods or from evils? Then, surely, he would seem to depart from goods. Yet certainly it would have gone better with him, and those lines would not be sung so tearfully: \textit{All this I saw — set ablaze, Priam's life torn from him by violence, Jupiter's altar fouled with blood.} As though, in fact, anything better could have befallen him then at that hand of violence! But if he had died sooner, he would have escaped such an end altogether; as it was, at this hour he lost the sense of his evils.

\ciceroSection{86} For our intimate friend Pompey, when he was gravely ill at Naples, things took a turn for the better. The Neapolitans wore garlands, and so, naturally, did the people of Puteoli; from the towns the public crowds came pouring out to offer congratulations — a silly business, to be sure, and rather Greek, yet a fortunate one. Now if he had been carried off then, would he have departed from good things or from evil ones? Certainly from wretched ones. For he would not have waged war with his father-in-law, would not have taken up arms unprepared, would not have left his home, would not have fled from Italy; he would not, his army lost, have fallen stripped and defenseless onto the swords and into the hands of slaves; his children would not have been mourned, all his fortunes would not be in the possession of the victors. If he had met death then, he would have died in the fullness of his fortunes; by prolonging his life, how many, how great, how incredible the disasters he drained to the dregs! These are the things death escapes — and even if they have not happened, still, because they can happen, the same holds. But men do not reflect that such things can befall them: each one hopes for himself the fortune of Metellus, as though either the fortunate were more numerous than the unhappy, or there were anything certain in human affairs, or it were wiser to hope than to fear. But grant this very point — that men are deprived of good things by death:

\ciceroSection{87} are the dead, then, deprived too of the comforts of life, and is that a wretched thing? They must necessarily say so. But can one who does not exist be deprived of anything at all? The very word "deprivation" is a grim one, because this force lies beneath it: he had, he no longer has; he misses, looks for, lacks. These, I suppose, are the discomforts of one deprived: a man is deprived of his eyes — blindness is hateful; of his children — bereavement. This holds among the living; but the dead are deprived not only of life's comforts, but not even of life itself. I am speaking of the dead, who do not exist: we who do exist — are we deprived of horns, or of wings? Would anyone say so? Surely no one. Why is that? Because, when you do not have a thing that is suited to you neither by use nor by nature, you are not deprived of it, even if you are aware that you do not have it.

\ciceroSection{88} This must be pressed again and again, and driven home, once that other point has been established — about which, if souls are mortal, we cannot doubt: that the destruction in death is so total that not even the least trace of sensation is left. So, once this has been duly settled and fixed, that other matter must be shaken out, so that we may know what it is to be deprived, and no error be left lurking in the word. To be deprived, then, means this: to lack something that you would wish to have; for wanting is contained in deprivation — except when the word is used in some other sense, as when one is said to be "without" a fever. For "to be without" is also used in another way, when you do not have something and are aware that you do not have it, even though you bear it easily. In this sense one is not said to be "deprived" in death; for then there would be nothing to grieve over. What is said is "to be deprived of a good" — which is an evil. But not even the living man is deprived of a good if he does not need it; yet in a living man it can still be understood that you are deprived of a kingdom — though even this cannot be said of you with sufficient precision; it could be said of Tarquin, when he had been driven from his kingdom. But in a dead man it cannot even be understood. For deprivation belongs to one who feels, and there is no feeling in the dead: there is, therefore, not even deprivation in the dead. Yet what need is there to philosophize about this, when we see that the matter scarcely calls for philosophy?

\ciceroSection{89} How many times have not only our commanders, but whole armies as well, rushed to a death they did not doubt was certain! And if that death were feared, Lucius Brutus would not have fallen in battle while barring the return of the tyrant whom he himself had driven out; the elder Decius fighting the Latins, his son fighting the Etruscans, his grandson fighting Pyrrhus would not have thrown themselves upon the enemy's weapons; Spain would not have seen the Scipios falling for their country in a single war, nor Cannae have seen Paulus and Geminus, nor Venusia Marcellus, nor Litana Albinus, nor Lucania Gracchus. Is any of these wretched today? No, nor even then, after his last breath; for no one can be wretched once sensation has been destroyed.

\ciceroSection{90} "But that very thing is hateful — to be without sensation." Hateful, if it were deprivation; but since it is plain that nothing can exist in one who does not himself exist, what can there be that is hateful in one who is neither deprived nor feels? Granted, I have said this far too often; but it is because in this lies all the soul's contraction that comes from the fear of death. For whoever has seen clearly what is plainer than daylight — that, soul and body alike consumed, the whole living creature destroyed, and the destruction made complete, that animal which once existed has become nothing — such a man will see plainly that there is no difference between Hippocentaur, who never existed, and King Agamemnon; and that Marcus Camillus minds this civil war now no more than I, while he was alive, minded the capture of Rome. Why, then, should Camillus have grieved, had he supposed these things would come to pass some three hundred and fifty years afterward, and why should I grieve, were I to suppose that some nation will take possession of our city ten thousand years from now? Because love of country is so great that we measure it not by our own sensation, but by the very safety of that country. And so death does not deter the wise man —

\ciceroSection{91} death, which, on account of life's uncertain chances, hangs over him every day, and, on account of life's brevity, can never be far off — does not deter him from taking thought, for all time to come, for the commonwealth and for his own, even though he counts posterity itself, of which he will have no sensation, as something that concerns him. And so a man may, even while judging the soul to be mortal, set in motion works that are eternal — not from a craving for glory, which he will not feel, but for virtue, which glory necessarily follows, even if you do not aim at it. If nature truly stands so that, just as our birth brought us the beginning of all things, so death brings the end, then, as nothing concerned us before our birth, so nothing will concern us after death. And where in this can there be any evil, since death concerns neither the living nor the dead?

\ciceroSection{92} As for the dead, they either no longer exist at all, or what comes will never touch them. Those who make death lighter wish it to be just like sleep — as if anyone would want to live ninety years on the terms that, once he had finished sixty, he should sleep away the rest; not even his own family would want that for him, let alone the man himself. Endymion, indeed — if we care to listen to the old tales — when he fell asleep, who knows how long ago, on Latmos, a mountain of Caria, has not yet, I think, woken up. Do you suppose, then, that he is troubled when the Moon is in distress — the Moon by whom he is believed to have been lulled to sleep, that she might kiss him as he slept? What should trouble a man who does not even feel? You have in sleep the very image of death, and you put it on every day: and do you still doubt that there is no sensation in death, when in its likeness you see there is no sensation at all?

\ciceroSection{93} Let us banish, then, this near-senile foolishness — that to die before one's time is a wretched thing. What time, after all? Nature's? But nature gave us the use of life as she gives us money, with no day fixed for repayment. What grounds have you for complaint, then, if she calls it in when she pleases? It was on that condition that you received it. The same people think that if a small child dies it must be borne with composure, and if a baby dies in the cradle there is no ground for complaint at all. And yet from the infant nature exacted more harshly what she had given. "He had not yet tasted the sweetness of life," they say; "but this one was already hoping for great things, and had begun to enjoy them." Yet in all other matters we judge it better to come by some part of a thing than none at all: why should it be otherwise in life? (Though Callimachus is not wrong to say that Priam wept far more often than Troilus.) The fortune of those who die at the end of a full age, by contrast, is called blessed.

\ciceroSection{94} Why? Because, I imagine, no man's life, were it granted longer, could be made more pleasant; for surely nothing is sweeter to a man than wisdom — and wisdom, whatever else it takes away, old age at least brings. But what age is long, or what is long for a man at all? Has not old age, pursuing close behind \textit{boys at one moment, young men at another, in mid-course, overtaken them unawares}? But because we have nothing beyond it, we call this long. All these spans, just as they are allotted to each thing by its due measure, are called long or short accordingly. Beside the river Hypanis, which flows from the European side into the Black Sea, Aristotle says that certain little creatures are born which live a single day. One of these that has died at the eighth hour has died at an advanced age; one that has died at sunset is decrepit, all the more so if it died on the longest day of summer. Set our longest span beside eternity: we shall be found to be of almost the same brevity as those little creatures.

\ciceroSection{95} Let us, then, despise all this foolishness — for what lighter name can I give to such lightness? — and place the whole power of living well in strength and greatness of soul, in the contempt and disdain of all human affairs, and in every virtue. For as it is, we are unmanned by the softest reflections, so that if death arrives before we have laid hold of the Chaldeans' promises, we think ourselves robbed of certain great goods, cheated and abandoned.

\ciceroSection{96} But if we hang in suspense from waiting and longing, if we are racked, if we are tormented — by the immortal gods, how pleasant that journey ought to be, once it is finished, when no further care, no anxiety will remain! How Theramenes delights me! How lofty is his spirit! For though we weep as we read, that famous man does not die a pitiable death. When he had been thrown into prison by order of the Thirty Tyrants, he drank off the poison as a thirsty man drinks, then flung out what was left from the cup so that it rang; and at the sound it made, he said with a smile, "Here's to the handsome Critias!" — Critias, who had been his most savage enemy. For the Greeks at their banquets are accustomed to name the man to whom they mean to pass the cup. That outstanding man made his jest with his last breath, when he already held in his vitals the death he had drawn in; and truly he foretold for the man to whom he had drunk that very death which followed soon after. Who would praise this even-tempered greatness of soul in the face of death itself, if he judged death an evil?

\ciceroSection{97} For Socrates goes to the same prison and, a few years later, to the same cup, condemned by the same crime in his judges as Theramenes was in his tyrants. What, then, is the speech which Plato gives him to make before the judges when he had already been sentenced to death? "A great hope holds me, judges," he says, "that it goes well with me in being sent to my death. For one of two things must be true: either death takes away all sensation entirely, or it is a passage by death from this place to some other. Therefore, if sensation is extinguished and death is like that sleep which sometimes, even without the visions of dreams, brings the most peaceful rest — good gods, what a gain it is to die! Or how many days could be found that one would set above such a night! And if all the everlasting time to come is like that night, who is happier than I?

\ciceroSection{98} But if what is said is true — that death is a migration to those shores which they who have departed from life inhabit — then it is far happier still. To think that, once you have escaped from those who wish to be counted as judges, you should come to those who are truly called judges — Minos, Rhadamanthus, Aeacus, Triptolemus — and should meet with those who have lived justly and in good faith: can this journey seem a trifling thing to you? And to be allowed to converse with Orpheus, Musaeus, Homer, Hesiod — at how much, after all, do you reckon that? For my part, I would be willing to die many times over, if it could be done, that I might be permitted to see for myself the things I speak of. And with what delight I should be filled when I met Palamedes, when I met Ajax, when I met others who were brought down by an unjust verdict! I would test, too, the wisdom of the supreme king who led the greatest forces to Troy, and of Ulysses and Sisyphus; and for pursuing these inquiries, just as I did here, I should not be condemned to death. — And you, too, judges who have acquitted me, do not fear death.

\ciceroSection{99} For nothing evil can befall a good man, neither living nor dead, nor will his affairs ever be neglected by the immortal gods, nor has this fate come upon me by chance. And as for those by whom I was accused or by whom I was condemned, I have no reason to be angry with them — except that they believed they were doing me harm." And so he spoke in this fashion; but nothing was better than his final words: "But it is time," he said, "to go from here now — for me to die, for you to live on. Which is the better, the immortal gods know; no mortal, I think, knows." Truly I would far rather have this man's spirit than the fortunes of all those who passed judgment on him. And yet — what he denies that anyone but the gods can know, whether it is the better, he knows himself; for he said so earlier — but he holds to the end to that principle of his, that he affirms nothing.

\ciceroSection{100} But let us hold to this: that we judge nothing an evil which is given by nature to all, and let us understand that, if death is an evil, it is an everlasting evil. For death seems to be the end of a wretched life; if death is wretched, there can be no end. But why do I cite Socrates or Theramenes, men outstanding in the glory of virtue and wisdom, when a certain Spartan — whose name has not even been handed down — so despised death that, as he was being led to it, condemned by the ephors, and his face was cheerful and glad, an enemy of his said to him, "Do you despise the laws of Lycurgus?" and he answered, "On the contrary, I am deeply grateful to the man who has fined me with a penalty I can pay off without a loan and without borrowing at interest." What a man, worthy of Sparta! To me, indeed, one of so great a spirit seems to have been condemned an innocent man. Such men our own state has produced beyond number.

\ciceroSection{101} But why should I name commanders and leading men, when Cato writes that legions often marched eagerly to a place from which they did not expect to return? With the same spirit the Spartans fell at Thermopylae — on whom Simonides wrote: \textit{Tell, stranger, that you saw us lying here at Sparta, while we obeyed our country's sacred laws.} And what does Leonidas, their commander, say? "Press on with stout hearts, Spartans; today perhaps we shall dine among the dead." This was a brave people, while the laws of Lycurgus held their force. One of them, when a Persian enemy had said boastingly in a parley, "You will not see the sun for the multitude of our javelins and arrows," replied, "Then we shall fight in the shade."

\ciceroSection{102} I cite men: what of the Spartan woman, then? When she had sent her son into battle and heard that he had been killed, she said, "It was for this that I bore him — that there might be one who would not hesitate to meet death for his country." So be it: the Spartans are brave and hardy; the discipline of the state has great power. What of it? Do we not marvel at Theodorus of Cyrene, a philosopher of no mean standing? When King Lysimachus threatened him with the cross, he said, "Threaten those horrors, I beg you, to your courtiers in their purple: it makes no difference to Theodorus whether he rots on the ground or up in the air." His words remind me that I should say something also about burial and the tomb — no difficult matter, especially now that what was said a little earlier about feeling nothing is understood. As for what Socrates thought on the subject, it is plain in that book in which he dies — a book of which we have already said so much.

\ciceroSection{103} For when he had discussed the immortality of the soul and the time of dying was already pressing, and Crito asked him how he wished to be buried, he said, "I have spent much effort, my friends, in vain; for I have not persuaded our friend Crito that I shall fly away from here and leave nothing of myself behind. Nevertheless, Crito, if you can catch up with me, or come upon me anywhere, bury me as you see fit. But, believe me, none of you will overtake me when I have departed from here." Finely said, indeed, by one who both gave his friend leave and showed that he cared nothing for this whole matter.

\ciceroSection{104} Harsher was Diogenes — who held the same view, but with a Cynic's roughness — for he ordered that he be thrown out unburied. Then his friends: "To the birds and beasts?" "By no means," he said, "but set a stick beside me, that I may drive them off." "How will you be able to? You will not feel." "Then what harm will the mangling of the beasts do me, when I feel nothing?" Finely said by Anaxagoras, who, when he was dying at Lampsacus and his friends asked whether he wished to be carried back to his homeland at Clazomenae if anything should happen to him, replied, "There is no need; for from everywhere the road to the underworld is the same length." And on the whole question of burial there is one thing to hold to: that it pertains to the body, whether the soul has perished or still thrives. But in the body it is plain that, whether the soul is extinguished or has slipped away, no sensation remains. Yet everything is full of errors.

\ciceroSection{105} Achilles drags Hector, bound to his chariot: he thinks, I suppose, that Hector is being torn and feels it. So he takes his revenge, as he imagines; but she mourns it as a most bitter thing: \textit{I saw — what I endured most grievously to see — Hector being dragged by a four-horse chariot.} What Hector? Or how long will he be Hector? Better is Accius, and an Achilles wise for once: \textit{No — rather I gave back the body to Priam; Hector I carried off.} It was not Hector, then, that you dragged, but the body that had been Hector's.

\ciceroSection{106} See, another rises from the earth who will not let his mother sleep: \textit{Mother, I call on you — you who relieve with sleep the care that hung suspended; you have no pity on me. Rise, and bury your son —!} When these lines are sung in subdued and tearful strains that bring grief upon whole theaters, it is hard not to judge the unburied wretched. \textit{Before the beasts and birds —} he fears that they will make poor use of his mangled limbs; of their being burned, he has no dread. \textit{Nor let the half-eaten remnants, with the bones laid bare, be foully dragged across the ground, smeared with gore —} I do not understand what he fears, when he pours out such fine seven-foot verses to the pipe. We must hold, then, that nothing is to be cared for after death — though many take vengeance even on dead enemies. In Ennius, Thyestes curses, in verses splendid enough, that first Atreus may perish by shipwreck: this is harsh, surely, for such a death is not without grievous sensation; but the rest is empty: \textit{Let him himself, impaled on the topmost jagged rocks, disemboweled, hanging by his side, splattering the rocks with foul matter, with gore and black blood —}

\ciceroSection{107} The very rocks will not be more wholly devoid of all sensation than that man hanging from the cliff, on whom this fellow thinks himself to be wishing torment. Those things would be hard, if he could feel them; without feeling, they are nothing. And then this, utterly empty: \textit{Let him have no tomb to receive him, no harbor for the body, where, once human life is laid aside, the body may rest from its ills.} You see in how great an error these lines are caught up: he supposes that there is a harbor for the body, and that the dead man rests within his tomb. Great is the fault of Pelops, who never schooled his son nor taught him how far each thing was worth caring for.

\ciceroSection{108} But why should I take note of the opinions of individuals, when I may survey the various errors of whole nations? The Egyptians embalm their dead and keep them in the house; the Persians even smear them over with wax before burial, so that the bodies may last as long as possible. It is the custom of the Magi not to bury the bodies of their own people unless they have first been mangled by wild beasts; in Hyrcania the common folk keep public dogs, the nobles dogs of their own household — and we know that breed of dogs to be a fine one, but each man, according to his means, provides the creatures by which he is to be torn, and they reckon that the best of burials. Chrysippus gathers a great many other such customs, inquisitive as he is in every branch of history, but some are so revolting that speech shuns and recoils from them. This whole topic, then, is to be despised where we ourselves are concerned, not neglected where our own people are concerned — yet so, that we who are living should feel that the bodies of the dead feel nothing;

\ciceroSection{109} but how much is to be granted to custom and reputation — that the living should attend to, yet so as to understand that it has nothing to do with the dead. And surely death is met with the most untroubled mind when the dying man can console himself in his own praises. No one has lived too short a time who has discharged the perfect duty of perfect virtue. For me too there were many seasons fit for death. Would that I could have met it then! For nothing further was being gained, the obligations of life had been heaped full, only struggles against fortune remained. So if reason itself should fail to bring us to disregard death, let a life well lived bring us to seem to have lived enough and more than enough. For though sensation may have departed, still the dead are not deprived of their own proper goods of praise and glory, even though they do not feel them. For although glory has nothing in itself for which it should be sought, still it follows virtue like a shadow.

\ciceroSection{110} But the judgment of the multitude about good men, if it is ever sound, is more to be praised than they are to be called blessed on that account. Yet I cannot say — however this will be received — that Lycurgus and Solon are without glory for their laws and their ordering of the state, or Themistocles and Epaminondas for their valor in war. For Neptune will sooner overwhelm Salamis itself than the memory of the trophy at Salamis, and Leuctra will sooner be lifted out of Boeotia than the glory of the battle of Leuctra. And far more slowly still will fame desert Curius, Fabricius, Calatinus, the two Scipios, the two Africani, Maximus, Marcellus, Paulus, Cato, Laelius, and countless others; and whoever has seized upon some likeness to these men — measuring it not by the fame of the crowd but by the true praise of good men — will, if circumstance so requires, walk to death with a confident mind, in which we have learned that there is either the highest good or no evil at all. And in the midst of his own prosperity he will even wish to die; for the heaping up of goods cannot be so pleasant as their withdrawal is grievous. This sentiment seems to be conveyed in that saying of the Spartan who, when the Rhodian Diagoras, a famous Olympic victor, had seen his two sons crowned victors at Olympia on a single day, approached the old man and, congratulating him, said:

\ciceroSection{111} "Die now, Diagoras — for you will not be ascending into heaven." Great honors, these, and perhaps too great, the Greeks think them — or rather thought them then; and the man who said this to Diagoras, judging it a very great thing that three Olympic victors should come forth from one house, thought it useless to him that he should linger longer in life, exposed to fortune. Now to you I had given, as it seemed to me, an answer in few words — enough for the purpose, since you had granted that the dead are in no evil — but I pressed on to say more for this reason: that here, in longing and grief, lies the greatest consolation. For the grief we take up on our own behalf and for our own sake we ought to bear in measure, lest we seem to love ourselves; but that suspicion racks us with intolerable pain — if we imagine that those of whom we have been bereft exist with some sensation amid the very evils the crowd imagines. This opinion I wished to shake out of myself by the roots, and for that reason I have perhaps gone on rather long. — You, long?

\ciceroSection{112} Not in my view, at any rate. For the earlier part of your discourse made me wish to die, the later part brought me at one moment merely not to refuse it, at another merely not to be troubled by it; but by the whole of your discourse this much at least has surely been accomplished — that I no longer count death among evils. Do we then want even the orators' epilogue? Or do we now abandon that art altogether? Do not abandon it, the art you have always adorned, and rightly so — for it, if we want to speak the truth, had adorned you. But what, pray, is this epilogue? I am eager to hear it, whatever it is.

\ciceroSection{113} In the schools they are wont to bring forward, on the subject of death, the judgments of the immortal gods — and not, indeed, inventing them themselves, but on the authority of Herodotus and many others. First are celebrated Cleobis and Bito, the sons of the priestess of Argos. The story is well known. When it was her right to be carried in a chariot to a solemn and appointed sacrifice — the shrine being a fair distance from the town — and the draught animals were delaying, then the young men I have just named laid aside their clothing, anointed their bodies with oil, and put themselves to the yoke. So the priestess was conveyed to the shrine, the chariot drawn by her sons, and she is said to have prayed to the goddess to grant them, as a reward for their devotion, whatever was the greatest thing that could be given to a man by a god; and after the young men had feasted with their mother and given themselves to sleep, in the morning they were found dead.

\ciceroSection{114} Trophonius and Agamedes are said to have used a similar prayer; when they had built the temple of Apollo at Delphi, in worship of the god they asked as the wage of their work and labor no small thing, indeed: nothing definite, but whatever was best for a man. To them Apollo showed that he would grant it on the third day after that day; and when that day dawned, they were found dead. They say the god gave his judgment — and a god, too, to whom the other gods had granted that he should foretell beyond all the rest. There is also told a certain little fable about Silenus, who, when he had been taken captive by Midas, is recorded to have given him this gift in return for his release: he taught the king that it is by far best for a man not to be born, and the next best to die as soon as possible.

\ciceroSection{115} This sentiment Euripides used in his \textit{Cresphontes}: \textit{For it would have befitted us, gathering in throngs, to mourn at the house where someone has been brought forth into the light, reckoning up the many ills of human life; but the man who had ended his heavy labors in death — him to follow forth with all praise and gladness.} Something similar occurs in the \textit{Consolation} of Crantor; for he tells of a certain Elysius of Terina who, while grieving deeply for the death of his son, came to an oracle of the dead to ask what had been the cause of so great a calamity, and on three tablets these three little verses were given him: \textit{Men wander through life with minds in ignorance: Euthynous, by the will and power of the fates, has come into death. Thus it was more profitable that an end be made, both for him and for you.}

\ciceroSection{116} Drawing on these and similar authorities, they confirm that the case has been judged by the immortal gods. Alcidamas, indeed, a rhetorician of the ancients and among the most distinguished, even wrote a praise of death, which consists of an enumeration of human ills; he lacked the reasons that are more searchingly assembled by the philosophers, but abundance of language he did not lack. Glorious deaths met for one's country are wont to seem to the rhetoricians not only glorious but even blessed. They go back to Erechtheus, whose daughters too eagerly sought death for the lives of the citizens; they recall Codrus, who flung himself into the midst of the enemy in a servant's dress, so that he might not be recognized had he been in royal array — for an oracle had been given that, if the king were killed, Athens would be victorious; Menoeceus is not passed over, who likewise, when an oracle had been pronounced, lavished his blood upon his country; for Iphigenia at Aulis bids herself be led off to be sacrificed, that the enemy might be drawn into shedding their own. From there they come to nearer instances: Harmodius is on every lip, and Aristogiton; the Spartan Leonidas, the Theban Epaminondas flourish. Our own men they do not know — to enumerate them would be a great labor: so many are they for whom we see that death was a thing to be desired, with glory. Since this is so, must one nonetheless employ great eloquence, and harangue as if from a higher platform, so that men may either begin to wish for death or at the least cease to fear it?

\ciceroSection{117} For if that final day brings not extinction but a change of place, what is more to be desired? But if it utterly destroys and obliterates, what is better than to fall asleep amid the labors of life and so, closing the eyes, to be lulled into everlasting slumber? And if that be so, Ennius's words are better than Solon's. For this poet of ours says: \textit{Let no one honor me with tears, nor make my funeral with weeping!} But that wise man: \textit{Let my death not lack tears: let us leave to our friends grief, that they may keep my funeral with lamentation.}

\ciceroSection{118} But as for us, if anything of the kind should happen, such that it seems announced by a god that we are to depart from life, let us obey gladly and with thanks, and reckon that we are being released from custody and freed from our chains — to sail back either into our eternal and truly our own home, or to be free of all sensation and trouble; but if nothing is announced, let us nonetheless be of such a mind as to think that day, dreadful to others, an auspicious one for us, and to count nothing among evils that has been ordained either by the immortal gods or by nature, the parent of all. For it was not at random nor by chance that we were sown and created; surely there was a certain power that took thought for the human race, and would not beget or nurture that which, when it had drained off all its labors, would then fall into the everlasting evil of death: let us think rather that a harbor and a refuge have been prepared for us.

\ciceroSection{119} Would that we might be carried there with sails full-spread! But if we are driven back by contrary winds, still we must be borne to the same place, only a little later. And can what is necessary for all be wretched for one alone? You have your epilogue, lest you think anything passed over or left out. Indeed I have — and that epilogue, too, has made me firmer still. Excellent, I said. But now let us yield something to your health; tomorrow, however, and for as many days as we shall be at Tusculum, let us pursue these matters, and above all those that bring relief to sicknesses, fears, and desires — which is the richest fruit of all philosophy.

\ciceroBook{2}

\ciceroSection{1} Neoptolemus, in Ennius, says that he must do some philosophy — but only a little, for he has no taste for it on the whole. I, however, Brutus, judge that I for my part must philosophize — for what better could I do, especially with nothing else on hand? — but not a little, as he would have it. For in philosophy it is hard for a man to know few things if he does not know most things, or all. Few cannot be picked out except from many; and a man who has grasped a few will pursue the rest with the same zeal.

\ciceroSection{2} Yet in a busy life, and one that — as Neoptolemus's then was — is given over to soldiering, even those few things often do great good and bear fruit: if not so great a harvest as can be gathered from philosophy entire, still enough to free us, now and then and in some measure, from desire or distress or fear. Take that discussion lately held by me at Tusculum: it seemed to produce a great contempt for death, which counts for not a little in freeing the mind from fear. For a man who dreads what cannot be avoided can in no way live with a tranquil mind; but a man who fears death not only because it must come, but because death holds nothing to dread, has secured himself a great safeguard for a happy life.

\ciceroSection{3} I am well aware, however, that many will speak zealously against this — which there was no way to avoid except by writing nothing at all. For consider my speeches, which I wished to win the approval of the crowd's judgment (since oratory is a popular art, and the test of eloquence is the approval of the listeners): there were nonetheless some who would praise nothing except what they felt confident of being able to imitate, and who set the same limit to speaking well that they set to their own hopes; who, overwhelmed by a wealth of thoughts and words, declared they preferred leanness and famine to richness and abundance — and from this arose the so-called Attic school, unknown to the very men who professed to follow it, and now fallen silent, all but laughed off the Forum itself.

\ciceroSection{4} What, then, are we to expect, when the people, whose support we relied on before, are now of no use to us at all? For philosophy is content with few judges, deliberately shunning the crowd of her own accord — to that very crowd both suspect and hated — so that anyone wishing to disparage her wholesale could do it with the people's blessing, or anyone trying to invade the particular school I most follow could find ample reinforcements in the doctrines of the other philosophers. To those who disparage philosophy entire I gave my answer in the \textit{Hortensius}; and what had to be said on behalf of the Academy I believe was set out with sufficient care in the four books of the \textit{Academica}. Yet so far am I from being unwilling to have writing directed against me that I positively wish for it. For in Greece herself philosophy would never have stood in such honor had she not flourished through the contentions and disagreements of the most learned men.

\ciceroSection{5} For which reason I urge all who are capable of it to wrest the glory of this field too from a Greece now flagging, and to carry it over into this city — just as our ancestors, by their zeal and industry, carried over all the other arts worth pursuing. The glory of oratory, indeed, drawn up from humble beginnings, has reached its summit; so that now — as nature ordains in nearly all things — it grows old, and seems destined in a short time to come to nothing. Let philosophy be born into Latin letters from this present age, and let us give her our help, and let us bear being refuted and contradicted ourselves. Those who take this hard are men bound, as if pledged and consecrated, to certain fixed and predetermined doctrines, and constrained by that necessity to defend, for consistency's sake, even positions they do not usually approve. We, who pursue what is probable and cannot advance beyond what strikes us as plausible, are ready both to refute without obstinacy and to be refuted without anger.

\ciceroSection{6} And if these studies are carried over to our own people, we shall have no need even of the Greek libraries, in which the multitude of books is endless because of the multitude of those who have written. For the same things are said by many, and from this they have crammed their books full. The same will happen with ours too, if more men flock to these studies. But let us, if we can, rouse those who, liberally educated and bringing to the work the elegance of disputation as well, philosophize by method and on a settled path.

\ciceroSection{7} For there is a certain class of men who wish to be called philosophers, and they are said to have a good many books in Latin. These I do not despise — for I have never read them; but since the very men who write them profess that they write neither with distinction nor with arrangement nor with elegance nor with ornament, I pass over reading that holds no pleasure. What the followers of that school say and think, no one even moderately learned is unaware of. And so, since they themselves take no trouble over the manner of their speaking, I do not see why they should be read by anyone except those of like mind among themselves.

\ciceroSection{8} For just as everyone reads Plato and the rest of the Socratics, and after them those who set out from these — even readers who do not approve of those doctrines, or do not pursue them with much enthusiasm — whereas almost no one but his own followers takes up Epicurus or Metrodorus, so these Latin writers are read only by those who think their doctrines correctly stated. But to my mind, whatever is committed to writing deserves to be commended to the reading of all the educated; and if I myself fall short of achieving this, I do not for that reason think it should any the less be done so.

\ciceroSection{9} And so I have always been pleased by the practice of the Peripatetics and the Academy of arguing both sides on every question — not only because there is no other way to discover what is plausible in any matter, but also because it is the greatest exercise in speaking. Aristotle first employed it, and after him those who followed. In our own day Philo, whom I often heard, made it his rule to teach the precepts of the rhetoricians at one time and those of the philosophers at another. Drawn to this practice by my friends, I devoted to it whatever time was given me at Tusculum. And so, having given the morning to declamation, as I had the day before, in the afternoon we went down to the Academy. The discussion held there I set out not as if recounting it, but in nearly the same words in which it was carried on and argued. So this, then, was the conversation we struck up as we walked, and it was led off from some such opening as this:

\ciceroSection{10} I cannot say how much I was delighted — or rather helped — by yesterday's discussion of yours. For though I am conscious of never having been too greedy of life, still there was now and then thrown across my mind a kind of fear and grief, when I considered that some day there would be an end of this light and a loss of all the goods of life. From this sort of distress I have been so set free — believe me — that I think there is nothing I need worry about less. — No wonder at all in that;

\ciceroSection{11} for this is what philosophy does: she heals minds, strips away empty anxieties, frees us from desires, drives off fears. But this power of hers is not equally effective with all; she avails much when she has embraced a fitting nature. For not only does Fortune favor the brave, as the old proverb has it, but far more does reason, which by certain precepts, as it were, strengthens the force of courage. Nature plainly bore you lofty and high-minded and contemptuous of human things; and so a discourse against death found ready lodging in a brave mind. But do you suppose these same arguments avail with the men themselves who devised and argued and wrote them — beyond a very few? For how rare is the philosopher to be found whose character, whose mind and life, are so ordered as reason demands? who reckons his teaching not a display of knowledge but a law of life? who obeys his own self and submits to his own decrees?

\ciceroSection{12} One may see some so frivolous and full of boasting that it would have been better for them never to have learned; others greedy for money, some for glory, many slaves to their appetites — so that their lives clash wonderfully with their words. And this seems to me utterly shameful. For just as a man who has set up as a grammarian and yet speaks like a barbarian, or who sings out of tune when he wishes to be held a musician, is the more disgraceful for failing in the very thing whose mastery he professes, so a philosopher who errs in the conduct of life is the more disgraceful because he stumbles in the duty he claims to teach, and, having professed the art of living, goes wrong in living. — Is there not reason to fear, then, if it is as you say, that you are decking philosophy out with a false glory? For what greater proof could there be that she does no good than that certain finished philosophers live shamefully?

\ciceroSection{13} That, however, is no proof at all. For just as not all cultivated fields are fruitful — and that line of Accius is false: \textit{Though good seed be sown in poorer ground, it will yet shine bright by its own nature} — so not all cultivated minds bear fruit. And, to keep to the same comparison: as a field, however fertile, cannot be fruitful without cultivation, so neither can a mind without teaching; each is feeble without the other. And the cultivation of the mind is philosophy: she pulls up vices by the roots, prepares minds to receive sowing, and entrusts to them — and, so to speak, plants — what, once grown, will bear the richest fruit. Let us proceed, then, as we have begun. Say, if you will, what you would like to have argued.

\ciceroSection{14} I judge pain the greatest of all evils. — Greater even than disgrace? — I do not dare to say that, and I am ashamed to have been knocked off my opinion so quickly. — You would have more reason to be ashamed if you held to it. For what is less worthy of you than to think anything worse than disgrace, dishonor, baseness? To escape these, what pain is there that ought not merely not to be refused, but sought out, undergone, welcomed of one's own accord? — That is exactly my view. — Then granting that pain need not be the supreme evil, it is certainly an evil. — Do you see, then, how much of the terror of pain you have shed at a brief reminder?

\ciceroSection{15} I see plainly, but I want more. — I shall try, for my part; but it is a great matter, and I need a mind that does not fight back. — That you shall have. For as I did yesterday, so now I shall follow reason wherever it leads me. First, then, let me speak of the weakness of many men and of the various schools of the philosophers. Foremost among them by authority and by antiquity, the Socratic Aristippus did not hesitate to call pain the supreme evil. Next, to this enervated and womanish view Epicurus showed himself a docile enough pupil. After him, Hieronymus of Rhodes said that to be free of pain is the supreme good — so great an evil did he reckon pain to be. The rest, apart from Zeno, Aristo, and Pyrrho, held much the same as you did just now: that it is indeed an evil, but that there are other things worse.

\ciceroSection{16} So the very thing that nature herself and a certain noble virtue spit out at once — to keep you, that is, from calling pain the supreme evil, and to drive you from that opinion by setting disgrace against it — in just that thing philosophy, the mistress of life, persists through so many ages. What duty, what praise, what honor will be worth so much that a man would wish to win it at the cost of bodily pain, if he has convinced himself that pain is the supreme evil? What infamy, in turn, what baseness would a man not endure to escape pain, if he has decided that pain is the supreme evil? And who is not wretched — not only then, when he is crushed by the worst pains, if in these lies the supreme evil, but even when he knows it can befall him? And to whom can it not befall? So it comes about that absolutely no one can be happy.

\ciceroSection{17} Metrodorus, indeed, thinks the man perfectly happy whose body is well constituted and who is assured it will always be so. But who is there to whom such assurance can belong? Epicurus, for his part, says things that seem to me, at least, to be courting laughter. For he affirms somewhere that if the wise man is burned, if he is tortured — you are waiting, perhaps, for him to say: he will bear it, he will endure it, he will not give in; a great praise indeed, by Hercules, and worthy of the very one by whom I swore, of Hercules himself. But for Epicurus, a harsh and hard-bitten man, this is not enough: if the wise man is in the bull of Phalaris, he will say, "How sweet this is, how little I care for it!" Sweet, even? Is it not enough if it is not bitter? Why, that very thing — that being tortured is sweet to anyone — is not usually said even by those who deny that pain is an evil: harsh, hard, hateful, against nature, they call it, yet not an evil. This man, who calls it the one evil and the worst of all evils, supposes the wise man will call it sweet.

\ciceroSection{18} I do not ask of you that you treat pain in the same terms Epicurus uses of pleasure — Epicurus, a devotee of pleasure, as you know. Let him say, if he likes, that the wise man, even inside the bull of Phalaris, would feel exactly what he would feel on a soft couch; I do not credit wisdom with such power against pain. If a man is brave in enduring it, that is enough for duty; that he should be glad of it as well, I do not demand. For pain is beyond doubt a grim thing, harsh, bitter, hostile to nature, hard to suffer and to bear. Look at Philoctetes: a man whose groans we must forgive.

\ciceroSection{19} For he had seen Hercules himself on Oeta, howling at the magnitude of his pains. So the arrows he had received from Hercules brought this man no comfort then, when \textit{a serpent's bite, its venom soaked into the veins of the flesh, sets loose its foul torments.} And so he cries out, longing for help, longing to die: \textit{Oh, who will fling me from this towering crag into the salt waves below! Now, now I am undone; the violence of the wound, the fever of the sore, destroys my life.} It seems hard to say that a man who is forced to cry out like this is not in the grip of evil, and a great evil at that.

\ciceroSection{20} But let us look at Hercules himself, broken then by pain, even as he sought immortality through death itself. What cries he sends up in Sophocles, in the \textit{Trachiniae}! When Deianira had clothed him in the tunic dyed with the Centaur's blood, and it had fastened to his flesh, he says: \textit{Oh, how much that is grievous to tell, harsh to endure, have I borne, drained dry in body and in soul! Neither Juno's relentless terror nor the grim Eurystheus brought such ruin on me as one mad daughter, the child of Oeneus. She it was who snared me, all unknowing, in this robe of fury, which clings to my sides and tears my flesh with its bite, and bearing hard upon me drinks the breath of my lungs. Already it has drunk down all my discolored blood. So my body, devoured by this hideous calamity, has wasted to nothing; bound in a woven plague, I am destroyed. No enemy's hand did this, nor the earthborn mass of the Giants, nor the Centaur with his twofold-bodied charge struck this blow against my body, no Greek force, no foreign monstrousness, no savage race banished to the ends of the earth — through all of which I roamed, driving out every wildness. No: a man, I am destroyed by a woman, by a woman's hand. O son, claim now in truth this name of son for your father; let no love of your mother prevail over me as I die. Seize her, drag her here to me, torn from those dutiful hands; now let me see whether you count her or me the dearer.}

\ciceroSection{21} \textit{Go on, be bold, my son! Weep over your father's torments, have pity: the nations will lament our miseries. Oh, that I should pour out a girl's wailing from my mouth — I whom no one ever saw groan at any evil! My manhood, unmanned, is struck down and dies. Come close, my son, stand by me, look upon the pitiable, the disemboweled body of your mangled father! Look, all of you — and you, father of the gods of heaven, hurl, I beg, your flashing bolt of lightning upon me! Now, now the whirling points of anguish-bearing pain wrench at me, now the burning creeps along. O hands once victorious,}

\ciceroSection{22} \textit{O chest, O back, O the swelling muscles of my arms — was it under your grip that the Nemean lion once gnashed and gasped out his last breath in agony? Did this right hand subdue Lerna by slaughtering the foul serpent? Did this hand crush the twin-bodied band? Did it lay low the devastating beast of Erymanthus? Did this hand drag up from the dark region of Tartarus the three-headed hound, offspring of the Hydra? Did this hand kill the dragon that with coiling, ever-multiplying twists kept guard with watchful gaze over the golden-bearing tree? Many another labor my victorious hand has worked through, and no man has ever taken the spoils of my glory.} Can we despise pain, when we see Hercules himself suffer it so past all bearing?

\ciceroSection{23} Let Aeschylus come forward — not a poet only, but a Pythagorean too; for so we have been told. How does Prometheus in his work bear the pain he meets with for the theft on Lemnos — \textit{the source from which fire is famed to have been parceled out to mortals in secret; which fire the cunning Prometheus stole by guile, and so paid the penalty to Jupiter by the supreme decree of fate?} Paying that penalty, then, nailed fast to the Caucasus, he speaks these words: \textit{Offspring of the Titans, kindred of my blood, born of Heaven, behold me bound to harsh rocks and fettered, as fearful sailors in the night, in their dread, lash their ship against the roaring strait. Thus has Saturn's son, Jupiter, pinned me here, and the will of Jupiter pressed into his service the hands of Mulciber. These wedges he drove in with cruel craftsmanship, splitting my limbs apart; pierced through by such skill, wretched, I dwell in this stronghold of the Furies.}

\ciceroSection{24} \textit{Now every third grim day, with mournful swoop, the satellite of Jupiter mangles me with hooked talons and tears me apart in savage feeding. Then, gorged and glutted to the full on my rich liver, it pours out a vast screeching, and soaring on high it fawns with feathered tail upon my blood. But when my consumed liver has swelled and grown back, then once more it returns, ravenous, to its foul feast. So I feed this guardian of my dismal torment, which fouls me, living, with unending misery. For, as you see, fettered in Jupiter's chains, I cannot drive the dread bird from my breast.}

\ciceroSection{25} \textit{So, bereft, I take upon myself these tormenting plagues, in love with death, seeking the end of my suffering; but far from death the will of Jupiter drives me back. And this ancient calamity, breeding grief, heaped up through the dread ages, has been pinned into my body; melted from it by the sun's burning, drops fall that drip without cease upon the rocks of the Caucasus.} It seems, then, we can scarcely call a man so afflicted anything but wretched; and if wretched now, then pain is surely an evil. — So far, in fact, you are arguing my own case for me; but I will see to that presently.

\ciceroSection{26} Meanwhile, where do those verses come from? I do not recognize them. — I will tell you, by Hercules; for it is a fair question. Do you not see how much leisure I have? — What of it? — You were often, I suppose, when you were at Athens, in the schools of the philosophers. — Indeed I was, and gladly. — Then you noticed that, though no one at that time was especially copious, they nonetheless wove verses into their discourse. — Yes, and a great many, from Dionysius the Stoic. — You speak rightly. But he did it as if reciting by rote, with no discrimination, no elegance: Philo would bring in his own measured lines, well-chosen poems, and set them aptly in place. And so, once I came to love this kind of declamation — an old man's habit, if you like — I make eager use, for my part, of our own poets; but wherever they fall short, I have translated much from the Greek, so that Latin discourse should lack no ornament in this kind of argument. But do you see what mischief the poets bring?

\ciceroSection{27} They bring on the bravest men lamenting, they soften our spirits, and then they are so sweet that they are not merely read but learned by heart. And so, when a bad upbringing at home and a sheltered, pampered life have been joined by the poets as well, they sap the whole sinew of virtue. Rightly, then, are the poets cast out by Plato from the state he imagined, when he was searching out the best morals and the best constitution. But we, schooled of course by Greece, read these things and learn them by heart from boyhood, and count this our liberal education and culture. Yet why are we angry at the poets?

\ciceroSection{28} Teachers of virtue have been found — philosophers — who declared that the greatest evil is pain. And yet you, young man, when a little while ago you said this was your view, and I asked you whether pain was a greater evil even than disgrace, you abandoned that view at a word. Put the same question to Epicurus: he will say that moderate pain is a greater evil than the utmost disgrace; for in disgrace itself, he holds, there is no evil at all, unless pains follow upon it. What pain, then, follows Epicurus, when he says this very thing — that the greatest evil is pain? No disgrace greater than this do I look for from a philosopher. And so you gave me answer enough when you replied that disgrace seemed to you a greater evil than pain. For if you will only hold to this, you will understand how firmly pain must be resisted; and the question is not so much whether pain is an evil, as how the spirit is to be steeled to bear pain. The Stoics spin out their little arguments for why it is not an evil —

\ciceroSection{29} as if the labor were over a word and not over the thing. Why do you cheat me, Zeno? For when you flatly deny that to be an evil which seems to me horrible, I am caught, and I want to know how that which I judge most wretched of all is not even an evil. Nothing is evil, he says, except what is base and vicious. You return to your quibbles; for you do not take away the thing that was tormenting me. I know that pain is not moral wickedness; stop teaching me that. Teach me this instead: that whether I am in pain or not in pain makes no difference. It makes no difference at all, he says — to living happily, at least, which rests in virtue alone; but it is nonetheless to be rejected. Why? It is harsh, against nature, hard to endure, grim, cruel.

\ciceroSection{30} Here is a wealth of words — to be able to call by so many names the one thing we all call by a single word: evil. You define pain for me; you do not abolish it, when you call it harsh, against nature, scarcely to be borne and endured; and you do not lie. But it was not fitting, while glorying in your words, to give way in fact. So long as nothing is good but what is honorable, nothing evil but what is base — that is a wish, not a teaching. The better and truer doctrine is this: that all the things nature recoils from are among evils, all the things she takes to herself among goods. Once this is laid down and the wrangling over words is set aside, that other thing which these men rightly embrace will stand out so far above all else — what we call the honorable, the right, the becoming, which we sometimes embrace under the name of virtue — that everything else besides, all that is counted among the goods of the body and of fortune, will seem trivial and slight, and so no evil at all; and not even all of them gathered into one place would bear comparison with the evil of disgrace.

\ciceroSection{31} Therefore, if — as you granted at the outset — disgrace is worse than pain, then pain is plainly nothing. For so long as it seems to you base and unworthy of a man to groan, to wail, to lament, to be broken, to be unmanned by pain — so long as honor, dignity, and decency are present, and you, fixing your eyes upon them, hold yourself in check — pain will surely yield to virtue and lose its force under the spirit's resolve. For either there is no virtue at all, or all pain is to be despised. Will you have prudence, without which no virtue can even be conceived? What then? Will it suffer you to do anything that gains nothing and labors in vain? Or will temperance allow you to do anything without measure? Or can justice be cultivated by a man who, under the force of pain, betrays what was entrusted to him, gives up his confederates, abandons his many duties? And courage, with its companions — greatness of spirit, gravity, endurance, contempt for human affairs — how will you answer to them?

\ciceroSection{32} When you are struck down and prostrate, deploring your lot in a voice of lamentation, will you hear the words: "What a brave man!"? No — a man in that state no one would even call a man. Courage, then, must be given up, or pain must be buried. Do you grasp this much, then — that if you lose something from your Corinthian bronzes, you can still have the rest of your furnishings safe, but that virtue, if you lose a single one — though virtue cannot in fact be lost, yet if you once confess you lack a single one, you will have none at all?

\ciceroSection{33} Can you, then, call a man brave, or great of spirit, or enduring, or grave, or contemptuous of human things — or call that Philoctetes such — ? No, I would rather take leave of him; but he, at least, is no brave man, who \textit{lies in his dank shelter, which, echoing back his wailing, his complaint, his groans, his roaring, returns the mournful cries from its own silence.} I do not deny that pain is pain — for why else would courage be wanted? — but I say it is to be crushed by endurance, if only there is such a thing as endurance; if there is none, why do we dress up philosophy, or why do we glory in its name? Pain pricks — let it even stab: if you are unarmed, offer your throat; but if you are guarded by the arms of Vulcan, that is, by courage, resist. For unless you do, this guardian of your dignity will leave you and abandon you.

\ciceroSection{34} The laws of the Cretans — laid down whether by Jupiter or by Minos on the authority of Jupiter, as the poets relate — and likewise the labors of Lycurgus, train the young with hardship: by hunting, by running, by hunger and thirst, by cold and heat. At Sparta, indeed, boys are received at the altar with such a scourging that much blood flows from their flesh — sometimes even, as I heard when I was there, to the point of death; yet not one of them ever cried out, not one so much as groaned. What then? Boys can do this, and grown men will not be able to? Custom prevails, and reason will not prevail?

\ciceroSection{35} There is some difference between toil and pain. They are close neighbors, to be sure, yet they differ. Toil is a kind of performance, by mind or body, of some heavier work and duty; pain, on the other hand, is a harsh motion in the body, foreign to the senses. These two the Greeks — whose tongue is supposed to be richer than ours — call by a single name. And so they call hardworking men "lovers" or rather "devotees" of pain; we, more fittingly, call them laborious: for it is one thing to toil, another to feel pain. O Greece, sometimes poor in the very words you always think you have in abundance! It is one thing, I say, to feel pain, another to toil. When Gaius Marius had his varicose veins cut, he felt pain; when in great heat he led his column on the march, he toiled. And yet there is between these a certain likeness, for the habit of toil makes the endurance of pain easier.

\ciceroSection{36} And so those who shaped the constitutions of Greece wanted the bodies of the young hardened by toil; and the Spartans extended this even to their women, who in other cities are hidden away in the shade of their walls in the softest manner of living. The Spartans, by contrast, wanted nothing of that sort to hold for the daughters of Laconia, \textit{for whom the wrestling-ground, the Eurotas, sun, dust, toil, and the discipline of soldiering are more a study than barbarous fertility.} Therefore in these laborious exercises pain too breaks in from time to time: they are shoved, struck, thrown down, they fall; and the toil itself draws a kind of callus over pain.

\ciceroSection{37} As for soldiering — I mean ours, not the Spartans', whose line advances to measure and pipe, and no exhortation is given them without anapaestic feet — first, you see where our armies get their name; then, what toil there is, how great, in the column: to carry rations for more than half a month, to carry whatever they want for use, to carry the stake (for the shield, sword, and helmet our soldiers count as no more a burden than their shoulders, arms, and hands: they say the soldier's arms are his limbs; and indeed they are borne so aptly that, if the need arises, they can throw off their burdens and fight with their weapons unencumbered, as with limbs). What of this — the drilling of the legions? What of that running, charging, shouting — how great a toil it is! From this comes the spirit, ready for wounds in battle. Bring up an unpracticed soldier with the same spirit:

\ciceroSection{38} he will look like a woman. Why is there so great a difference between a new army and a veteran one, as we have learned by experience? The age of the recruits is generally better; but to bear toil, to despise a wound — habit teaches that. And what is more, we often see the wounded carried off the field, and the raw, unpracticed man, however slight the blow, raising the most shameful wails; while the practiced veteran, braver for that very reason, asks only for a doctor to bind him up: \textit{O Patroclus,} he says, \textit{coming to you I seek your help and your hands, before I meet the foul death dealt me by an enemy's hand … nor can the flowing blood by any means be stanched, if by your wisdom death can the more be turned aside. For the sons of Aesculapius fill the colonnades with their wounded; one cannot get near.} — Surely this man, at least, is Eurypylus. What a practiced fellow!

\ciceroSection{39} And no less practiced is the other: where so much grief is being recounted, see how unweepingly he answers, and even brings a reason why he should bear it with an even mind: \textit{He who prepares destruction for another should know it is prepared for himself, that he shares an equal ruin.} Patroclus, I suppose, will lead him off to lay him on a couch and bind his wound — if he were a man at all; but I saw nothing less. For he asks what has happened: \textit{Speak out, speak out, how the Argive cause holds up in battle.} — Words cannot tell as much as the deeds command of toil. Rest, then, and bind your wound. Even if Eurypylus could, Aesopus could not: \textit{When Hector's fortune, turned against us, made our keen line waver …} and the rest he sets out though in pain; for so unbridled is military glory in a brave man. Will a veteran soldier, then, be able to do this, and a learned and wise man not be able?

\ciceroSection{40} He, indeed, will do it better — and not a little better at that. But so far I am speaking of the habit of training, not yet of reason and wisdom. Old women often bear hunger for two or three days; take away an athlete's food for a single day, and he will implore Jupiter — Jupiter of Olympia, the very one for whom he trains — and cry out that he cannot bear it. Great is the force of habit: hunters spend the night out in the snow on the mountains; the Indians let themselves be burned; boxers, battered with the gauntlet, do not so much as groan.

\ciceroSection{41} But why these men, who reckon a victory at the Olympic games the equal of the old consulship? Gladiators — either ruined men or barbarians — what blows they endure! How those who have been well trained would rather take a blow than shamefully avoid it! How often it is plain that they want nothing more than to satisfy their owner or the crowd! Even when finished off with wounds they send to their owners to ask what they wish: if they have given satisfaction, they are content to go down. What ordinary gladiator ever groaned, ever changed his expression? Which of them, not merely standing but even down, ever lay shamefully? Which, once down and ordered to take the steel, ever drew in his neck? So great is the power of training, of discipline, of habit. So a Samnite can do this — a filthy man, worthy of that life and that place — and shall a man born for glory have any part of his soul so soft that he cannot brace it by discipline and reason? The spectacle of the gladiators is wont to seem cruel and inhuman to some, and perhaps it is so as it is now conducted; but when it was the guilty who fought it out with the steel, there could perhaps be many a stronger lesson for the ears, but for the eyes none stronger, against pain and death.

\ciceroSection{42} I have spoken of training and habit and rehearsal. Come now, let us look at reason, unless you have something to say to all this. — That I should interrupt you? I would not even wish it; so does your discourse lead me to belief. — Whether, then, it be an evil to feel pain or not, let the Stoics see to it, who, by certain twisted little syllogisms and tiny conclusions that do not reach down to the senses, want it proved that pain is no evil. I, for my part, hold that the thing, whatever it is, is not as great as it appears, and I say that men are moved more violently by a false vision and image of it, and that the whole of its pain is bearable. Where, then, shall I begin? Or shall I touch briefly again on the same things I said just now, so that my discourse may go on the more easily to greater length?

\ciceroSection{43} Among all men, then, this is agreed — not learned men only but the unlearned too — that to bear pain with endurance is the mark of brave and great-souled and patient men who rise above the merely human; nor was there ever anyone who did not think a man who bore it so deserved praise. What is asked of the brave, then, and is praised when it is done — to dread its coming, or not to bear its presence, is that not shameful? And yet consider: since all right dispositions of the soul are called virtues, this is not properly the name of all of them, but they have all been named from the one that excelled the rest. For virtue (\textit{virtus}) was named from "man" (\textit{vir}); and the property most peculiar to a man is courage (\textit{fortitudo}), whose two greatest functions are the contempt of death and of pain. We must use these, then, if we wish to be possessed of virtue, or rather if we wish to be men, since virtue borrowed its name from men. Perhaps you will ask how — and rightly: for it is just such a medicine that philosophy professes.

\ciceroSection{44} Up comes Epicurus, a man by no means bad, or rather a thoroughly good one; he advises as much as he understands. "Disregard pain," he says. Who says this? The same man who calls pain the greatest evil? Hardly consistent enough. Let us listen. "If the pain is at its greatest," he says, "it must be brief." Run that past me again! For I do not quite understand what you mean by "greatest," what by "brief." "Greatest" is that than which nothing is higher; "brief" that than which nothing is briefer. I despise the magnitude of a pain from which the brevity of time will release me almost before it has come. But what if the pain is as great as Philoctetes'? "It seems to me a fine and great pain, certainly, but still not the greatest; for nothing is in pain but his foot: the eyes can hurt, the head, the sides, the lungs can hurt, all of him can hurt; it is far, then, from the greatest pain." Therefore, he says, a long-lasting pain has more pleasure in it than annoyance.

\ciceroSection{45} Now I cannot say that so great a man understands nothing, but I think we are being mocked by him. The greatest pain — and I call it greatest even if another is greater by ten atoms — I deny is necessarily brief, and I can name many good men who for many years are tortured by the most severe pains of gout. But the clever fellow never fixes the limit either of magnitude or of duration, so that I might know what he means by "greatest" in pain, what by "brief" in time. Let us pass him by, then, as one who says nothing at all, and force him to confess that the remedies for pain are not to be sought from a man who has called pain the greatest of all evils, however brave a little face he may put on amid his own colic and his strangury. The medicine must therefore be sought elsewhere — and most of all, if we are seeking what is most consistent, from those for whom what is honorable is the greatest good, what is shameful the greatest evil. In their presence you will surely not dare to groan and toss about;

\ciceroSection{46} for through their voice Virtue herself will speak with you: "You — when you have seen boys at Sparta, young men at Olympia, barbarians in the arena, taking the heaviest blows and bearing them in silence — if some pain happens to pluck at you, will you cry out like a woman, and not bear it steadily and calmly?" It cannot be; nature does not allow it. I hear you. Boys bear it, led by glory; others bear it from shame, many from fear — and yet are we to fear that nature does not allow what is endured by so many and in so many places? Nature not only allows it but even demands it: for she has nothing more excellent, nothing she more desires, than what is honorable, than praise, than dignity, than honor. By these several names I mean to declare one single thing, but I use several to mean it as fully as I can. And what I mean is that the thing far best for a man is what is to be chosen for its own sake, springing from virtue or set in virtue itself, praiseworthy of its own accord — which I would sooner call the only good than not the greatest good. And as with the honorable, so the opposite holds of the shameful: nothing so foul, nothing so to be spurned, nothing more unworthy of a man.

\ciceroSection{47} And if you are persuaded of this — for at the outset you said that there seemed to you more evil in disgrace than in pain — it remains for you to command yourself. Though how this is said, I do not quite know. As though we were two, one to command and one to obey! And yet it is not unaptly said. For the soul is parceled into two parts, one of which shares in reason, the other lacks it. When, therefore, the rule is given that we should command ourselves, this is the rule: that reason should curb recklessness. There is in nearly everyone's soul a certain softness by nature, something cast down, lowly, in a way unstrung and languid. If there were nothing else, nothing would be more deformed than man. But ready at hand is reason, the mistress and queen of all, which, striving by its own power and advancing further, becomes perfected virtue. That this should rule that part of the soul which ought to obey — this is what a man must see to. In what way? you will say.

\ciceroSection{48} As a master rules a slave, or a commander a soldier, or a parent a child. If that part of the soul which I called soft conducts itself most shamefully, if it gives itself over to womanish wailing and tears, let it be bound and shackled by the watch of friends and kinsmen; for we often see those broken by shame whom no reasoning could overcome. These, then, like household slaves, are to be held all but in chains and under guard; but those who will be firmer, though not yet of the strongest sort, must be put in mind, like good soldiers recalled to their posts, to guard their dignity. In the \textit{Niptra} that wisest of the Greeks does not lament too much in his wound, or rather laments with measure: \textit{Go softly,} he says, \textit{and with calm effort, lest the jolting seize a greater pain} (Pacuvius does this better than Sophocles;

\ciceroSection{49} for in Sophocles Ulysses laments quite tearfully over his wound); yet even to him, groaning lightly, those very men who carry the wounded hero, looking to the gravity of his person, do not hesitate to say: \textit{You too, Ulysses — though we see you grievously struck — are almost too soft of soul, you who are used to passing your life under arms.} The prudent poet understands that the habit of bearing pain is a teacher not to be despised.

\ciceroSection{50} And he, not immoderately, in great pain: \textit{Hold me, hold! The sore overwhelms me; uncover it! Alas, wretched me — I am racked.} He begins to give way, then at once breaks off: \textit{Cover it, get away now, now! Let me be! For by handling and shaking it you make the savage pain greater.} Do you see how the pain fell silent — not the body's, soothed, but the soul's, chastened? And so at the very end of the \textit{Niptra} he rebukes others too, and that as he is dying: \textit{To complain of adverse fortune is fitting, not to lament; that is a man's office — weeping was added by a womanish nature.}

\ciceroSection{51} The softer part of this man's soul obeyed reason as a soldier who knows shame obeys a stern commander. But the man in whom wisdom shall be perfected — one whom thus far, at least, none of us has seen, though the philosophers' doctrines set out what he will be like, if ever he comes to be — he, then, or that reason which will be in him perfect and complete, will so command the lower part as a just father commands worthy sons; with a nod he will accomplish what he wishes, with no toil, no annoyance; he will rouse himself, raise himself up, equip himself, arm himself, so as to stand against pain as against an enemy. What are these arms? Tension, resolve, and inward speech, when he says to himself:

\ciceroSection{52} "Let nothing be shameful, nothing slack, nothing unmanly." Let noble images pass before the soul. Set up Zeno of Elea, who endured everything sooner than betray his fellow conspirators in the plot to destroy the tyrant. Think of Anaxarchus, the disciple of Democritus, who, when he had fallen into the hands of King Timocreon on Cyprus, neither begged off any kind of torture nor refused it. Callanus the Indian, an untaught barbarian born at the foot of the Caucasus, was burned alive by his own choice; while we, if a foot aches, if a tooth — but suppose the whole body in pain — cannot bear it. For there is a certain belief, soft and frivolous — and no more so in pain than the very same belief in pleasure — under which, melting and dissolving in our own softness, we cannot bear the sting of a bee without an outcry.

\ciceroSection{53} But Gaius Marius now — a man of the countryside, yet plainly a man — when he was being operated on, as I said above, at first forbade them to tie him down, and no one before Marius is said to have been cut while left unbound. Why, then, did others follow him afterward? His example carried the day. Do you see, then, that the evil belongs to belief, not to nature? And yet that same Marius showed that the bite of the pain was sharp, for he did not offer the second leg. Thus he both bore the pain like a man and, like a man, refused to bear more of it without compelling cause. The whole thing, then, lies in this: that you command yourself. And I have shown what kind of command it is. This very reflection — on what is most worthy of patience, of fortitude, of greatness of soul — not only keeps the soul in check but somehow even makes the pain itself milder.

\ciceroSection{54} For as it happens in battle that the coward, the timid soldier, the moment he has caught sight of the enemy, throws away his shield and runs as fast as he can, and for that reason sometimes perishes though his body is whole, while nothing of the kind befalls the man who held his ground — so those who cannot bear the look of pain throw themselves down and lie there, beaten and breathless; but those who stood firm come away the victors more often than not. For the soul has certain likenesses to the body. As burdens are borne more easily by bodies that are braced, but crush bodies that are slack, so the soul by its own tension throws off the full pressure of every weight,

\ciceroSection{55} but by slackness is so weighed down that it cannot lift itself up. And if we seek the truth, in the pursuit of every duty the soul's tension must be brought to bear; it is the one guardian, so to speak, of duty. But this same thing must above all be watched in pain: that we do nothing abjectly, nothing timidly, nothing like a coward, nothing slavishly or like a woman; and first of all, let that famous cry of Philoctetes be repudiated and cast out. To groan is sometimes granted to a man, and that rarely; to wail is not granted even to a woman. And this, surely, is the lamentation that the Twelve Tables forbade to be used at funerals.

\ciceroSection{56} Nor does the brave and wise man ever even so much as groan, unless perhaps to brace himself toward firmness, as runners in the stadium cry out as loud as they can. Athletes do the same when they train; and boxers, even as they strike their opponent, groan as they swing their gloves — not because they are in pain or give way in spirit, but because in pouring out the voice the whole body is braced, and the blow comes home harder. What of the man who wants to shout louder? Is it enough for him to brace his sides, his throat, his tongue, from which we see the voice drawn out and poured forth? With his whole body, and with every nail, as the saying goes, he serves the straining of his voice.

\ciceroSection{57} By Hercules, I myself saw Marcus Antonius touch the ground with his knee while he was speaking for himself with all his force under the Lex Varia. For as catapults of stones and the other engines for hurling missiles send forth heavier shots the more they are drawn taut and pulled back hard, so the voice, so the runner's pace, so the blow, falls the heavier the more strenuously it is launched. And since the force of this tension is so great, if a groan in pain will serve to steady the soul, we shall use it; but if it is a wailing groan, weak, abject, tearful, the man who has surrendered to it I could scarcely call a man. Even if that groan brought some relief, we might still consider what befits a brave and high-spirited man; but when it lessens the pain not at all, why should we choose to be shameful to no purpose? For what is more shameful in a man than weeping like a woman?

\ciceroSection{58} And this precept, which is given concerning pain, reaches more widely. For against all things, not pain alone, the same tension of soul must be set in resistance. Anger flares up, desire is roused: we must take refuge in the same citadel, we must take up the same arms. But since we are speaking of pain, let us leave those aside. Toward bearing pain calmly and steadily, then, it does the most good to consider with the whole heart, as the saying goes, how noble a thing it is. For we are by nature, as I said before — and it must be said more than once — most eager for nobility and most desirous of it; and if we have so much as glimpsed some light of it, there is nothing we are not ready to bear and to suffer through in order to attain it. From this onrush and drive of souls toward true praise and nobility come those dangers that men face in battle: brave men do not feel their wounds in the line, or feel them and prefer to die rather than be moved so much as one step from their place of honor.

\ciceroSection{59} The Decii saw the flashing swords of the enemy as they hurled themselves into their ranks. For them the nobility and glory of their death lightened all fear of wounds. Do you suppose Epaminondas groaned when he felt his life flowing out together with his blood? For he was leaving his country in command over the Spartans, when he had received it in servitude to them. These are the consolations, these the salves, of the greatest pains.

\ciceroSection{60} You will say: what of times of peace, what of home, what of the sickbed? You call me back to the philosophers, who do not often go out into the battle line. Among them one truly frivolous man, Dionysius of Heraclea, having learned from Zeno to be brave, was unlearned by pain. For when he was suffering from his kidneys, in the very midst of his wailing he kept crying out that what he had once held about pain was false. And when his fellow pupil Cleanthes asked him what reasoning had drawn him from his opinion, he answered: "Because, if after giving so little effort to philosophy I still could not bear pain, that would be argument enough that pain is an evil. But I have spent very many years in philosophy, and I cannot bear it. Therefore pain is an evil." Then Cleanthes, they say, struck the ground with his foot and quoted a line from the \textit{Epigoni}: \textit{Do you hear this, Amphiaraus, hidden beneath the earth?} He meant Zeno, from whom he grieved to see the man degenerate.

\ciceroSection{61} But not so our friend Posidonius, whom I myself often saw — and I will tell what Pompey used to relate: that when he had come to Rhodes on his way back from Syria, he wished to hear Posidonius; but when he learned that the man was gravely ill, his joints in violent pain, he still wished to call on the most distinguished of philosophers. And when he had seen and greeted him and attended him with words of honor, and had said how grieved he was that he could not hear him, the other replied: "But you can; nor will I allow pain of the body to make so great a man come to me in vain." And so Pompey would tell how, lying there, he argued gravely and at length on this very point — that nothing is good except what is noble — and how, as the torches of pain were brought close to him, he said again and again:

\ciceroSection{62} "You accomplish nothing, pain! Troublesome though you are, I will never confess that you are an evil." And in general all labors that are illustrious and ennobled become at once even bearable. Do we not see, among peoples where there is great honor in those games called gymnic, that no pain is shunned by those who go down into that contest? And among those where the praise of hunting and riding flourishes, that those who seek it flee no pain? What shall I say of our own ambitions, of our craving for office? Through what flame have they not run, those men who once gathered these honors vote by single vote? And so Africanus always had Xenophon, the disciple of Socrates, in his hands, and praised above all that saying of his: that the same labors do not weigh equally on a commander and a soldier, because the honor of command itself makes the labor lighter.

\ciceroSection{63} But for all that, it comes about that even among the unwise crowd a belief in nobility holds sway, though they cannot see it in itself. And so they are moved by reputation and by the judgment of the multitude, supposing that to be noble which is praised by most. As for you, even if you are in the eye of the multitude, I would still not have you stand by its judgment, nor think that finest which it thinks so. You must use your own judgment. If you please yourself by approving what is right, then you will have conquered not only yourself — which is what I urged a little while ago — but everyone and everything.

\ciceroSection{64} Set this before yourself, then: that greatness of soul, and a kind of heightening of the soul to its utmost, which shows itself most in the contempt and disdain of pain, is the one finest thing of all — and the finer if it does without the crowd and, while catching at no applause, takes delight in itself alone. Indeed, to me everything seems more praiseworthy that is done without display and without the people for witness — not that the people are to be shunned, for all good deeds want to be set in the light, but still no theater is greater for virtue than one's own conscience.

\ciceroSection{65} And let us reflect on this above all: that this endurance of pain, which I have said so often must be steadied by the soul's tension, should show itself even and uniform in every kind of trial. For often many a man who, whether from desire of victory or of glory, or even to hold fast his right and his liberty, has taken wounds bravely and borne them — that same man, once the tension is dropped, cannot bear the pain of illness; for the pain he had borne so easily he had borne not by reason or by wisdom, but by zeal rather and by love of glory. And so certain barbarous and savage peoples can fight it out with the sword most fiercely, yet cannot fall ill like men. The Greeks, on the other hand, men not brave enough, but wise as far as men's capacity goes, cannot look an enemy in the face, yet bear sickness with patience and humanity. But the Cimbri and the Celtiberians exult in battle and lament in illness. For nothing can be uniform that does not proceed from settled reason.

\ciceroSection{66} But when you see that those who are led on by zeal or by belief are not broken in the pursuit and attainment of their aim by pain, you should judge either that pain is not an evil, or — even if, granting that whatever is harsh and contrary to nature may be called an evil, it still be so slight that it is so overwhelmed by virtue as to appear nowhere at all. Meditate on these things, I beg you, day and night. For this reasoning will spread more widely and take up a place considerably larger than the single matter of pain. For if we do all things for the sake of fleeing baseness and attaining nobility, we shall be free to despise not only the goads of pain but even the thunderbolts of fortune — especially since that refuge from yesterday's discussion stands ready.

\ciceroSection{67} For as, if some god were to say to a man at sea, with pirates in pursuit, "Throw yourself from the ship; one is at hand to catch you — either a dolphin, as it caught Arion of Methymna, or those Neptunian horses of Pelops, said to have swept their chariots aloft through the waves; they will catch you and carry you wherever you wish" — and he would lay aside all fear: so, when harsh and hateful pains press hard, if they are so great that they are not to be borne, you see where one must flee for refuge. This much, more or less, I thought should be said at this time. But perhaps you remain of the same mind. — Not in the least; and within these two days I hope I have been freed of my fear of the two things I most dreaded. — Tomorrow, then, at the water-clock; for so we agreed, and I see this cannot but be owed to you. — Just so; and that one before noon, this one at this same hour. — So we shall do, and indulge your excellent pursuits.

\ciceroBook{3}

\ciceroSection{1} What can the reason be, Brutus, that I should suppose — given that we are composed of soul and body — an art was sought out for tending and guarding the body, and its usefulness consecrated to the discovery of the immortal gods, while the medicine of the soul was neither so longed for before it was found, nor so cultivated after it was known, nor welcome and approved by so many, but to many even suspect and hateful? Is it because we judge the heaviness and pain of the body with the soul, but do not feel the soul's disease with the body? So it comes about that the soul passes judgment on itself precisely when the very thing by which judgment is made is sick.

\ciceroSection{2} If nature had brought us forth such that we could gaze upon and discern her very self, and with her as our guide — the best of guides — complete the course of life, there would surely be no reason for anyone to look for system and teaching. As it is, she has given us small sparks of fire, which, swiftly corrupted by bad habits and bad opinions, we so smother that nowhere does the light of nature shine through. For there are inborn in our natures the seeds of the virtues, and if these were allowed to grow, nature herself would lead us to the happy life. But as it is, the moment we are brought forth into the light and taken up, we live at once amid every kind of depravity and the utmost perversity of opinions, so that we seem almost to have drunk in error with our nurse's milk. And when we are restored to our parents and then handed over to teachers, we are then steeped in such varied errors that truth gives way to emptiness, and nature herself to opinion grown firm.

\ciceroSection{3} Then too the poets come in. Putting forward as they do a great show of learning and wisdom, they are heard, read, learned by heart, and cling deep within our minds. And when to these is added, as a kind of supreme teacher, the people, and the whole multitude that on every side agrees in vice, then we are thoroughly stained with the depravity of opinion and fall away from nature — so that those men seem to us to have seen nature's force best who have judged that nothing is better for a human being, nothing more to be sought after, nothing more excellent than offices, commands, and popular glory. Toward this every best man is borne, and while he seeks that true honor which alone nature most deeply desires, he busies himself in utter emptiness and chases not the eminent form of virtue but a shadowed image of glory. For glory is something solid and clearly stamped, not a shadow: it is the agreeing praise of good men, the uncorrupted voice of those who judge well of outstanding virtue; it answers back to virtue as an echo, and because it is for the most part the companion of right action, it is not to be spurned by good men.

\ciceroSection{4} But that other thing, which wishes to pass for an imitation of glory — reckless and ill-considered, and for the most part a praiser of faults and vices — popular reputation, by a pretense of honor corrupts the form and beauty of the real thing. By this blindness, when men have longed for certain splendid things too, and yet did not know what they were or where or of what kind, some have utterly overturned their own cities, others have perished themselves. And these men, indeed, in seeking the best, are deceived not so much by their will as by mistaking their road. What of those who are carried away by the craving for money, by the desire for pleasures, and whose minds are so disturbed that they are not far from madness — which is the lot of all the foolish? Is no cure to be applied to them? Is it because the soul's sicknesses do less harm than the body's, or because, while bodies can be cured, there is no medicine for souls?

\ciceroSection{5} But in fact the soul's diseases are both more destructive and more numerous than the body's. For these very diseases are hateful precisely because they belong to the soul and trouble it; and the sick soul, as Ennius says, \textit{always wanders, and can neither bear nor endure, and never ceases to crave}. By these two diseases alone, to leave the others aside — distress and craving — what diseases in the body, after all, could be heavier? And how can it be allowed that the soul cannot heal itself, when the soul is the very inventor of the body's medicine; when, for the healing of bodies, the bodies themselves and nature count for much, and not all who have let themselves be treated immediately recover as well — whereas souls that have wished to be healed and have obeyed the precepts of the wise are healed without any doubt at all?

\ciceroSection{6} There is, surely, a medicine of the soul: philosophy. Its help is not to be sought from outside, as in diseases of the body; rather we must labor with all our resources and strength to be able to heal ourselves. As for philosophy in general — how greatly it ought to be sought and cultivated — enough, I think, has been said in the \textit{Hortensius}. And on the greatest matters I have left off hardly at all from arguing and writing ever since. In these books are set out the things that were debated by me with my friends at my place at Tusculum. But since in the two earlier books we spoke of death and of pain, the third day of discussion will make up this third volume.

\ciceroSection{7} For when we had gone down into our Academy, the day already declining toward afternoon, I asked one of those present to propose a subject for discussion. The matter then proceeded thus. — It seems to me that distress falls upon the wise man. — And do the other disturbances of the soul as well — fears, desires, fits of anger? These are roughly of the kind the Greeks call \textit{[Greek: pathe]}. I could call them "diseases," and that would be word for word, but it would not suit our usage. For to pity, to envy, to exult, to rejoice — all these the Greeks call diseases, motions of a soul not obedient to reason; while we, I think, would rightly call these same motions of an agitated soul "disturbances," but "diseases" not quite in keeping with custom — unless something else seems right to you. — No, I agree, in that way.

\ciceroSection{8} Do you think, then, that these things fall upon the wise man? — That is precisely my view. — Why, then, that vaunted wisdom is not to be valued highly, if it differs little from madness. — What? Does every commotion of the soul seem to you madness? — Not to me alone; but, as I am often given to wonder at, I find that our ancestors too held this view, many ages before Socrates, from whom flowed all this philosophy that concerns life and conduct. — In what way, then? — Because the word "madness" \textit{(insania)} signifies a sickness and disease of the mind — that is, an unsound and sick soul, which they called \textit{insania}, "unsoundness."

\ciceroSection{9} (All the disturbances of the soul, however, the philosophers call diseases, and they deny that any fool is free of these diseases. Now those who are diseased are not sound; and the souls of all the foolish are diseased: therefore all the foolish are mad.) For they held that the soundness of souls lay in a certain tranquillity and steadiness; a mind empty of these things they called \textit{insania}, on the ground that in a disturbed soul, as in a disturbed body, soundness cannot exist.

\ciceroSection{10} No less acute was this: that they named an affection of the soul which lacks the light of the mind \textit{amentia}, mindlessness, and the same thing \textit{dementia}. From which it must be understood that those who fixed these names on things felt the very thing that the Stoics, taking it from Socrates, have carefully retained — that all the foolish are unsound. For a soul that is in some disease — and these disturbed motions, as I said just now, the philosophers call diseases — is no more sound than a body that is in disease. So it comes about that wisdom is the soundness of the soul, but folly is a kind of unsoundness, which is madness and likewise mindlessness; and these things are marked out far better by Latin words than by Greek.

\ciceroSection{11} This will be found in many other places as well; but that another time — now, what presses. The whole of what we are inquiring after, then — what it is and of what kind — the force of the word itself declares. For since we necessarily understand those to be sound whose mind is disturbed by no motion, as if by a disease, those who are otherwise affected must necessarily be called unsound. And so nothing is better than what we have in the usage of the Latin tongue, when we say that those have "gone out of their own control" who are carried away unbridled either by desire or by anger — although anger itself is a part of desire; for thus it is defined: anger is the desire of taking revenge. Those, then, who are said to have gone out of their own control are said to be so because they are not under the control of the mind, to which nature has granted the rule over the whole soul. Why the Greeks call it \textit{[Greek: mania]} I could not easily say; yet we distinguish the thing itself better than they. For this madness, which, joined to folly, extends more widely, we mark off from frenzy. The Greeks mean to do so, but they are not equal to it in their word: what we call frenzy they call \textit{[Greek: melancholia]} — as though the mind were moved only by black bile, and not often by anger too severe, or by fear, or by grief; in this sense we say that Athamas, Alcmaeon, Ajax, and Orestes were frenzied. A man so affected the Twelve Tables forbid to be master of his own affairs; and so it is not written "if he be mad," but "if he be frenzied." For they judged that folly, though lacking steadiness — that is, soundness — could nonetheless look after a middling discharge of duties and the common, customary conduct of life; but frenzy they reckoned to be a blindness of the mind in all things. And although this seems to be greater than madness, yet it is of such a kind that frenzy can fall upon the wise man, while madness cannot. But this is another question; let us return to what we set out to do.

\ciceroSection{12} You said, I think, that distress seemed to you to fall upon the wise man. — And indeed that is my view. — Human enough, that you should think so. For we are not born of flint; rather there is in our souls something naturally tender and soft, which distress shakes as a storm shakes the sea. And not without reason that famous Crantor, who in our Academy was among the foremost in renown, said: "I do not at all agree with those who so greatly praise that I-know-not-what insensibility, which can neither exist nor ought to. May I not fall ill; but if I do," he said, "let feeling be there, whether something be cut or torn from the body. For that state of feeling no pain comes only at a great price — a brutishness in the soul, a stupor in the body."

\ciceroSection{13} But let us take care that this is not the talk of men who flatter our weakness and indulge our softness; let us instead dare not only to cut away the branches of our miseries, but to tear out all the fibers of the roots. Yet something will perhaps be left behind — so deep run the stocks of folly; but only what is necessary will be left. Hold this much for certain: unless the soul is healed — which cannot happen without philosophy — there will be no end of miseries. Therefore, since we have begun, let us hand ourselves over to it for treatment: we shall be healed, if we wish to be. And I will go further still: for I shall set out not only about distress — though that indeed first — but about every disturbance of the soul, as I have put it, every disease, as the Greeks would have it. And first, if you please, let us proceed in the manner of the Stoics, who are wont to bind their arguments tightly; then we shall range about in our own way.

\ciceroSection{14} He who is brave is also confident — since "self-assured" \textit{(confidens)}, by a bad habit of speaking, is counted a fault, though the word is drawn from "confiding," which belongs to praise. But he who is confident surely does not take fright; for confidence is at odds with fearing. Now upon whomever distress falls, upon him falls fear as well; for the very things whose presence puts us in distress, those same things, when they hang over us and approach, we fear. So it comes about that distress is at war with fortitude. It is likely, then, that whoever is liable to distress is liable also to fear, and indeed to a broken and dejected state of soul. And the man upon whom these fall is also liable to be enslaved, to confess himself beaten, if ever it should come to that. And whoever admits these must necessarily admit timidity and cowardice as well. But these do not fall upon the brave man: therefore neither does distress. But no one is wise unless he is brave: therefore distress will not fall upon the wise man.

\ciceroSection{15} Moreover, it must be that whoever is brave is also great of soul; whoever is great of soul, unconquerable; whoever is unconquerable, looks down on human affairs and reckons them as set beneath himself. But no one can look down on those things on account of which he can be afflicted with distress; from which it follows that the brave man is never afflicted with distress. And all the wise are brave: distress, therefore, does not befall the wise man. And just as an eye that is disturbed is not properly disposed to perform its office, and the remaining parts, or the whole body, when shifted from their standing, fail in their function and their office, so a soul that is disturbed is not fit to carry out its office. Now the office of the soul is to make good use of reason; and the soul of the wise man is always so disposed that it makes the best use of reason; it is therefore never disturbed. But distress is a disturbance of the soul: the wise man, therefore, will always be free of it.

\ciceroSection{16} This too is likely: that whoever is temperate — whom the Greeks call \textit{[Greek: sōphrōn]}, and that virtue \textit{[Greek: sōphrosynē]}, which I myself am accustomed to call now temperance, now moderation, and sometimes even modesty — though I am not sure whether that virtue might rightly be called \textit{frugalitas}, thrift, which carries a narrower force among the Greeks, who call thrifty men \textit{[Greek: chrēsimous]}, that is, merely useful; but our word is broader; for all self-restraint, all harmlessness (which among the Greeks has no name in common use, though it could have \textit{[Greek: ablabeia]}; for harmlessness is a disposition of the soul such as harms no one) — thrift, then, embraces the other virtues as well. And were it not so great, were it confined within the narrow limits in which most people suppose it confined, the surname of Lucius Piso would never have been so highly praised.

\ciceroSection{17} But since neither the man who out of fear has abandoned his post, which is cowardice, nor the man who out of greed has failed to return a deposit entrusted to him in secret, which is injustice, nor the man who out of recklessness has managed an affair badly, which is folly, is wont to be called thrifty, on that ground thrift has gathered into itself three virtues, courage, justice, and prudence (though this indeed is common to the virtues, for they are all bound and yoked together among themselves): thrift, then, is itself the remaining and fourth virtue. For its proper work seems to be to govern and calm the motions of the appetite, and, ever opposed to desire, to preserve a measured constancy in everything. The vice contrary to it is called depravity.

\ciceroSection{18} Thrift, \textit{frugalitas}, I take it, comes from \textit{frux}, the fruit of the earth, than which nothing is better; depravity, \textit{nequitia}, from this (though this perhaps will be rather forced — but let us try; let us be thought to have played, if it comes to nothing) — from the fact that in such a man there is naught of worth, \textit{nequicquam}, from which he is likewise called a man of nothing. The man who is thrifty, then — or, if you prefer, moderate and temperate — must be constant; and whoever is constant, untroubled; whoever is untroubled, free of every disturbance, and therefore of distress too. And these belong to the wise man: distress, then, will be absent from the wise man. And so it is not without wit that Dionysius of Heraclea argues against the lines in which Achilles complains in Homer, in this fashion, as I take it: \textit{And deep within, my heart swells up with bitter wrath, when I recall myself stripped of my honor and of all my glory.} Is a hand in good order when it is swollen, or is any other limb that is swollen and distended not in a faulty state?

\ciceroSection{19} So, then, a soul puffed up and swelling is in a faulty state. But the soul of the wise man is always free of fault, never swells, never grows distended; whereas the soul of an angry man is of just that kind: the wise man, therefore, is never angry. For if he is angry, he also craves; for it is proper to the angry man to crave to burn into the one by whom he thinks himself wronged the greatest possible pain. And whoever has craved this must, if he has attained it, rejoice greatly. From which it follows that he takes pleasure in another's misfortune; and since this does not befall the wise man, neither does it befall him to be angry. But if distress befell the wise man, anger would befall him too; since he is free of anger, he will be free of distress as well.

\ciceroSection{20} And indeed, if the wise man could fall into distress, he could fall into pity too, and into envying (I did not say envy, which exists when one is envied; but from envying, \textit{invidere}, "envying," \textit{invidentia}, can rightly be formed, so that we may escape the ambiguity of the word envy. That word is drawn from gazing too hard at another's fortune, as in the \textit{Melanippus}: \textit{Who, then, has envied the flower of my children?} It seems bad Latin, but Accius wrote splendidly; for as we say "to see" a thing, so "to envy" the flower is more correct than "to envy at" the flower. Custom forbids us this;

\ciceroSection{21} but the poet kept his own right and spoke more boldly) — both pity, then, and envy befall the same man. For whoever grieves at another's adversities grieves also at another's prosperity, as Theophrastus, lamenting the death of his comrade Callisthenes, is tormented by Alexander's prosperous fortunes, and so says that Callisthenes had fallen in with a man of the highest power and the highest fortune, but one who knew not how it befits a man to use prosperity. And just as pity is distress arising from another's adversities, so envying is distress arising from another's prosperity. Upon whomever, then, pity befalls, envying too befalls him; but envying does not befall the wise man: therefore neither does pity. And if the wise man were wont to take things hard, he would be wont to feel pity too. Distress, then, is absent from the wise man.

\ciceroSection{22} These things are stated by the Stoics in this way, and drawn to their conclusion rather too tightly. But sometimes they must be stated more amply and more at large; one ought, however, to use above all the doctrines of those who employ the boldest and, so to speak, the most manly reasoning and judgment. For the Peripatetics, our intimate friends, than whom nothing is richer, nothing more learned, nothing weightier, do not at all win my approval with their mean averages, whether of the soul's disturbances or of its diseases. For every evil, even a moderate one, is an evil; and what we are arguing is that in the wise man there be no evil at all. For just as the body, even if it is moderately sick, is not sound, so in the soul that moderation lacks soundness. And so our countrymen, splendidly, as in many other things, named trouble, anxiety, and anguish \textit{aegritudo}, distress, on account of its likeness to sick bodies.

\ciceroSection{23} By much the same word the Greeks call every disturbance of the soul; for they call it \textit{[Greek: pathos]}, that is, a disease, whatever is a turbid motion in the soul. We do better: for distress is most like sick bodies, but desire is not like sickness, nor is unbounded gladness, which is a pleasure of the soul exalted and exultant. Fear itself, too, is not very like a disease, although it is neighbor to distress; but, properly, just as sickness has its name in the body, so distress has, in the soul, a name not sundered from pain. The origin of this pain, then, must be unfolded by us, that is, the cause that produces distress in the soul as sickness in the body. For just as physicians think that, once the cause of a disease has been found, the cure has been found, so we, once the cause of distress has been discovered, shall discover the means of healing.

\ciceroSection{24} The whole cause, then, lies in opinion — and not the cause of distress only, but of all the other disturbances as well, which are four in kind, and more in their parts. For since every disturbance is a motion of the soul either devoid of reason, or spurning reason, or not obeying reason, and since that motion is roused in a twofold way, by the opinion of either good or evil, the four disturbances are evenly distributed. For two arise from the opinion of good: of these, the one, exultant pleasure, that is, gladness raised beyond measure, from the opinion of some great present good; the other, craving, which may rightly be called desire too, which is an unbounded reaching after a supposed great good, not heeding reason —

\ciceroSection{25} these two kinds, then, exultant pleasure and desire, are stirred up by the opinion of goods, as the remaining two, fear and distress, by the opinion of evils. For fear is the opinion of a great impending evil, and distress is the opinion of a great present evil, and indeed a fresh opinion of such an evil that it seems right to be anguished at it — that is, that the one who grieves supposes he ought to grieve. Against these disturbances, which folly looses upon the life of men and goads on like so many Furies, we must fight back with all our strength and all our resources, if we wish to pass tranquilly and peaceably through this span that has been granted to life. But the rest another time; now let us drive off distress, if we can. For let this be our purpose, since you said that distress seems to you to befall the wise man — which I in no way think; for it is a foul thing, wretched, detestable, to be fled with every exertion, under full sail, so to speak, and at the oar.

\ciceroSection{26} For what sort of man does that one seem to you — \textit{sprung from Tantalus, born of Pelops, who once, having seized in marriage Hippodamia from her father, King Oenomaus}—? He, surely, was the great-grandson of Jupiter. So utterly cast down, then, so utterly broken? \textit{Come not near me, strangers}, he says, \textit{stop right where you stand, lest my contagion or my shadow harm the good. So great is the power of crime that clings within my body.} Will you, Thyestes, condemn yourself and rob yourself of the light because of the power of another's crime? What of this? Do you not think that son of the Sun unworthy of his own father's light? \textit{The eyes have shrunk back, the body has wasted away with leanness, tears have eaten through the cheeks with bloodless moisture, neglect upon the chin makes the beard bristle hideous with filth, and the unshorn squalor blackens the rough breast.} These ills, O most foolish Aeetes, you yourself added to your own lot; they were not in those that chance had brought upon you, and that too when the evil had grown old, once the swelling of the soul had subsided — for distress, as I shall show, lies in the opinion of a fresh evil — but you grieve, clearly, from longing for your kingdom, not for your daughter. Her you hated, and rightly perhaps; your kingdom you could not bear to be without with an even mind. But it is a shameless grief, that of a man wasting himself in mourning because he is not allowed to lord it over the free.

\ciceroSection{27} Dionysius the tyrant, indeed, driven out of Syracuse, taught boys at Corinth: so utterly could he not do without command. And what was more shameless than Tarquin, who waged war against those who had refused to bear his arrogance? When he could not be restored to his kingdom by the arms either of the Veientes or of the Latins, he is said to have betaken himself to Cumae, and in that city to have been worn out by old age and distress. Do you, then, think this can happen to a wise man — that he be crushed by distress, that is, by wretchedness? For while every disturbance is wretchedness, distress is downright torture. Desire has its burning, exultant gladness its frivolity, fear its abjectness, but distress has graver things — wasting, racking, affliction, foulness; it tears the soul, eats it out, and utterly destroys it. Unless we strip this off so as to cast it away, we cannot be free of wretchedness.

\ciceroSection{28} And this much, at least, is clear: that distress arises when something has so appeared to us that some great evil seems to be present and to press upon us. Epicurus, however, holds that distress is by nature the opinion of evil, so that whoever fixes his gaze on some greater evil, if he supposes it has befallen him, is straightway in distress. The Cyrenaics think that distress is produced not by every evil, but by one unhoped for and unforeseen. And this indeed is no small thing toward increasing distress: for all things sudden seem the heavier. From this those lines too are justly praised: \textit{When I begot them, then I knew they would die, and for that very lot I raised them. And further, when I sent them to Troy to defend Greece, I knew that I was sending them into deadly war, not to a feast.}

\ciceroSection{29} This forethought, then, about the evils to come softens their arrival, when you have seen them coming long beforehand. And so the words Theseus speaks in Euripides are praised; for I may, as I often do, turn them into Latin: \textit{For since I had remembered hearing this from a learned man, I used to turn over future miseries within myself — a bitter death, or the grieving exile of banishment, or some other weight of evil I would always be rehearsing — so that, if any cruelty came down upon me by chance, no sudden onset of care should tear me unprepared.}

\ciceroSection{30} But what Theseus says he heard from a learned man, Euripides is speaking of himself. For he had been a pupil of Anaxagoras, who, when his son's death was reported to him, is said to have replied: I knew that I had fathered a mortal. This saying makes plain that such things are bitter to those who have not thought them through beforehand. So this much, at least, is not in doubt: that everything held to be evil is heavier when it comes unforeseen. And therefore, although it is not this one thing that produces the greatest distress, still, since the mind's foresight and preparation count for much in lessening pain, let a man ponder always all that can befall a human being. And surely this is that surpassing, that divine wisdom — to have grasped human affairs deeply and handled them thoroughly, to be astonished at nothing when it happens, and to judge, before it comes about, that nothing is impossible. \textit{For every man, precisely when his fortunes stand at their best, ought then above all to rehearse within himself how he is to bear adversity's hardship. Returning from abroad, let him always turn over in his mind dangers, losses, a son's wrongdoing, a wife's death, a daughter's illness — that these are the common lot, so that none of them may ever strike his mind as something new; and whatever turns out beyond his hope, let him count all of it so much gain.} So when Terence said this so aptly, having taken it from philosophy, shall we — from whose very springs it was drawn — not both say it better and hold to it more steadily?

\ciceroSection{31} For this is that ever-unchanging countenance which Xanthippe is said to have been in the habit of declaring belonged to her husband Socrates: she had always seen the same face on him as he left the house and as he returned. And it was no such brow as that of the famous old Marcus Crassus, who in his whole life, Lucilius says, laughed but once, but a calm and serene one; so we have been told. And rightly was his countenance always the same, since no change ever came over the mind from which it is fashioned. So I do indeed accept from the Cyrenaics these weapons against the blows and turns of chance, by which the assaults of misfortune, as they come on, are broken through long forethought; and at the same time I judge that this evil belongs to opinion, not to nature; for if it lay in the thing itself, why should things foreseen come lighter?

\ciceroSection{32} But there is something subtler to be said on these same matters, once we have first looked at the view of Epicurus. He holds it inevitable that all those who think themselves in evil should be in distress, whether that evil was foreseen and expected beforehand or has grown old upon them. For evils, he says, are not lessened by length of time, nor made lighter by being thought of in advance; and it is even foolish to rehearse a future evil, or one that perhaps is not even going to come about. Every evil, when it has arrived, is hateful enough; but the man who has been forever supposing that some misfortune could happen makes that evil everlasting for himself; and if it is not even going to come, then he has taken on a self-imposed wretchedness for nothing. So in either case he is tormented — by suffering the evil, or by brooding on it.

\ciceroSection{33} The relief of distress, however, he places in two things: drawing the mind away from dwelling on what troubles it, and calling it back to the contemplation of pleasures. For he holds that the mind can obey reason and follow where reason leads. Reason, then, forbids it to gaze at troubles, drags it off from bitter reflections, blunts its keenness for contemplating miseries; and when it has sounded the retreat from these, it drives the mind on again and spurs it to behold and to handle with the whole mind the various pleasures with which he holds the wise man's life to be filled, both by the memory of past ones and by the hope of those to come. We have said these things in our own manner, the Epicureans say them in theirs; but let us look at what they say, and let the manner go.

\ciceroSection{34} To begin with, they are wrong to find fault with forethought about things to come. For there is nothing that so blunts and lightens distress as the continual reflection, throughout one's whole life, that nothing can happen which is impossible — as the pondering of the human condition, the law of life and the practice of submitting to it, which brings about not that we should always grieve, but that we should never grieve. For a man who reflects on the nature of things, on the variety of life, on the weakness of the human race, does not grieve when he reflects on these things, but then most of all is discharging the office of wisdom; for he gains both ends at once — in considering human affairs he enjoys the proper task of philosophy, and against adverse chances he is healed by a threefold consolation: first, that nothing has befallen him except what he had long reflected could befall, a reflection which more than anything else thins out and washes away all troubles; next, that he understands human things must be borne in a human way; and last, that he sees there is no evil but fault, and there is no fault when something has come about that no man could have prevented.

\ciceroSection{35} As for that calling-back which he offers, when he draws us away from looking at our evils, it is worth nothing. For it is not in our power, when these things we suppose to be evils are jabbing at us, to feign or to forget. They tear, they harry, they apply their goads, they bring up fire, they will not let us draw breath — and you order us to forget, which is against nature, while you wrench away the help that nature itself has given against inveterate pain? Slow indeed is that medicine, yet great, which length of time and the passing days bring. You bid me think of good things, forget the evils. You would be saying something — something worthy of a great philosopher, even — if you held those things to be good which are most worthy of a human being. Were Pythagoras to say it to me, or Socrates, or Plato:

\ciceroSection{36} why do you lie low, why do you grieve, why do you give way and yield to fortune? Fortune may perhaps have been able to pluck at you and to prick you, but it surely had no power to break your strength. There is great force in the virtues; rouse them, if by chance they are sleeping. At once fortitude, foremost among them, will stand by you, and will compel you to be of so great a spirit that you despise all that can befall a human being and count it as nothing. Temperance will stand by you — the same as moderation, which I called frugality a little while ago — and it will not suffer you to do anything basely or worthlessly. And what is more worthless or more base than an effeminate man? Justice itself will not allow you to do such things, though it might seem to have the least place in this case; yet it will say that you are doubly unjust, since you both reach for what is another's — you who, born a mortal, demand the condition of immortals — and take it hard that you have given back what you received only on loan.

\ciceroSection{37} And prudence — what will you answer to her, when she teaches that virtue is sufficient unto itself, both for living well and likewise for living happily? For if virtue hung bound to externals and did not both spring from itself and turn back again to itself, embracing all that is its own and seeking nothing from elsewhere, I do not see why it should seem either to deserve such ardent praise in words or to be so greatly sought in fact. — If it is to these good things that you call me back, Epicurus, I obey, I follow, I take you yourself for my guide; I even forget my evils, as you bid, and the more easily because I hold that they are not to be reckoned among evils at all. But you divert my thoughts to pleasures. Which ones? Of the body, I suppose, or such as are entertained, by memory or by hope, on account of the body. Is there anything else? Do I read your view rightly? For your followers are in the habit of denying that we understand what Epicurus means.

\ciceroSection{38} This is what he means — and this is what that sharp little old man Zeno, the keenest of them, used to argue and to declaim in a loud voice while I was listening at Athens: that the man is happy who enjoys present pleasures and is confident that he will go on enjoying them, either throughout his whole life or through the greater part of it, with no pain intervening; or, if pain should intervene, then, if it is at its height, it will be brief, and if more drawn out, it will hold more of the sweet than of the bad. A man who reflects on these things will be happy, especially since he is content with goods perceived before and dreads neither death nor the gods. There you have the form of the happy life according to Epicurus, set out in Zeno's words, so that nothing can be denied. What then?

\ciceroSection{39} Will the setting forth and contemplation of this kind of life be able to relieve a Thyestes, or an Aeëtes, of whom I spoke a little before, or a Telamon driven from his country, exiled and destitute? Of him this astonishment was uttered: \textit{Is this that Telamon whom glory but now raised to the sky, on whom men gazed, before whose face the Greeks turned their own?}

\ciceroSection{40} But if, as the same poet says, a man's spirit has collapsed together with his fortunes, then the medicine must be sought from those weighty philosophers of old, not from these devotees of pleasure. For what abundance of goods do these men speak of? Grant, by all means, that the highest good is to feel no pain — though that is not called pleasure, but there is no need to take in everything now — is that the thing to which we are to be diverted in order to lighten our grief? Let it be granted, by all means, that the highest evil is to feel pain: does it follow that the man who is not in that evil, if he is free of evil, thereby straightway enjoys the highest good?

\ciceroSection{41} Why do we shift and dodge, Epicurus, instead of admitting that we mean by pleasure the very thing you yourself, once you have brazened it out, are accustomed to call pleasure? Are these your own words or not? In that book, indeed, which contains the whole of your teaching — for I shall now play the part of a translator, lest anyone suppose I am inventing — you say this: For my part I have no notion of what that good is, if I strip away the pleasures perceived by taste, strip away those perceived through the things of love, strip away those that come to the ears from songs, strip away too those sweet motions perceived by the eyes from beautiful shapes, or whatever other pleasures are engendered in the whole man through any of the senses. Nor indeed can it be said that the mind's gladness alone is among the goods. For the mind's gladness I know to be this: the hope that, by gaining all those things I named above, a man's nature shall be free of pain.

\ciceroSection{42} And this, indeed, in these very words — so that anyone may understand what pleasure Epicurus knew. Then, a little further down: I have often asked, he says, of those who were called wise men, what they would have left to count among the goods, if they took those things away, unless they meant to pour out empty noises; from them I could learn nothing. If they will spout virtues and wisdoms, they will be saying nothing but the road by which those pleasures are produced which I named above. What follows is in the same vein, and the whole book, which is on the highest good, is crammed with words and sentiments of this sort.

\ciceroSection{43} Is it to this kind of life, then, that you will recall the famous Telamon, to lighten his distress? And if you see one of your own broken with grief, will you hand him a sturgeon rather than some little treatise of Socrates? Will you urge him to listen to a water-organ rather than to Plato? Will you spread out before him a show of flowers and colors to gaze upon? Will you hold a little nosegay to his nostrils? Will you burn perfumes and bid him be wreathed with garlands and with roses? And if you add but one thing more—why then you will have wiped away every trace of mourning.

\ciceroSection{44} These things Epicurus must confess, or else strike out of his book the very passages I have just quoted word for word—or rather throw the whole book away, for it is crammed full of pleasures. We must ask, then, how we are to free of his distress a man who speaks like this: \textit{By Pollux, it is fortune that fails me now, more than my birth. For a kingdom once was mine—that you may know from how high a place, with what great resources, and from what an estate fortune has slipped and fallen.} What of it? Is a cup of mulled wine to be thrust upon this man, that he may stop his weeping—or something of that kind? Look, here from the other side, by the same poet: \textit{From the heights of fortune, in want of your help, Hector—} this man we ought to assist; for he is asking for aid: \textit{What protection should I seek, or what pursue, in what aid of exile or of flight should I now put my trust? Of citadel and city I am bereft. Where shall I turn? To whom apply? For me there stand no altars of my fathers' house; they lie shattered and flung apart, the shrines burnt out in the flames, the high walls scorched and standing disfigured, and the fir-wood warped—} you know what follows, and above all this: \textit{O father, O fatherland, O house of Priam, temple shut behind your high-resounding hinge! I saw you, while still the eastern wealth stood by, with your carved and paneled ceilings, royally set out in gold and ivory.}

\ciceroSection{45} What an extraordinary poet!—though he is held in contempt by these chanters of Euphorion. He feels that all things sudden and unlooked-for are heavier to bear; and so, after heaping up the royal riches that seemed bound to last forever, what does he add? \textit{All this I saw set ablaze, the life wrenched out of Priam by violence, the altar of Jupiter defiled with blood.}

\ciceroSection{46} A magnificent poem! For it is mournful in its matter, its words, and its measures alike. Let us take his distress from him. How? Let us lay him on a feather mattress, bring in a girl with a lyre, give him perfume, set a saucer of incense burning, look out for a little sweet drink and something to eat? Are these, then, the goods by which the gravest distresses are to be drawn off? For you yourself said a moment ago that you could not so much as conceive of any goods but these. That a man ought to be called back from grief to the contemplation of goods—on that I should agree with Epicurus, if only we agreed on what a good is. Someone will say: What, then? Do you suppose Epicurus intended any such thing, or that his doctrines were a charter for lust? No indeed, not in the least; for I see much that he says with severity, much with distinction. And so, as I have often said, the question is about his sharpness of mind, not about his morals; however he may scorn the very pleasures he has just praised, I shall still remember what he holds to be the highest good. For he did not lay it down in the bare word "pleasure"; he explained what he meant: taste, he says, and the embrace of bodies, and games and songs, and those shapes by which the eyes are pleasantly stirred. Am I inventing this? Am I lying? I long to be refuted. For what do I labor at, except that the truth in every inquiry should be brought to light?

\ciceroSection{47} "But the same man says that pleasure does not grow once pain is removed, and that the highest pleasure is to feel no pain." Three great errors in a few words. The first, that he is at war with himself. For just now he could not so much as suspect anything to be a good unless the senses were, as it were, tickled with pleasure; but now the highest pleasure is to be free of pain. Could a man be more at war with himself? The second error: that whereas in nature there are three states—one to feel joy, another to feel pain, a third to feel neither joy nor pain—he thinks the first and the third one and the same, and does not distinguish pleasure from the absence of pain. The third error he shares with certain others: that whereas virtue is the thing most to be sought, and philosophy was devised for the very sake of attaining it, he has cut off the highest good from virtue.

\ciceroSection{48} "But he often praises virtue." Gaius Gracchus too, when he had poured out the largest doles and drained the treasury, nonetheless defended the treasury in words. Why should I listen to words when I see deeds? That famous Lucius Piso Frugi had always spoken against the grain law. When the law was carried, consular though he was, he came forward to receive his grain. Gracchus noticed Piso standing in the assembly, and asked him, in the hearing of the Roman people, how he could be consistent with himself, seeking grain under the very law he had argued against. "I should not like it, Gracchus," he said, "to be your pleasure to divide my goods among the people man by man; but if you do it, I will ask my share." Did this grave and wise man not declare plainly enough that the public patrimony was being squandered by the Sempronian law? Read the speeches of Gracchus, and you will call him a guardian of the treasury.

\ciceroSection{49} Epicurus denies that one can live pleasantly unless one lives with virtue; he denies that fortune has any power over the wise man; he prefers a sparing diet to a lavish one; he denies that there is any time at which the wise man is not happy. All worthy of a philosopher—but at war with pleasure. "That is not the pleasure he means." Let him mean whatever he likes; surely he means the pleasure in which no part of virtue is present. Come now—if we do not understand his pleasure, do we not understand pain either? I say, then, that it is not for a man who measures the highest evil by pain to make any mention of virtue.

\ciceroSection{50} And certain Epicureans—excellent men, for no class of people is less given to malice—complain that I argue against Epicurus with too much zeal. As if, I suppose, we were contending over public honor or rank. To me the highest good seems to lie in the soul, to them in the body; to me in virtue, to them in pleasure. And they put up a fight, and even call on their neighbors to back them—and there are many ready to come flocking at once; whereas I am the kind of man who will say that I shall make no trouble of it, but count as settled whatever they have done.

\ciceroSection{51} For what is this? Is it the Punic War we are debating? On that very question Marcus Cato held one view and Lucius Lentulus another, yet there was never any quarrel between them. These men go at it too hot-temperedly—especially since the doctrine they defend is not exactly a high-spirited one, on behalf of which they would not dare to speak in the Senate, nor in the assembly, nor before the army, nor before the censors. But I will deal with these people another time, and in such a spirit that I shall start no contest at all: to those who speak the truth I shall yield readily. I will only give them this warning: granting it were perfectly true that the wise man refers everything to the body—or, to put it more decently, that he does nothing but what is to his advantage, or that he refers everything to his own benefit—since these are not things to win applause, let them rejoice in their own bosom and stop talking grandly. There remains the doctrine of the Cyrenaics;

\ciceroSection{52} they hold that distress arises when something happens unexpectedly. That is indeed a great point, as I said above; I know that Chrysippus too holds this view—that what has not been foreseen beforehand strikes the harder. But not everything lies in this. It is true that a sudden onset of the enemy throws us into rather more confusion than one expected, and a sudden storm at sea alarms sailors more violently than one foreseen, and there are many cases of the kind. But when you consider carefully the nature of the unexpected, you will find nothing else but that all sudden things seem greater—and that for two reasons. First, because there is no room to weigh how great the things are that befall; and second, because, when it appears that the harm could have been guarded against, had it been foreseen, the evil, as if incurred through one's own fault, makes the distress sharper.

\ciceroSection{53} That this is so, time makes plain, which as it passes so softens the wound that, though the same ills remain, the distress is not only eased but in most cases removed. Many Carthaginians were slaves at Rome, and Macedonians after King Perses was taken; I myself, when I was a young man, saw in the Peloponnese certain Corinthians. All of these might have bewailed their lot in those very lines about Andromache: \textit{All this I saw}—but by then, perhaps, they had finished their song. For such was their look, their speech, all the rest of their bearing and posture, that you would have called them Argives or Sicyonians; and the ruined walls of Corinth, suddenly come upon, moved me more than the Corinthians themselves, over whose minds long brooding had drawn the callus of habit.

\ciceroSection{54} We have read the book of Clitomachus, which he sent, after the fall of Carthage, to console the captives, his fellow citizens; in it is written down a discussion of Carneades which he says he had entered in his notebook. The proposition had been set down that the wise man would seem to be in distress when his fatherland was captured; and what Carneades argued against this was written out. So great, then, is the medicine the philosopher applies to a present calamity, while for one grown old it is not even called for; nor, if that same book had been sent to the captives some years later, would it have healed wounds, but only scars. For little by little, step by step, as it advances, grief is worn thin—not because the thing itself is wont, or able, to be changed, but because experience teaches what reason ought to have taught: that those things are smaller which had seemed greater. What need is there, then, someone will say, of reasoning at all, or of that consolation we are accustomed to use when we wish to lighten the grief of mourners?

\ciceroSection{55} For this is roughly what we then have ready to hand: that nothing ought to seem unforeseen. Or how will a man bear his misfortune more tolerably, who has recognized that some such thing must befall a human being? For this kind of talk takes nothing from the sheer sum of the evil; it only brings home that nothing has happened which should not have been expected. And yet that mode of speech is not without force in consolation—indeed, I rather think it has the most. So those unexpected things do not have such power that all distress arises from them; for they strike, perhaps, more heavily, but they do not bring it about that the things which befall appear greater. They appear greater because they are recent, not because they are sudden.

\ciceroSection{56} There is, then, a twofold method of finding the truth, not only in the things that seem evil, but also in those that seem good. For either we inquire of what kind and how great the nature of the thing itself is—as sometimes about poverty, whose burden we lighten by arguing and showing how few and how small are the things that nature requires—or we lead our discourse away from subtlety of argument to examples. Here Socrates is recalled, here Diogenes, here that line of Caecilius: \textit{Often, even under a shabby cloak, there is wisdom.} For since the force of poverty is one and the same, what can possibly be said why it was bearable to Gaius Fabricius, while others deny that they can endure it?

\ciceroSection{57} Similar, then, to that second kind is the method of consolation which teaches that what has befallen us is merely human. For that line of argument does not only convey an understanding of the human lot; it signifies as well that what others have borne and still bear is bearable. The talk is of poverty: many poor men are recalled who endured it patiently. The talk is of holding office in contempt: many are brought forward who went unhonored, and were the happier for that very thing; and the lives of those who set private leisure above public business are praised by name; nor is that anapaest of the most powerful of kings passed over in silence, who praises an old man and calls him fortunate because he is to reach his last day without glory and without renown.

\ciceroSection{58} In the same way, by recalling examples, the loss of children too is celebrated, and the grief of those who bear it too heavily is eased by the examples of others. So the endurance of others makes what has befallen us seem far smaller than it had been reckoned to be. Thus it comes about, little by little, that to those who think it over it becomes plain how far opinion has lied. This is just what that Telamon declares: \textit{I, when I begot them} —— and Theseus: \textit{I rehearsed in advance the miseries that would be mine} —— and Anaxagoras: \textit{I knew that I had begotten a mortal.} For all these men, long pondering human affairs, understood that such things were by no means to be dreaded as the common crowd supposes. And to me, indeed, it seems that very nearly the same thing happens to those who reflect beforehand as to those whom time heals — except that a certain reasoning cures the one, and nature herself the others, once the point that holds the whole matter is grasped: that the evil which was supposed to be the greatest is by no means so great as to be able to overturn the happy life.

\ciceroSection{59} The upshot of this, then, is that the blow is greater for coming unforeseen — not, as those others suppose, that when equal misfortunes befall two men, only the one to whom his misfortune came unforeseen is afflicted with distress. And so some men in their mourning are said to have borne it the more heavily once they had heard of this common condition of mankind, that we are born under such a law that no one can be forever exempt from evil. For which reason Carneades, as I see our Antiochus write, used to take Chrysippus to task for praising that famous song of Euripides: \textit{No mortal lives whom grief and sickness do not touch; many must bury their children, and beget others in their place, and death is appointed for all. These bring anguish to the human race in vain: earth must be rendered back to earth, and then life must be reaped, for all, like the corn. So Necessity commands.}

\ciceroSection{60} He denied that this kind of speech had anything whatever to do with lightening distress. For this very thing, he said, was to be grieved at — that we had fallen into so cruel a necessity; while that other speech, drawn from the recital of others' misfortunes, was suited only to consoling the ill-natured. But to me it seems far otherwise. For the necessity of bearing the human condition forbids us, as it were, to fight against a god, and reminds us that we are men — a thought that greatly lightens grief; and the enumeration of examples is brought forward, not to gratify the minds of the ill-natured, but so that the one who mourns may judge he must bear what he sees many have borne with measure and tranquillity.

\ciceroSection{61} For those who collapse and cannot hold together under the magnitude of their distress must be propped up by every means. From this Chrysippus thinks distress itself was called \textit{[Greek: lupē]}, as it were a dissolution of the whole man. All of this can be torn out once the cause of distress is laid open, as I said at the outset; for there is no other cause but the opinion and judgment of some great evil present and pressing. And so bodily pain, whose sting is sharpest, will be borne through with the hope of a good held out before us, and a life lived honorably and splendidly brings such consolation that those who have so lived are either untouched by distress or pricked only very lightly by pain of mind. But when to this opinion of a great evil there is added the further opinion that one ought, that it is right, that it belongs to one's duty to bear hard what has befallen — then at last that grievous disturbance of distress is produced.

\ciceroSection{62} From this opinion come those varied and detestable kinds of mourning: squalor, women's tearing of cheeks, beatings of breast and thigh and head; hence that Agamemnon of Homer's, the same in Accius too, \textit{rending in his grief his unshorn hair again and again}; on which there is that witty remark of Bion's, that the most foolish of kings did just as well to tear out his hair in mourning, as though grief were lightened by going bald.

\ciceroSection{63} But men do all these things supposing that it ought to be done so. And so Aeschines, too, inveighs against Demosthenes for having sacrificed victims on the seventh day after his daughter's death. But with what rhetoric, with what fullness! What sentiments he gathers, what words he hurls about! So that you may understand a rhetorician is allowed anything. None of this would anyone approve, did we not have implanted in our minds the notion that all good men ought to mourn as heavily as possible at the loss of their own. From this it comes about that, in pains of the mind, some men seek out solitudes, as Homer says of Bellerophon: \textit{Wretched, mourning, he wandered over the Aleian plains, eating his own heart, shunning the tracks of men;} and Niobe is imagined turned to stone on account, I suppose, of her unbroken silence in grief, while Hecuba they think is imagined turned into a dog on account of a certain bitterness and frenzy of mind. But there are others whom it often delights, in their mourning, to speak with solitude itself, like that nurse in Ennius: \textit{The longing has seized me, wretched, to speak out now to heaven and earth the miseries of Medea.}

\ciceroSection{64} All these things men do in their grief, thinking them right, true, and owed; and that this is done as if by a judgment of duty is most clearly shown by the fact that, if any happen — while wishing to be in mourning — to have done something more humanely or to have spoken more cheerfully, they call themselves back again to sadness and accuse themselves of a fault, that they have left off grieving. Boys, indeed, their mothers and tutors are even accustomed to chastise — and not with words only, but with blows — if anything is done or said by them more cheerfully amid the household's mourning, and they force them to weep. What of this? When the slackening of mourning has at last followed, and it has been understood that nothing is gained by mourning, does not the very thing make plain that the whole of it was a matter of the will?

\ciceroSection{65} What of that man in Terence who punishes himself — that is, the \textit{[Greek: heauton timoroumenos]}? \textit{I have resolved, Chremes, to wrong my son this much the less for just so long as I make myself wretched.} Here he resolves to be wretched. But does anyone resolve upon anything against his will? \textit{Indeed I would judge myself deserving of any evil} — he judges himself deserving of evil unless he is wretched. You see, then, that the evil is one of opinion, not of nature. What of those whom the situation itself forbids to mourn? As in Homer the daily slayings and deaths of many bring an abating of grief; he speaks thus: \textit{For we see too many falling, and in every daylight, so that no one can be free from mourning. The more fitting it is, then, to commit the slain to their graves with a firm spirit, and to end mourning with a single day's tears.}

\ciceroSection{66} Therefore it is within our power to cast off pain when we will, in deference to the occasion. Or is there any occasion — since the matter is in our own power — to which we may not defer for the sake of laying aside care and distress? It was well established that those who had seen Gnaeus Pompeius falling under his wounds, even at that most bitter and wretched spectacle, fearing for themselves because they saw themselves hemmed in by the enemy's fleet, did nothing else at the time but urge on the rowers and seek their safety in flight; and only after they had come to Tyre did they begin to be cast down and to lament. Fear, then, could drive distress away from these men; will reason not be able to do so for a wise man? And what is there that avails more for laying aside pain than the realization that nothing is gained by it and that it was taken up in vain? If, then, it can be laid down, it can also not be taken up; we must confess, therefore, that distress is taken up by will and by judgment.

\ciceroSection{67} And this is shown by the endurance of those who, having often suffered many things, more easily bear whatever befalls, and judge that they have by now grown hardened against fortune, like that man in Euripides: \textit{If now for the first time this grievous day had dawned on me, and I had not voyaged on so wretched a sea, there would be cause for grieving — just as colts, when the bridle is first thrown on, are thrown into a panic by the strange new touch; but, broken now by my miseries, I have gone numb.} Since, then, the weariness of misery makes distress milder, it must be understood that the thing itself is not the cause and fountain of mourning.

\ciceroSection{68} Do not the highest philosophers, who have not yet attained wisdom, understand that they are in the greatest of evils? For they are unwise, and there is no greater evil than unwisdom. And yet they do not mourn. Why so? Because to this kind of evil there is not attached that opinion that it is right and fitting and belongs to one's duty to bear hard the fact that one is not wise — the same opinion we do attach to this distress, in which mourning resides, which is the greatest of all distresses.

\ciceroSection{69} And so Aristotle, accusing the old philosophers who had supposed philosophy brought to perfection by their own genius, says that they were either most foolish or most boastful; but that he himself saw, since in a few years a great addition had been made, that in a short time philosophy would be plainly completed. Theophrastus, again, is said to have accused nature as he was dying, because she had given long life to stags and crows, to whom it made no difference at all, and to men, to whom it made the greatest difference, so brief a span; whereas if their age could have been longer, it would have come to pass that human life would be schooled in every accomplished art and every branch of learning. He complained, then, that he was being snuffed out just when he had begun to see those things. What of this? Among the rest of the philosophers, does not each best and weightiest of them confess that he is ignorant of many things, and that there are many things he must learn over and over again?

\ciceroSection{70} And yet, though they understand that they are stuck in the middle of folly, than which nothing is worse, they are not weighed down by distress; for no opinion of a dutiful pain is mixed in. What of those who think that men ought not to mourn? Such was Quintus Maximus, who carried out for burial his son of consular rank; such was Lucius Paulus, who lost two sons within a few days; such was Marcus Cato, whose son died after his designation as praetor; such were the rest, whom we have gathered together in the \textit{Consolation}.

\ciceroSection{71} What else calmed these men, if not that they did not think it the part of a man to mourn and grieve? Therefore that very thing to which others, thinking it right, are wont to surrender themselves in distress — that these men, thinking it shameful, drove distress away. From which it is understood that distress lies not in nature but in opinion. Against this the following is urged: who is so deranged as to grieve of his own free will? Nature, they say, brings on the pain — the pain to which your own Crantor thinks one must yield; for it presses and bears down, and cannot be withstood. And so that Oileus in Sophocles, who had earlier consoled Telamon over the death of Ajax, when he heard the news about his own son, was broken. About his altered mind it is put thus: \textit{Nor is anyone endowed with wisdom so great that, while easing the affliction of others with his words, he is not himself broken by his own sudden disaster when fortune, changed, turns its assault against him — so that those very sayings and counsels addressed to others fall away.} When they argue in this way, what they are bent on establishing is that nature can in no way be resisted; yet these same men admit that heavier distresses are taken on than nature compels. What madness is it, then —? that we too may put the same question to them.

\ciceroSection{72} But there are several causes for taking on pain. First, that opinion of evil, upon the seeing and being convinced of which distress necessarily follows. Next, men suppose that they do something pleasing to the dead if they mourn them heavily. To this is added a kind of womanish superstition; for they imagine they will more easily satisfy the immortal gods if they confess themselves stricken by their blow, prostrate and laid low. But how these things conflict with one another most men do not see. For they praise those who die with an even mind, while they think those who bear another's death with an even mind deserving of blame — as if it could possibly come about, what is wont to be said in a lover's talk, that anyone should love another more than himself.

\ciceroSection{73} That is a fine thing, and, if you ask me, a right and true one too: that those who ought to be dearest to us we should love as much as our very selves; but that we should love them more is in no way possible. Nor is it even to be wished for in friendship, that he should love me more than himself, or I him more than myself; a confounding of life and of all its duties would follow, if it were so. But of this elsewhere; for now it is enough not to lay our wretchedness to the loss of friends, lest we love them more than they themselves, if they could perceive it, would wish, and certainly more than our very selves. As for their saying that most people are not relieved at all by consolations, and adding that the consolers themselves confess they are wretched when fortune turns its assault upon them, both points are dissolved. For these are flaws not of nature but of fault. And folly may be accused as fully as one likes. For both those who are not relieved invite themselves to wretchedness, and those who bear their own misfortunes otherwise than they prompted others to do are no more faulty than nearly all those who, being greedy, find fault with the greedy, or, being vainglorious, with the lovers of glory. For it is the very mark of folly to discern the faults of others and forget one's own.

\ciceroSection{74} But surely this above all is to be brought out: since it is agreed that distress is removed by length of time, this power lies not in the day but in long reflection. For if both the matter is the same and the man is the same, how can anything about the pain be altered, if nothing is altered either about that on account of which he grieves or about the one who grieves? It is long reflection, then, that there is no evil in the matter, that heals the pain, not mere length of time. Here they bring me their moderate degrees of feeling. If these are natural, what need is there of consolation? For nature herself will set the limit; but if they rest on opinion, let the whole opinion be removed. I think enough has been said that distress is an opinion of present evil, in which opinion this is contained: that it is fitting to take on distress.

\ciceroSection{75} To this definition Zeno rightly adds that the opinion of present evil be fresh. But this word they interpret in such a way that they would have it called fresh not only when the thing has happened a little while before, but for as long as there is a certain force in that imagined evil, so that it is vigorous and keeps a kind of greenness — for so long it is called fresh. So that Artemisia, the wife of Mausolus king of Caria, who built that famous tomb at Halicarnassus, lived in mourning as long as she lived, and, worn out by it, wasted away. To her that opinion was fresh every day; and it is not called fresh only then, when it has dried up with the lapse of time. These, then, are the duties of those who console: to remove distress utterly, or to calm it, or to draw off as much of it as possible, or to suppress it and not let it spread further, or to divert it to other things.

\ciceroSection{76} There are those who think the one duty of a consoler is to show that the thing is no evil at all, as Cleanthes holds; there are those who hold it is no great evil, as the Peripatetics; there are those who would lead a man away from evils to goods, as Epicurus; there are those who think it enough to show that nothing unforeseen has happened, as the Cyrenaics think nothing an evil. Chrysippus, however, judges that the chief thing in consoling is to remove from the mourner that opinion by which he supposes himself to be discharging a just and owed duty. There are also those who gather all these kinds of consoling together — for one man is moved one way, another another — as I myself, in my \textit{Consolation}, threw nearly everything together into a single consolation; for the mind was in a swollen state, and every cure was being tried upon it. But the right time must be taken in diseases of the mind no less than in those of the body; as that Prometheus of Aeschylus, to whom, when it had been said: \textit{And yet, Prometheus, I judge that you hold this — that reason can heal anger,} replied: \textit{Yes, provided one applies the remedy in good season and does not press the festering wound with a heavy hand.}

\ciceroSection{77} In consolations, then, the first medicine will be to teach either that there is no evil, or a very small one; the second, to speak both of the common condition of life and, particularly, of anything that needs to be argued about the case of the mourner himself; the third, that it is the height of folly to be consumed by mourning to no purpose, when you understand that nothing can be gained by it. For Cleanthes, indeed, consoles the wise man, who has no need of consolation. For if you have persuaded a mourner that nothing is an evil except what is shameful, you will have taken from him not his mourning but his folly; and that is no fit time for teaching. And yet Cleanthes does not seem to me to have seen this well enough: that distress can sometimes be taken on out of that very thing which Cleanthes himself admits to be the highest evil. For what shall we say when Socrates had persuaded Alcibiades, as we are told, that he was nothing of a man, and that there was no difference between Alcibiades, born in the highest station, and any common porter — and when Alcibiades, distraught and weeping, begged Socrates as a suppliant to impart virtue to him and drive out his baseness — what shall we say, Cleanthes? That in that matter, which afflicted Alcibiades with distress, there was no evil?

\ciceroSection{78} What of Lyco's words — of what sort are they? He, making light of distress, says it is stirred by small things, by the discomforts of fortune and of the body, not by the evils of the mind. What then? That over which Alcibiades grieved — did it not consist of the evils and faults of the mind? About the consolation of Epicurus enough has been said before.

\ciceroSection{79} Not even that consolation is the firmest, though it is both in common use and often of profit: "this has not befallen you alone." It is of profit, as I said, but not always nor for everyone; for there are those who reject it. But it matters how it is applied. For we must set forth how each of those who bore things wisely bore them, not by what misfortune each was afflicted. Chrysippus's consolation is the firmest with respect to truth, but difficult for the moment of distress. It is a great task to prove to a mourner that he mourns by his own judgment and because he thinks it fitting so to act. Surely, then, just as in lawsuits we do not always use the same issue — for so we call the kinds of disputed cases — but adapt ourselves to the occasion, to the nature of the controversy, to the person, so in the easing of distress one must see what cure each man is able to receive.

\ciceroSection{80} But somehow or other the discussion has strayed from what was set down by you. For you had asked about the wise man, to whom either nothing can seem an evil that is free of baseness, or so small an evil that it is overwhelmed by wisdom and scarcely shows — a man who adds nothing to his distress by opinion, nor takes anything on, and does not think it right that he should be racked to the utmost and consumed by mourning, than which nothing could be more perverse. Yet reason has taught us, as it seems to me at least — although this very point was not properly under inquiry at this time, namely whether anything is an evil that could not also be called shameful — it has nonetheless taught us to see this: that whatever of evil there is in distress is not natural, but contracted by a voluntary judgment and by an error of opinion.

\ciceroSection{81} And I have treated that kind of distress which alone is the greatest of all, so that, with it removed, we should judge the remedies for the rest not greatly worth seeking. For there are settled things wont to be said about poverty, settled things about a life unhonored and inglorious; there are separate, settled schools of argument about exile, about the ruin of one's country, about slavery, about infirmity, about blindness, about every chance to which the name of calamity is wont to be given. These the Greeks parcel out into separate schools and separate books; for they seek out the work — though full discussions are a delight;

\ciceroSection{82} and yet, just as physicians, in tending the whole body, heal even the smallest part if it has come to ache, so philosophy, when it has removed distress entire, removes it also if any error has arisen from some quarter, if poverty has bitten, if disgrace has stung, if exile has spread some darkness over a man, or if any of the things I have just named has arisen. Although there are consolations proper to each single matter, about which you shall hear, indeed, whenever you wish. But we must return to the same source: that all distress is far from the wise man, because it is empty, because it is taken on to no purpose, because it springs not from nature but from judgment, but from opinion, but from a kind of invitation to grieve, when we have decreed that so it ought to be done.

\ciceroSection{83} When this has been stripped away — which is wholly voluntary — that mourning distress will be removed, yet a certain bite and small contraction of the mind will be left. Let them by all means call this natural, provided the name of distress is absent — that grave, foul, deadly name, which can in no way coexist with wisdom and, so to speak, dwell with it. But what roots distress has, how many, how bitter! All of which, once the trunk itself is overthrown, must be plucked out, and, if it shall be necessary, by separate discussions. For this leisure remains to us, of whatever sort it is. But there is one principle behind all distresses, though many names. For to envy belongs to distress, and to be jealous, and to disparage, and to pity, and to be in anguish, to mourn, to grieve, to be afflicted with trouble, to lament, to be harassed, to feel pain, to be in vexation, to be cast down, to despair.

\ciceroSection{84} All these the Stoics define, and the words I have named belong to single things and do not, as they seem to, signify the same things, but differ in something — which perhaps we shall treat in another place. These are those fibers of the roots which I spoke of at the outset, to be tracked down and all plucked out, so that none may ever be able to spring up. A great task and a difficult one — who denies it? But what that is glorious is not also arduous? Yet philosophy professes that it will accomplish this, if only we accept its cure. But of these things, indeed, thus far; the rest, as often as you wish, will be ready for you both in this place and in others.

\ciceroBook{4}

\ciceroSection{1} On many counts, Brutus, I have long admired the genius and the virtues of our countrymen, but most of all in those pursuits which they sought out only quite late and carried over into this state from Greece. For while from the very first rise of the city the auspices, the rites, the assemblies, the right of appeal, the council of the fathers, the ordering of horse and foot, the whole art of war had been established by something like divine guidance — partly through the institutions of the kings, partly through laws as well — there followed, once the commonwealth was freed from the dominion of the kings, an advance toward every excellence so admirable, a course so swift, as to be scarcely believable. This, however, is not the place to speak of the manners and institutions of our ancestors and of the discipline and balance of the state; I have said enough on these matters, with sufficient care, in other places, and above all in the six books I wrote on the commonwealth.

\ciceroSection{2} But at this point, as I turn over in my mind the pursuit of learning, much occurs to me to suggest that these studies too were fetched from abroad — not merely sought after, but kept alive and cultivated. For there stood almost within their sight, a man of surpassing wisdom and renown, Pythagoras, who was in Italy in the very times when Lucius Brutus freed his country — that illustrious founder of your own nobility. And since the teaching of Pythagoras flowed far and wide, it seems to me to have soaked into this state of ours; and this is not only probable by conjecture but is shown by certain traces as well. For who could suppose that, when the part of Italy then called Magna Graecia flourished with cities of the greatest size and power, and in them the name first of Pythagoras himself and afterward of the Pythagoreans stood so high, the ears of our countrymen were closed to those most learned voices?

\ciceroSection{3} Indeed, I believe it was out of admiration for the Pythagoreans that King Numa too was reckoned by later generations a Pythagorean. For knowing the discipline and institutions of Pythagoras, and having received from their ancestors the report of that king's justice and wisdom, yet ignorant of the dates and periods through their great antiquity, they believed that a man of such surpassing wisdom must have been a pupil of Pythagoras. So much, then, for conjecture. As for the traces of the Pythagoreans, though many could be gathered, I shall use only a few, since that is not my business at present. For since they are said to have been in the habit of conveying certain precepts in verse, and somewhat secretly, and of leading their minds, by song and the lyre, away from the strain of thought toward tranquillity, Cato — a most weighty authority — said in his \textit{Origines} that among our ancestors there was this custom at banquets: that the guests, each in turn as they reclined, should sing to the pipe the praises and virtues of famous men. From which it is plain that there were then both songs set to the sounds of voices and poems.

\ciceroSection{4} Though indeed even the Twelve Tables declare that already in those days the composing of verse was customary, since they laid it down by law that it should not be composed to another man's injury. Nor is it a small proof of cultivated times that the lyre is played before the gods upon their sacred couches and before the banquets of the magistrates — which was the very mark of the discipline of which I am speaking. To me, indeed, even the song of Appius Caecus, which Panaetius warmly praises in a certain letter addressed to Quintus Tubero, seems Pythagorean. Many things, too, in our own institutions are drawn from theirs; but I pass them over, lest we seem to have learned elsewhere what we are thought to have discovered ourselves.

\ciceroSection{5} But to bring the discussion back to its purpose: in how short a time, how many poets, and how great, and what orators arose! So that it readily appears our countrymen could achieve everything the moment they had begun to wish it. But of the other pursuits I shall speak in another place, should there be occasion — and have often spoken already. The pursuit of wisdom is indeed an old thing among us; yet before the age of Laelius and Scipio I cannot find men I could name by name. In their youth I see that the Stoic Diogenes and the Academic Carneades were sent by the Athenians as envoys to the Senate; and though these men had never touched any part of public life, and the one was a Cyrenaean, the other a Babylonian, surely they would never have been roused from their schools and chosen for that office, had there not in those days been, in certain leading men, a real zeal for learning. And while these men committed other things to writing — some the civil law, some their own speeches, some the records of their ancestors — that most ample of all the arts, the discipline of living well, they pursued by their lives rather than by their letters.

\ciceroSection{6} And so of that true and elegant philosophy which, descending from Socrates, has persisted to this day among the Peripatetics, and among the Stoics who say the same things in another way, while the Academics arbitrated their disputes — of this there are scarcely any, or very few, Latin records, whether on account of the magnitude of the matters and the busy lives of men, or because they did not think such things could be made acceptable to the unlearned. Meanwhile, with these men silent, Gaius Amafinius came forward to speak; and when his books were published, the crowd was stirred and gave itself above all to that doctrine — whether because it was very easy to grasp, or because they were enticed by the seductive lures of pleasure, or because, since nothing better had been put before them, they held to what there was.

\ciceroSection{7} After Amafinius, many rivals of the same way of thinking wrote much and took possession of all Italy; and the very thing that is the greatest proof those doctrines are not expounded with subtlety — that they are so easily learned by heart and approved by the unlearned — this they take to be the firm foundation of their teaching. But let each man defend what he thinks; for judgments are free. I shall hold to my own practice, and, bound by no laws of any single school which I must necessarily obey in philosophy, I shall always seek out what is most probable in each matter. This I have often pursued on other occasions, and lately, with much zeal, at my place at Tusculum. And so, the discussions of three days having been set out, this fourth day is brought to a close in the present book. For when we had gone down into the lower walk, as we had done on the previous days too, the matter was carried on thus: Let him speak, if anyone wishes, on the topic he would have discussed.

\ciceroSection{8} — It does not seem to me that the wise man can be free of every disturbance of the soul. — From distress, at least, he seemed to be free by yesterday's discussion — unless perhaps you were agreeing with me for the moment's sake. — Not at all; for your argument was thoroughly approved by me. — You do not suppose, then, that distress can befall the wise man? — I do not think so at all. — And yet, if that cannot disturb the soul of the wise man, nothing will be able to. For consider: would fear throw it into confusion? But fear is of those things, when absent, of which distress is felt when present; remove distress, then, and fear is removed. There remain two disturbances: exuberant gladness and desire. And if these do not befall the wise man, the mind of the wise man will always be tranquil.

\ciceroSection{9} — That is exactly how I understand it. — Which, then, do you prefer? That we set sail at once, or that, like men just leaving harbor, we row a little first? — What do you mean by that? For I do not understand. — Because when Chrysippus and the Stoics discuss the disturbances of the soul, they are largely taken up with dividing and defining them, while that part of their discourse is exceedingly meager by which they might heal the soul and keep it from being turbulent; whereas the Peripatetics bring forward much for the calming of the soul and pass over the thorns of dividing and defining. So I was asking whether I should spread the sails of my discourse at once, or first drive it forward a little with the oars of the logicians. — That way, by all means; for the whole of what I seek will be more complete from both together. — That is indeed the better course;

\ciceroSection{10} but you shall ask afterward, if anything proves rather obscure. — I shall do so; and yet you, as is your habit, will state these very obscure points more plainly than the Greeks state them. — I shall make the effort; but there is need of a concentrated mind, lest everything slip away if any one thing escapes us. Since what the Greeks call \textit{[Greek: pathē]} we prefer to call disturbances rather than diseases, in explaining them I shall follow that old classification, first of Pythagoras and then of Plato, who divide the soul into two parts: one they make to share in reason, the other to have no part in it. In the part that shares in reason they place tranquillity, that is, a placid and quiet steadiness; in that other part they place turbid motions, both of anger and of craving, contrary and hostile to reason.

\ciceroSection{11} Let this, then, be our source; yet in describing these disturbances let us use the definitions and divisions of the Stoics, who seem to me to handle this question most acutely. This, then, is Zeno's definition: that a disturbance — what he calls \textit{[Greek: pathos]} — is a commotion of the soul averted from right reason and contrary to nature. Some define it more briefly as a more violent impulse — and they hold that "more violent" means one that has departed further from the steadiness of nature. The kinds of disturbances they hold to arise from two supposed goods and from two supposed evils; thus there are four — from goods, desire and gladness, gladness being of present goods, desire of future ones; from evils, they reckon, arise fear and distress, fear directed at the future, distress at the present. For the things that are feared in their coming afflict us with distress in their arrival.

\ciceroSection{12} Now gladness and desire turn upon the supposition of goods, since desire, enticed and inflamed, is carried off toward what seems good, while gladness is borne aloft and exults when it has now gained something craved. For by nature all men pursue those things which seem good and flee their opposites; and so, the instant the appearance of anything that seems good is set before us, nature herself impels us to obtain it. When this happens steadily and prudently, the Stoics call that kind of striving \textit{[Greek: boulēsin]}; let us call it rational wish. They hold that it exists in the wise man alone, and define it thus: rational wish is that which desires a thing in accordance with reason. But that which, against the opposition of reason, is roused more violently, is desire, or unbridled craving, which is found in all the foolish.

\ciceroSection{13} Likewise, when we are so moved as to be in some good, this happens in two ways. For when the soul is moved by reason placidly and steadily, that is called joy; but when the soul exults emptily and without measure, that may be called exuberant or excessive gladness, which they define thus: an irrational elation of the soul. And since, just as by nature we strive after goods, so by nature we turn aside from evils, this turning aside, if it is done with reason, may be called caution, and is understood to exist in the wise man alone; but that which is without reason and accompanied by a low and broken loss of nerve, let it be named fear. Fear, then, is caution averted from reason.

\ciceroSection{14} As for a present evil, the wise man feels no such affection toward it; but the foolish feel distress, and are afflicted at supposed evils, casting down and contracting their souls in disobedience to reason. And so this is the first definition: that distress is a contracting of the soul against the opposition of reason. Thus there are four disturbances and three good states, since to distress no good state is opposed. But they hold that all disturbances arise from judgment and from supposition. And so they define them more precisely, that it may be understood not only how faulty they are but also how far they lie within our power. Distress, then, is a fresh supposition of a present evil, in which it seems right for the soul to be cast down and contracted; gladness, a fresh supposition of a present good, in which it seems right to be borne aloft; fear, a supposition of an impending evil that seems unbearable; desire, a supposition of a coming good, that it would be of use for it to be already present and at hand.

\ciceroSection{15} But the judgments and beliefs which I have said the disturbances to be — these, they say, are not the only seat of the disturbances; there are also the effects that the disturbances produce, as distress produces a kind of bite of pain, fear a sort of shrinking and flight of the mind, gladness an unrestrained merriment, desire an unbridled appetite. As for opinion, which we included in all the definitions above, they hold it to be a feeble assent.

\ciceroSection{16} But under each single disturbance several species of the same genus are ranged: under distress, envy — for in teaching I must use a less familiar word, since the ordinary "envy" is applied not only to the man who envies but also to the man who is envied — and rivalry, disparagement, pity, anguish, mourning, grief, tribulation, sorrow, lamentation, worry, vexation, affliction, despair, and any others of the same kind. Under fear are ranged sluggishness, shame, terror, dread, panic, faintness, confusion, alarm; under pleasure, malice that rejoices in another's misfortune, delight, exultation, and the like; under desire, anger, irascibility, hatred, enmity, discord, want, longing, and the rest of that sort. These they define in this way: envy, they say, is distress incurred on account of another's prosperity, where that prosperity does no harm to the one who envies.

\ciceroSection{17} For if a man grieves at the prosperity of one by whom he is himself injured, he would not rightly be said to envy, as Agamemnon could not be said to envy Hector; but the man whom another's advantages do no harm, and who nonetheless grieves that the other enjoys them — he surely envies. Rivalry, for its part, is spoken of in two senses, so that the word stands both for a virtue and for a vice; for the imitation of virtue is also called rivalry — though we make no use of that sense here, since it belongs to praise — and there is also the rivalry that is a distress, when another has gained possession of what one craved oneself but does without. Disparagement, again — by which I mean what is understood as jealousy \textit{[Greek: zēlotypia]} — is distress arising from the fact that another too possesses the very thing one craved oneself.

\ciceroSection{18} Pity is distress at the wretchedness of another suffering under an undeserved injury — for no one is moved to pity by the punishment of a parricide or a traitor; anguish is distress that presses hard; mourning is distress at the bitter death of one who was dear; grief is tearful distress; tribulation is toilsome distress; sorrow is racking distress; lamentation is distress with wailing; worry is distress with brooding; vexation is lasting distress; affliction is distress with torment of the body; despair is distress without any expectation that things will turn out better. As for what is ranged under fear, they define it thus: sluggishness is the fear of toil to come; ...

\ciceroSection{19} ... terror is fear that strikes a blow, from which it comes that blushing attends shame, and pallor and trembling and the chattering of teeth attend terror; dread is fear of an evil drawing near; panic is fear that drives the mind from its place — whence that line of Ennius: \textit{then panic drives all wisdom out of my breast, and I am unmanned}; faintness is fear that comes close upon panic and is, as it were, its companion; confusion is fear that shakes loose one's thoughts; alarm is fear that abides.

\ciceroSection{20} The species of pleasure they describe in this way: malice is pleasure at another's misfortune with no profit to oneself; delight is pleasure that soothes the mind by the sweetness of hearing; and as this pleasure of the ears is, so are those of the eyes, of touch, of smell, and of taste, all of one genus, pleasures that flow over the mind as if it were melted by them. Exultation is exuberant pleasure that carries itself off with overweening insolence.

\ciceroSection{21} What is ranged under desire is defined thus: anger is the desire to punish one who is thought to have done an undeserved injury; irascibility is anger newly born and just coming to be, which is called \textit{[Greek: thymōsis]} in Greek; hatred is anger grown inveterate; enmity is anger watching for the moment of revenge; discord is anger more bitter, conceived deep within the soul and the heart; want is insatiable desire; longing is the desire to see one who is not yet present. They draw this further distinction: that desire is for the things which are predicated of some person or persons — what the logicians call predicates \textit{[Greek: katēgorēmata]}, such as to have riches, to win offices — while want is for the things themselves, such as offices, such as money.

\ciceroSection{22} They say that the source of all the disturbances is intemperance, which is a falling away of the whole mind from right reason, so turned aside from reason's prescription that the appetites of the soul can in no way be governed or held in check. Just as temperance, then, calms the appetites and brings them to obey right reason and preserves the considered judgments of the mind, so its enemy intemperance inflames, throws into turmoil, and goads the soul's whole condition; and thus distresses and fears and all the remaining disturbances are begotten of it. Just as, when the blood is corrupted, or phlegm or bile is in excess, diseases and sicknesses arise in the body, so the turmoil of perverse beliefs, and their warring with one another, strips the soul of soundness and racks it with diseases.

\ciceroSection{23} From the disturbances, first, are produced the diseases, which they call \textit{[Greek: nosēmata]}, together with their contraries, which carry a faulty aversion and disgust toward particular things; then the sicknesses, which are called by the Stoics \textit{[Greek: arrōstēmata]}, and likewise the aversions set over against them as their opposites. At this point too much labor is spent by the Stoics, and most of all by Chrysippus, in comparing the diseases of the soul to a likeness of the diseases of the body. Let us pass over that discourse, as quite unnecessary, and handle thoroughly the points that hold the matter together.

\ciceroSection{24} Let it be understood, then, that a disturbance is always in motion, tossed inconstantly and turbulently by the beliefs that toss themselves about; but when this heat and agitation of the soul has grown inveterate and has settled, as it were, in the veins and marrow, then there arise both disease and sickness and those aversions which are contrary to those diseases and sicknesses. The things I am speaking of differ in thought but in fact are bound together, and they take their rise from desire and from gladness. For when money is craved and reason is not at once applied, like a kind of Socratic medicine to heal that craving, the evil seeps through into the veins and clings within the inmost parts, and there arise disease and sickness, which once inveterate cannot be torn out; and the name of that disease is avarice.

\ciceroSection{25} Likewise the other diseases — such as the craving for glory, such as womanizing, to give that name to what is called in Greek \textit{[Greek: philogynia]} — and the rest of the diseases and sicknesses are born in the same way. Their contraries are thought to be born of fear: such as hatred of women, of the kind found in the misogynist \textit{[Greek: misogynōi]} of Atilius; such as hatred toward the whole human race, which we have heard reported of Timon, who is called the misanthrope \textit{[Greek: misanthrōpos]}; such as inhospitality: all which sicknesses of the soul are born of a certain fear of the very things that men flee and hate.

\ciceroSection{26} They define sickness of the soul as a vehement opinion, clinging and deeply implanted, about a thing not worth desiring, as though it were greatly worth desiring. What is born of aversion they define thus: a vehement opinion, clinging and deeply implanted, about a thing not to be fled, as though it were to be fled; and this opinion is the judgment that one knows what one does not know. To sickness certain such things are ranged: avarice, ambition, womanizing, obstinacy, gluttony, drunkenness, daintiness, and any like them. Avarice is a vehement opinion about money, clinging and deeply implanted, as though it were greatly worth desiring; and the definition of the rest of the same genus is similar.

\ciceroSection{27} The definitions of the aversions are of this kind: inhospitality is a vehement opinion, clinging and deeply implanted, that a guest is greatly to be fled; and likewise hatred of women is defined, as in Hippolytus, and hatred of the human race, as in Timon. And to come to the likeness of bodily health and to make use of that comparison at last — though more sparingly than the Stoics are wont to do: just as some men are more prone to some diseases than others — and so we call certain men subject to colds, certain to colic, not because they actually have them now, but because they often do — so some are prone to fear, others to some other disturbance; whence in some there is said to be anxiety, from which they are called anxious, in others irascibility, which differs from anger. To be irascible is one thing, to be angry another, as anxiety differs from anguish; for not all who are anguished at times are anxious, nor are those who are anxious always anguished, just as there is a difference between drunkenness and a habit of drunkenness, and to be a lover is one thing, to be in love another. And this proneness of different men to different diseases ranges widely, for it pertains to all the disturbances;

\ciceroSection{28} it appears in many vices too, but the thing has no name. So men are called envious, malevolent, lustful, timid, and given to pity, because they are prone to those disturbances, not because they are always carried away by them. This proneness, then, of each man to his own particular kind, may from the likeness of the body be called a sickness, provided it be understood as a proneness to falling sick. But in good things — since some men are better fitted for some goods — let it be named aptitude; in bad things, proneness, to signify a tendency to slip; in neutral things, let it keep the former name. As there is in the body disease, sickness, and defect, so it is in the soul. They call disease the corruption of the whole body; sickness, disease together with weakness; defect, when the parts of the body are at variance with one another, from which come distortion of the limbs, contortion, deformity.

\ciceroSection{29} And so those two, disease and sickness, arise from a shaking and disturbance of the whole health of the body, while a defect is discerned by itself even when the health is sound. But in the soul we can separate disease from sickness only in thought, whereas faultiness is a settled disposition or condition, inconsistent throughout a whole life and at odds with itself. Thus it comes about that in the one case disease and sickness are produced by a corruption of beliefs, in the other by inconsistency and self-contradiction. For not every fault has equal divisions within it: in men who are not far from wisdom there is indeed a condition that disagrees with itself so long as it is unwise, but it is neither distorted nor depraved. Diseases and sicknesses, then, are parts of faultiness; but whether the disturbances are parts of the same is a question.

\ciceroSection{30} For faults are abiding conditions, while disturbances are in motion, so that they cannot be parts of abiding conditions. And just as in evils the likeness of the body touches the nature of the soul, so it does in goods. For in the body there are things of the first rank — beauty, strength, health, firmness, swiftness — and there are likewise such things in the soul. For as a sound proportion of the body, when those elements of which we are composed agree among themselves, is called health, so it is called health in the soul when its judgments and beliefs are in concord; and this is the virtue of the soul, which some say is temperance itself, others a thing that obeys the precepts of temperance and follows after it and has no proper appearance of its own — but whether it is the one or the other, it exists in the wise man alone. Yet there is a certain health of the soul that can fall even to the unwise, when by the treatment and cleansing of physicians the confusion of the mind is taken away.

\ciceroSection{31} And as the body has a certain fitting arrangement of its parts together with a certain pleasantness of color, and this is called beauty, so in the soul an evenness and constancy of beliefs and judgments, together with a certain firmness and stability, following upon virtue or containing the very force of virtue, is called beauty. Likewise, by words like those for the strength and sinews and effectiveness of the body, the strengths of the soul are named. The swiftness of the body is called speed, and this same thing is reckoned a merit of the intellect as well, because of the soul's running through many things in a brief span of time. This much is unlike between souls and bodies: that vigorous souls cannot be assailed by disease, but bodies can; yet the body's mishaps can happen without blame, those of the soul not so, for all its diseases and disturbances arise from a contempt of reason. And so they exist in human beings alone; for the beasts do something similar, but do not fall into disturbances.

\ciceroSection{32} Between the keen and the dull, however, there is this difference: that the gifted, like Corinthian bronze taking the rust, both fall into disease more slowly and recover more quickly, whereas the dull do not. Nor indeed does the gifted soul fall into every disease and disturbance; for there are none that are savage and monstrous, while certain ones have even a first appearance of humanity, such as pity, distress, and fear. Sicknesses and diseases of the soul, however, are thought to be harder to root out than those highest faults that are contrary to the virtues; for while the diseases remain, the faults cannot be removed, since they are not healed so quickly as the faults are taken away.

\ciceroSection{33} You now have what the Stoics argue, with their fine analysis, about the disturbances — what they call the logical part \textit{[Greek: logika]}, because it is treated more subtly. And since our discourse has now sailed clear of these, as out of jagged reefs, let us hold to the course of the rest of the discussion — provided that we have spoken on those points clearly enough, given the obscurity of the matter. — Quite clearly enough; but if any points must be examined more carefully, we shall inquire into them another time. For now we await the sails you spoke of just now, and the course.

\ciceroSection{34} Since — as I have said in other places about virtue, and shall often have to say, for most of the questions that bear on life and conduct are drawn from the spring of virtue — since, then, virtue is a constant and harmonious disposition of the soul, making praiseworthy those in whom it resides, and itself praiseworthy in its own right, of its own accord, even apart from any usefulness, from it there proceed honorable wishes, judgments, actions, and all right reason (though virtue itself, most briefly, may be called right reason). Contrary to this virtue, then, is faultiness — for so I prefer to name it rather than malice, which is what the Greeks call \textit{[Greek: kakia]}; for malice is the name of some particular fault, faultiness of all of them — and from this the disturbances are stirred up, which are, as I said a little before, turbid and agitated motions of the soul, turned away from reason and most hostile to a tranquil mind and life. For they bring on anxious and bitter distresses, and afflict and weaken souls with fear; the same disturbances inflame them with excessive craving, which we call now cupidity, now desire — a kind of impotence of the soul utterly at variance with temperance and moderation.

\ciceroSection{35} And if ever that craving has obtained the thing it longed for, then it is so carried away by elation that there is no consistency in what it does — like the man who judges excessive pleasure of the soul to be the highest error. The healing of those evils, then, lies in virtue alone. And what is there that is not only more wretched, but even fouler and more deformed, than a man cast down, weakened, prostrate by distress? Next to this wretchedness is the man who, fearing some approaching evil, hangs in suspense and is half dead with dread. To signify the force of this evil, the poets make a rock hang over Tantalus among the dead, for his crimes and his arrogance of soul and his haughty speech. This is the common penalty of folly. For over all whose minds shrink from reason there always hangs some such terror.

\ciceroSection{36} And as these are the wasting disturbances of the mind — I mean distress and fear — so the more cheerful ones, cupidity ever greedily reaching after something, and empty elation, that is, exultant gladness, do not differ much from madness. From this it is understood what manner of man he is whom we call now self-controlled, now modest and temperate, now constant and continent; sometimes we wish to refer these same terms, as to a head, to the name of frugality. For unless the virtues were contained under that name, that saying would never have become so widely current as now to hold the place of a proverb — that a frugal man does all things rightly. When the Stoics say the same of the wise man, they seem to speak too marvelously and too magnificently.

\ciceroSection{37} Therefore this man, whoever he is, who by moderation and constancy is quiet in soul and at peace with himself, so that he neither wastes away with troubles nor is broken by fear nor burns with desire in some thirsting pursuit nor melts away in exultation at empty elation — this man is the wise man we are seeking, this is the happy man, to whom nothing in human affairs can seem either unbearable enough to crush his spirit or joyful enough to carry him away. For what can seem great in human affairs to one who knows all eternity and the vastness of the whole world? What, either in the pursuits of men or in so brief a span of life, can seem great to the wise man, who keeps such constant watch with his soul that nothing can befall him unforeseen, nothing unexpected, nothing wholly new?

\ciceroSection{38} And this same man directs so keen a gaze in every direction that he always sees a seat and place for living without trouble and anguish, so that, whatever chance fortune may bring upon him, he bears it aptly and quietly. The man who does this will be free not only from distress but from all the other disturbances as well. And a soul empty of these makes men perfectly and absolutely happy, while the same soul, stirred up and dragged away from sound and certain reason, loses not only its constancy but its health besides. For this reason the reasoning and discourse of the Peripatetics is to be thought soft and unnerved, who say that souls must necessarily be disturbed, but apply a certain measure beyond which one ought not to go.

\ciceroSection{39} You apply a measure to a fault? Or is it no fault not to obey reason? Or does reason teach too little when it says that the thing is not good which you either crave ardently or, having gained it, exult over insolently, and again not evil, by which you either lie crushed or, to avoid being crushed, can scarcely keep your mind steady — and that all these things become either too sad or too glad through error, which error, if in fools it is thinned out by time, so that, though the matter remains the same, they bear old things differently from fresh ones, does not touch the wise man at all?

\ciceroSection{40} For what, after all, will that measure be? Let us seek the measure of distress, on which most of the labor is spent. It is written in Fannius that Publius Rupilius bore his brother's rejection for the consulship hard. But still he seems to have passed the measure, since for that reason he departed from life; he ought, then, to have borne it more moderately. And what if, while he bore that with moderation, the death of his children had come upon him? A new distress would have been born, but a moderate one. Yet a great addition would have been made. What if then came grievous pains of the body, the loss of property, blindness, exile? If distresses were added for each separate evil, the sum would become one that could not be endured.

\ciceroSection{41} The man who seeks a measure for a fault, then, acts much as if he supposed that a man who had flung himself from Leucate could check himself when he wished. For as he cannot, so a soul disturbed and incited can neither restrain itself nor halt where it wishes. And in general, things that are ruinous when they grow are also faulty when they are born;

\ciceroSection{42} but distress and the other disturbances, once enlarged, are certainly destructive: therefore even at the moment they are taken up they are already largely in the realm of destruction. For they drive themselves on, once reason has been left behind, and weakness indulges itself, and is carried out unawares into the deep, and finds no place to halt. For this reason there is no difference whether men approve moderate disturbances or moderate injustice, moderate cowardice, moderate intemperance; for whoever sets a measure to faults takes upon himself a part of the faults — which thing, hateful in itself, is the more troublesome because faults are on slippery ground, and once set in motion they slide down the incline and can in no way be held back. What of the fact that those same Peripatetics say that the disturbances we think should be uprooted are not only natural but even given by nature to good purpose?

\ciceroSection{43} Their argument runs like this. First, in many words they praise irascibility; they call it the whetstone of courage; they say that the onset of angry men against an enemy and against a bad citizen is far more forceful, and that the petty little reasonings of those who think along these lines — "this battle is rightly fought," "it is fitting to contend for the laws, for liberty, for the fatherland" — have no force at all unless courage has caught fire from anger. Nor indeed do they argue about warriors only: they hold that no command can be severe enough without some bitterness of irascibility; and the orator, in the end, they approve neither in accusing nor even in defending without the stings of irascibility — which, even if it is not present, they think must still be feigned in word and gesture, so that the orator's delivery may kindle the listener's anger. They deny, in short, that a man seems a man if he does not know how to be angry, and what we call gentleness they call by the faulty name of sluggishness.

\ciceroSection{44} Nor indeed do they praise this desire alone — for anger is, as I just defined it, the desire to take vengeance — but they say that this very kind of desire, or craving, was given by nature for the highest usefulness; for no one, they hold, can do anything splendidly except what gives him pleasure. Themistocles used to walk abroad at night, because he could not get to sleep, and to those who asked why he answered that he was roused from sleep by the trophies of Miltiades. Who has not heard of the vigils of Demosthenes? — who said he grieved whenever the early-rising industry of the craftsmen got the better of him. The very founders of philosophy, in the end, could never have made such great progress in their studies without a blazing craving for it. We are told that Pythagoras, Democritus, and Plato traveled to the ends of the earth; for they judged that wherever there was anything that could be learned, there they must go. Can we suppose all this could have been done without the highest ardor of craving?

\ciceroSection{45} Distress itself — which we have said must be fled like a foul and monstrous beast — they say was established by nature not without great usefulness, so that men, when affected by chastisements, reproaches, and disgraces, should feel pain at their own wrongdoing. For impunity for their offenses seems granted to those who bear disgrace and infamy without pain; it is better to be bitten by conscience. From this comes that line drawn from life by Afranius — for when the dissolute son cries, \textit{Ah, wretched me!}, the stern father answers, \textit{So long as he feels some pain, let him feel pain at whatever you please.}

\ciceroSection{46} The remaining kinds of distress too, they say, are useful: pity, for bringing aid and relieving the calamities of the undeserving; and that very emulation and disparagement is not useless, when a man sees either that he himself has not attained the same thing as another, or that another has attained the same thing as he. Indeed, if anyone were to do away with fear, then all carefulness in life would be done away with — the carefulness that is at its height in those who fear the laws, the magistrates, poverty, disgrace, death, pain. Yet they argue all this in such a way as to admit that these things must be pruned back, while denying that they either can or need to be uprooted entirely; and in nearly all matters they think the mean is best. When they set this out, do they seem to you to be saying nothing, or to be saying something? — To me, indeed, to be saying something; and so I am waiting to hear your reply to it. — I shall perhaps find one, but first this.

\ciceroSection{47} Do you see how great was the modesty of the Academics? For they say plainly what bears on the matter: the Stoics make answer to the Peripatetics; let those two fence it out, for all I care, since I have no need but to search out where that lies which seems likeliest to the truth. What, then, is there that meets us in this inquiry, from which something resembling the truth can be reached, beyond which the human mind cannot advance? The definition of a disturbance, which Zeno, I think, employed rightly. For he defines it thus: a disturbance is a movement of the soul turned away from reason and contrary to nature; or, more briefly, a disturbance is an appetite too vehement — and the too-vehement is to be understood as that which stands far off from the constancy of nature.

\ciceroSection{48} What could I say against these definitions? And most of this is the work of men reasoning prudently and acutely; that other talk belongs to the parade of the rhetoricians — "the fires of the spirit," "the whetstones of the virtues." Is it really so that a brave man, unless he has begun to seethe, cannot be brave? That is a gladiator's notion. And yet even in gladiators themselves we often see composure: they confer, they meet, they ask something, they make demands, so that they seem more at peace than angry. But in that line of work let there be, by all means, some Pacideianus of the temper that Lucilius reports: \textit{I will kill him, mark you, and win the day, if that is what you ask}, he says; \textit{but I expect it will go thus — I shall first take his blade in my own face before I plant my sword in the madman's belly and lungs. I hate the man; I fight in anger, and there is no longer any wait for us than while each fits the sword to his right hand; so utterly am I carried away by zeal and by hatred of him, by rage.} Yet without this gladiatorial irascibility we see Ajax in Homer come forward with much cheerfulness, when he was about to fight it out with Hector;

\ciceroSection{49} as he took up his arms, his very stride brought gladness to his comrades and terror to the enemy, so that Hector himself — as it stands in Homer — trembling in his whole breast, repented of having challenged him to combat. And these two, having conferred between themselves before they came to blows, did so gently and quietly, and did nothing angrily or madly even in the fight itself. As for that famous Torquatus, who won this cognomen, I do not suppose it was in anger that he stripped the torc from the Gaul; nor that Marcellus at Clastidium was brave because he was angry.

\ciceroSection{50} About Africanus indeed, since he is better known to us through the recent memory of him, I could even swear that he was not then inflamed with irascibility, when in the line of battle he covered Marcus Allienus the Pelignian with his shield and drove his sword into the enemy's breast. About Lucius Brutus I might perhaps be in doubt whether, on account of his boundless hatred of the tyrant, he charged the more unbridledly against Arruns; for I see that the two fell hand to hand by mutual blows. Why, then, do you bring in anger here? Does courage, unless it has begun to rave, have no impulses of its own? What of Hercules, whom that very thing you would have to be irascibility raised to heaven by his courage — do you suppose he closed with the Erymanthian boar or the Nemean lion in anger? Or did even Theseus seize the horns of the Marathonian bull in anger? Take care lest courage be in no way raging, and irascibility be wholly a thing of frivolity.

\ciceroSection{51} For there is no courage at all that is without reason. Human affairs must be despised, death must be counted as nothing, pains and toils must be reckoned bearable — when these things are settled by judgment and resolve, then there is that robust and steadfast courage; unless perhaps we suspect that whatever is done forcefully, sharply, spiritedly is done in anger. To me not even that famous Scipio, the chief pontiff, who declared this Stoic doctrine to be true — that the wise man is never a private citizen — seems to have been angry with Tiberius Gracchus then, when he abandoned the faltering consul and, himself a private man, as though he were consul, bade those who wished the commonwealth safe to follow him.

\ciceroSection{52} I do not know whether we ourselves have done anything bravely in the commonwealth; but if we have done anything, we certainly did not do it in anger. Or is there anything more like madness than anger? Ennius rightly called it the beginning of madness. The color, the voice, the eyes, the breath, the lack of control over what is said and done — what part of these has any soundness? What is fouler than Homer's Achilles, what than Agamemnon in the quarrel? As for Ajax, anger led him to frenzy and to death. Courage, then, does not call for irascibility to assist it; it is sufficiently equipped, prepared, and armed of itself. For in that way one might say drunkenness is useful for courage, and madness useful too, since both the insane and the drunk often do many things more vehemently. Ajax was always brave, yet bravest in his frenzy; for \textit{he did the greatest deed, when, as the Greeks were giving way, he saved the day by his own hand.} He restored the battle in his madness.

\ciceroSection{53} Shall we say, then, that madness is useful? Examine the definitions of courage: you will understand that it has no need of seething. Courage, then, is a disposition of the soul obedient to the supreme law in enduring things; or the preservation of a steadfast judgment in undertaking and repelling those things which seem fearful; or the knowledge of things fearful, and of their opposites, and of those wholly to be disregarded, which keeps a steadfast judgment about such matters; or, more briefly, as Chrysippus has it — for the earlier definitions were those of Sphaerus, a man among the foremost at defining well, as the Stoics think; for they are all on the whole much alike, but they bring out our common notions, some more, some less — how, then, does Chrysippus put it? Courage, he says, is the knowledge of things to be endured, or a disposition of the soul, in suffering and enduring, obedient to the supreme law without fear. However much we may harry these men, as Carneades used to, I fear they are the only true philosophers. For which of those definitions does not lay open that notion of ours which we all hold, covered and wrapped up, about courage? And once it is laid open, who is there who would look for anything more in the warrior, the commander, or the orator, or suppose that they can do anything bravely without raving?

\ciceroSection{54} What? Do not the Stoics, who say that all the foolish are mad, gather up these very points? Take away the disturbances, and irascibility above all: they will at once seem to be talking monstrosities. But as it is, they argue thus: they say that all fools are mad in the same sense that all filth stinks. — But not always. — Stir it: you will smell it. So too the irascible man is not always angry; provoke him, and presently you will see him in a fury. What? That warlike irascibility, when it has come home, what is it like with the wife, with the children, with the household? Is it useful then too? Is there anything, then, that a disturbed mind can do better than a steady one? Or can anyone be angry without a disturbance of the mind? Rightly, then, did our people, since all vices lay in men's characters, and none was fouler than irascibility, name the irascible alone "ill-tempered."

\ciceroSection{55} For an orator to be angry is least of all becoming; to feign anger is not unbecoming. Do I, in your eyes, seem angry then, when I speak somewhat sharply and forcefully in a case? What of this — when, the matters being now done and past, we write out our speeches, do we write in anger? Does anyone suppose so? \textit{Conquer!} — but do you imagine Aesopus ever acted in anger, or that Accius wrote in anger? Those things are performed splendidly, and indeed better by the orator, if only he is an orator, than by any actor — but they are performed gently and with a tranquil mind. And to praise desire — what kind of desire is it that does that? You bring forward to me Themistocles and Demosthenes, you add Pythagoras, Democritus, Plato. What? Do you call studies desire? These, even of the best things, such as the ones you bring forward, ought nonetheless to be settled and tranquil. And to praise distress, the one thing most of all to be loathed — to what philosophers, in the end, does this belong? But Afranius put it aptly: \textit{So long as he feels some pain, let him feel pain at whatever you please.} For he said it of a ruined and dissolute youth, whereas we are inquiring about a steadfast and wise man. And as for that very anger, let the centurion have it, or the standard-bearer, or the rest of those about whom there is no need to speak, lest we lay open the mysteries of the rhetoricians. For it is useful for one who cannot use reason to use a movement of the soul. But we, as I keep declaring, are inquiring about the wise man.

\ciceroSection{56} But emulation too, they say, is useful — and disparagement, and pity. Why should you pity rather than bring aid, if you can do it? Can we not be generous without pity? For we ourselves ought not to take on distresses for the sake of others, but rather to relieve others, if we can, of their distress. And to disparage another, or to be emulous with that faulty emulation which is akin to rivalry — what usefulness has it, when it belongs to the emulous man to be vexed at another's good which he himself does not have, and to the disparager to be vexed at another's good simply because another has it too? How could it be approved, to take on distress for the sake of gaining experience, if one wished to have a thing? For merely to wish to have it for oneself alone is the height of madness. And who could rightly praise the means of evils?

\ciceroSection{57} For who, in whom there is desire or craving, can fail to be lustful and grasping? In whom there is anger, fail to be irascible? In whom there is anguish, fail to be anxious? In whom there is fear, fail to be timid? Are we to suppose, then, that the wise man is lustful and irascible and anxious and timid? Of his preeminence much can indeed be said, at whatever length and breadth one likes, but most briefly in this way: that wisdom is the knowledge and understanding of things divine and human, and of the cause of each particular thing; from which it follows that he imitates the divine and counts all things human inferior to virtue. Into this man, then, as into a sea that lies open to the winds, you said that disturbance seemed to you to fall? What is there that could disturb so great a gravity and steadiness? Something unforeseen, perhaps, or sudden? But what of that kind can befall a man to whom nothing that can happen to a human being is unrehearsed? As for their saying that the excessive ought to be pared away and the natural left alone — what, after all, can be natural that can at the same time be excessive? For all these things have sprung from the roots of error, and they must be torn out and dragged up by the roots, not trimmed around the edges nor lopped off.

\ciceroSection{58} But since I suspect that you are asking not so much about the wise man as about yourself — for him you suppose to be free of every disturbance, and you wish the same for yourself — let us see how great are the remedies that philosophy applies to the diseases of the soul. For there is certainly a kind of medicine; nor was nature so hostile and so unfriendly to the human race that, having found so many things salutary for the body, she found none for the soul. Toward souls she has even deserved better in this: that the body's aids are applied from without, while the soul's health is enclosed within the souls themselves. But the greater and the more divine their excellence, the greater the care they require. And so reason well applied discerns what is best; neglected, it is entangled in many errors.

\ciceroSection{59} To you, then, I must now turn my whole discourse; for you pretend to ask about the wise man, but you are asking, perhaps, about yourself. The cures, then, of those disturbances which I have set out are various. For not every distress is calmed by a single method — one medicine must be applied to the man who mourns, another to the man who pities or who envies. And in all four disturbances there is also this distinction: whether the discourse is better directed at disturbance in general — which is a spurning of reason, or a too-vehement appetite — or at the particular passions, as at fear, desire, and the rest; and whether the better course is to argue that the thing from which the distress was taken up should not be borne hard, or that distress over anything whatever is to be removed altogether. So, if a man bears it hard that he is poor, the question is whether you should argue that poverty is no evil, or that a man ought to bear nothing hard. Surely this latter is better, lest, if you happen not to convince him about poverty, distress must be conceded; but once distress is removed by its own proper arguments, those we used yesterday, in a certain way the evil of poverty too is taken away.

\ciceroSection{60} But every disturbance of this kind may be washed away by that appeasement of the mind which comes when you teach that the thing from which gladness or desire arises is not a good, nor that from which fear or distress arises an evil. Yet still the sure and proper healing is this: if you teach that the disturbances are in themselves faulty and have in them nothing either natural or necessary — as we see distress itself made lighter when we cast at the grieving the reproach of a weak and unmanned spirit, and when we praise the gravity and steadiness of those who bear what is human without turmoil. This is wont to happen even to those who hold such things to be evils, yet think they must nonetheless be borne with an even mind. One man supposes pleasure to be a good, another money; yet the one can be called off from intemperance and the other from greed. But that other method and discourse, which at once removes the false opinion and draws off the distress, is the more useful of the two, though it rarely succeeds and is not to be applied to the common run of men.

\ciceroSection{61} There are, moreover, certain distresses which that medicine can in no way relieve — as, if a man bears it hard that there is in him no virtue, no spirit, no sense of duty, no honor; he is indeed tormented on account of evils, but some other cure must be brought to bear on him, and one of such a kind as could be common even to all the philosophers who differ on other matters. For all must agree that the agitations of the soul which are turned away from right reason are faulty, so that, even if those things were evils which stir fear or distress, or goods which stir craving or gladness, the agitation itself is nonetheless faulty. For we wish that man whom we call high-souled and brave to be steady, calm, grave, despising all that is human. But no one can be such while mourning, or fearing, or craving, or transported with delight. For these are the marks of those who count the outcomes of human affairs above their own souls.

\ciceroSection{62} For this reason, as I said before, there is for all the philosophers one method of healing: that nothing should be said about the character of the thing which disturbs the soul, but only about the disturbance itself. And so first, in craving itself, when the sole aim is that it be removed, the question is not whether the thing that stirs desire is a good or not, but desire itself is to be removed; so that, whether the honorable is the highest good, or pleasure, or the two of them joined, or those three classes of goods, then, even if the appetite were for virtue itself but too vehement, the same discourse is to be applied to all, to deter them. The whole calming of the soul, moreover, lies in setting human nature plainly in view; and that this may be the more easily discerned in clear outline, the common condition and the law of life must be unfolded in discourse.

\ciceroSection{63} And so it was not without cause that, when Euripides was producing the play \textit{Orestes}, Socrates is said to have called for the first three verses to be repeated: \textit{No speech is so terrible to tell, no stroke of fortune, no calamity sent down by the wrath of the gods, that human nature cannot bear and endure it.} And to persuade a man that what has befallen can and ought to be borne, a roll-call of those who have borne it is useful. Although the calming of distress was both set out in yesterday's discussion and treated in the book of \textit{Consolation}, which — for we were no wise men — I wrote in the midst of grief and sorrow; and what Chrysippus forbids, the applying of a remedy to the still-fresh swellings of the soul, that I did, and forced nature, so that the greatness of the sorrow might yield to the greatness of the medicine.

\ciceroSection{64} But to distress, of which enough has been argued, fear is neighbor, and of it a few words must be said. For fear is to a future evil what distress is to a present one. And so some called fear a certain part of distress, while others named fear a foretaste of trouble, since it was, as it were, the leader of the trouble to follow. By whatever methods, then, present things are borne, by those same the things to come are despised. For in both we must take care that we do nothing low, submissive, soft, unmanned, broken, abject. But although something must be said about the inconstancy, the weakness, the lightness of fear itself, it is of great profit to despise the very things that are feared. And so, whether it happened by chance or by design, it fell out most conveniently that on the first and the next day there was a discussion of those things which are most feared, death and pain. If these have been proved, we are in great part set free from fear. And of the opinion of evils, thus far.

\ciceroSection{65} Let us see now about the opinion of goods, that is, about gladness and craving. To me, indeed, in the whole reasoning that bears on the soul's disturbance, one fact seems to hold the cause: that all these disturbances are in our power, all taken up by judgment, all voluntary. This error, then, must be torn out, this opinion drawn off; and just as in supposed evils the things believed to be evils must be made bearable, so in goods the things counted great and gladdening must be made calmer. And this much is common to evils and goods: that, even if it is now hard to convince a man that none of the things which disturb the soul are to be held either among goods or among evils, still a different cure must be applied to a different motion, and the envious man corrected by one method, the lover by another, again the anxious man by another, the timid by yet another.

\ciceroSection{66} And it would be easy, following that method which is most approved concerning goods and evils, to deny that the fool can ever be affected with gladness, since he never has any good; but we are now speaking in the common way. Granted that those things are goods which are thought to be — honors, riches, pleasures, the rest — still, in the very enjoying of them, a gladness that exults and leaps about is shameful, just as, if laughter is permitted, a burst of cackling is nonetheless to be blamed. For the same fault lies in the soul's outpouring in gladness as in its contraction in sorrow, and craving in pursuing has the same lightness as gladness in enjoying; and just as those too cast down by trouble, so those too lifted up by gladness are rightly judged frivolous; and since to envy belongs to distress, and to take pleasure in others' misfortunes belongs to gladness, each is wont to be chastised by setting before it a certain savagery and brutishness; and just as it is becoming to be wary but not to fear, so it is becoming to rejoice but not to be glad, since for the sake of teaching we distinguish gladness from rejoicing.

\ciceroSection{67} I have already said above that the contraction of the soul can never rightly come to be, but its elation can. For the Hector of Naevius rejoices in one way: \textit{Glad am I to be praised by you, father, by a man who is himself praised} — and the man in Trabea rejoices in another: \textit{The bawd, won over by silver, will watch for my nod, for what I wish, what I desire. As I come, with a finger I will tap the door; the doors will swing open. When Chrysis catches sight of me unawares, she will come eager to meet me, longing for my embrace, and give herself to me.} How fine he thinks this, he will say himself: \textit{I shall outstrip Fortune herself in my fortunes.}

\ciceroSection{68} How shameful this gladness is, anyone who pays careful and close attention can see to the bottom. And just as those are shameful who are carried away by gladness when they are enjoying the pleasures of Venus, so disgraceful are those who, with inflamed soul, crave them. But that whole thing which is commonly called love — nor, by Hercules, do I find by what other name it could be called — is of such great lightness that I see nothing I would think comparable to it. Caecilius reckons that a man who does not hold it the highest god is either a fool or ignorant of things: \textit{In whose hand it lies whom he would have be mad, whom be sound of mind, whom be made whole, whom be cast into sickness; whom he would have loved in return, whom sought after, whom summoned.} O what an excellent corrector of life is poetry, which thinks love, the author of disgrace and lightness, should be set in the council of the gods!

\ciceroSection{69} I am speaking of comedy, which, if we did not approve of these disgraces, would not exist at all. What says, out of tragedy, that prince of the Argonauts? \textit{You saved me for love's sake rather than for honor's.} What then? This love of Medea — how great a blaze of miseries it kindled! And yet, in another poet, she dares to say to her father that she had taken for husband that man \textit{whom Love had given, who prevails more and is mightier than a father.}

\ciceroSection{70} But let us allow the poets their play, in whose tales we see Jupiter himself caught up in this disgrace; let us come to the philosophers, the masters of virtue, who deny that love is a thing of debauchery, and on this point quarrel with Epicurus, who, as my opinion goes, does not lie much. For what is this love of friendship? Why does no one love an ugly youth, nor a handsome old man? This custom, to my mind, was born in the gymnasia of the Greeks, in which such loves are free and licensed. Well, then, did Ennius say: \textit{The beginning of disgrace is the baring of bodies among fellow citizens.} And these loves, even granting that they are chaste — as I see they can be — are nonetheless full of trouble and anxiety, and the more so because they hold and curb themselves.

\ciceroSection{71} And — to pass over loves of women, to which nature has granted a wider license — who is in doubt about the rape of Ganymede, what the poets mean by it, or fails to understand what Laius both says and craves in Euripides? And then, what do the most learned men and the greatest poets publish of themselves in their poems and their songs? A brave man, well known in the affairs of his own state — and yet what verses Alcaeus writes about the love of young men! As for Anacreon, his entire poetry is amorous. But that Ibycus of Rhegium burned with love more fiercely than any of them is plain from his writings. And the loves of all these men we see to be born of lust. We philosophers have arisen — and indeed on the authority of our own Plato, whom Dicaearchus accuses not unjustly — to lend authority to love.

\ciceroSection{72} The Stoics, indeed, say that the wise man will love, and they define love itself as an effort to form a friendship out of the appearance of beauty. If there exists anywhere in nature such a love, free of anxiety, free of longing, free of care, free of sighing, then by all means let it be; for it is empty of all desire — and our discussion here is about desire. But if there is some love — and surely there is — which is little or nothing short of madness, of the sort there is in the \textit{Leucadia}: \textit{if indeed there be any god to whom I am a care} —

\ciceroSection{73} but this it was the business of all the gods to see to, the way this man should enjoy the pleasure of his love! \textit{Wretched me!} — nothing truer. And the other answers him aptly: \textit{Are you in your right mind, to wail so wildly?} So he seems mad even to his own people. And what tragedies he stages! \textit{You, holy Apollo, bring me aid; and you, all-powerful Neptune, I call upon; and you besides, you Winds!} He thinks the whole world will turn itself to relieving his love — Venus alone he shuts out as unfair: \textit{for why should I call on you, Venus?} He says that out of lust she cares for nothing — as though he himself were not, out of lust, both doing and saying such monstrous things.

\ciceroSection{74} — for a man in this condition, then, this is the cure to apply: that it be shown him how trivial, how contemptible, how utterly worthless is the thing he craves; how easily it may be obtained either from another source or in another way, or else passed over altogether. He must also sometimes be drawn off to other pursuits, concerns, cares, occupations; and finally, like a sick man who is not recovering, he must often be cured by a change of place.

\ciceroSection{75} Some even think that an old love must be driven out by a new one, as one nail drives out another. But above all he must be reminded how great a frenzy love is. For of all the disturbances of the mind there is surely none more violent — so that, even if you should be unwilling to bring against it those very charges, I mean defilements and seductions and adulteries, and finally incests, the foulness of all of which deserves accusation — yet, to leave these aside, the very disturbance of the mind in love is foul in itself.

\ciceroSection{76} For, to pass over those things that belong to its frenzy, consider the very triviality of those features in it that seem moderate: \textit{Injuries, suspicions, enmities, a truce, war — then peace again! If you should ask to make these uncertain things certain by sure reasoning, you would accomplish no more than if you set yourself to go mad by reason.} This inconstancy and changeableness of mind — whom would its very perversity not deter? And here too must be demonstrated what is said of every disturbance: that there is none that is not a matter of opinion, not taken up by judgment, not voluntary. For if love were natural, then all men would love, and would love always, and would love the same object; nor would shame deter one, reflection another, satiety another. As for anger — anger, so long as it disturbs the mind, leaves no room for doubt that it is madness; at whose impulse a quarrel of this sort breaks out even between brothers:

\ciceroSection{77} \textit{What man on earth has ever surpassed you in shamelessness? And who has surpassed you in malice?} — you know what follows; for in alternating verses the gravest insults are hurled between the brothers, so that it is easy to see they are the sons of Atreus — of him who plots a new kind of vengeance against his brother: \textit{A heavier mass, a greater evil I must brew, to crush and break that bitter heart of his.} To what, then, does this mass burst forth? Hear Thyestes: \textit{My own brother urges me on, that I, wretched, should chew my own sons between my jaws} — and serves him their flesh. For is there anything to which anger does not press on, just as frenzy does? And so we say properly of angry men that they have gone out of their power — that is, out of counsel, out of reason, out of mind; for the power over these ought to extend over the whole soul.

\ciceroSection{78} From such men those against whom they attempt to make their assault must either be withdrawn, until they collect themselves again — and what is collecting oneself but gathering the scattered parts of the mind back into their place? — or they must be begged and entreated that, if they have any power of taking vengeance, they put it off to another time, until the anger cools. And "to cool" surely signifies the heat of the mind stirred up against the will of reason. From this comes the praise of Archytas, who, when he had grown rather angry at his bailiff, said: "What a state I would have left you in, had I not been angry!"

\ciceroSection{79} Where, then, are those who call irascibility useful — can madness be useful? — or natural? Is there anything in accordance with nature that comes to be in defiance of reason? And how, if anger were natural, would one man be more irascible than another, or the desire to avenge come to an end before it had taken its vengeance, or anyone repent of what he had done in anger? As we see in the case of King Alexander, who, when he had killed his close companion Clitus, scarcely kept his hands from himself; so great was the force of his remorse. Once these things are known, who is there who can doubt that this movement of the mind too is wholly a matter of opinion and voluntary? For who could doubt that the diseases of the mind, such as avarice or the craving for glory, arise from the fact that the thing by which the mind is made sick is valued highly? From which it must be understood that every disturbance, too, lies in opinion.

\ciceroSection{80} And if confidence — that is, a firm assurance of the mind — is a kind of knowledge and a weighty opinion of one who does not assent rashly, then fear too is a distrust of an expected and impending evil; and if hope is the expectation of a good, fear must be the expectation of an evil. As fear, then, so the rest of the disturbances are concerned with evil. Therefore, as steadiness belongs to knowledge, so disturbance belongs to error. But those who are said to be by nature irascible, or compassionate, or envious, or anything of the kind, are constituted as if with an ill condition of mind — yet curable, as Socrates is said to have been. When Zopyrus, who professed to discern each man's nature from his form, had charged Socrates with many vices before a gathering, he was laughed at by the rest, who did not recognize those vices in Socrates; but he was rescued by Socrates himself, who said that they were indeed inborn in him, but had been cast out by reason.

\ciceroSection{81} Therefore, just as a man of the soundest constitution may yet seem by nature more inclined toward some particular illness, so one mind is more prone than another to particular vices. But of those who are said to be vicious not by nature but through fault, their vices consist of false opinions about things good and bad, so that one is more prone to some movements and disturbances than another. A long-settled habit, however, as in the body, is more painfully driven off than a disturbance; and a sudden swelling of the eyes is cured more quickly than a chronic inflammation is driven away.

\ciceroSection{82} But now that the cause of the disturbances is known — that they all arise from the judgments of opinions and from acts of will — let there now be an end to this discussion. We ought to recognize that, once the ends of good and evil are known, so far as they can be known by man, nothing greater or more useful can be wished from philosophy than the things we have discussed over these four days. For with death held in contempt and pain made lighter for the bearing, we have added the calming of distress, than which no evil is greater for a man. For although every disturbance of the mind is grievous and differs little from madness, still we are accustomed to say of the rest, when they are in some disturbance of fear or gladness or craving, only that they are agitated and disturbed; but of those who have surrendered themselves to distress, that they are wretched, afflicted, careworn, ruined.

\ciceroSection{83} And so it does not seem to have come about by chance, but to have been set forth by you on a plan, that we should discuss distress separately and the other disturbances apart; for in it lies the fountain and head of all miseries. But of distress and of the rest of the soul's diseases there is one and the same cure: that they are all matters of opinion and voluntary, and are undertaken because this seems to be the right thing to do. This error, as it were the root of all evils, philosophy promises to tear out by the very roots.

\ciceroSection{84} Let us give ourselves over, then, to her for cultivation, and let ourselves be healed. For while these evils sit lodged within us we cannot be happy — no, nor even sound. Either, then, let us deny that anything is accomplished by reason, when on the contrary nothing can be rightly done without reason — or, since philosophy consists in the marshalling of reasons, let us seek from her, if we wish to be both good and happy, every aid and help toward living well and happily.

\ciceroBook{5}

\ciceroSection{1} This fifth day, Brutus, will bring the Tusculan discussions to their close — the day on which we argued the very question that you, of all of them, most approve. For I have come to see, both from the book you wrote to me with the greatest care and from many conversations of yours, how thoroughly the view pleases you that virtue is sufficient of itself for the happy life. And though this is hard to prove, on account of fortune's torments, so various and so many, it is nonetheless the kind of thing at which we must labor, that it may be proved the more readily. For among all the matters that philosophy handles, there is nothing that is spoken of with greater weight or grandeur.

\ciceroSection{2} For since this was the cause that drove the first men to apply themselves to the study of philosophy — that, setting all else aside, they devoted themselves wholly to the search for the best condition of life — surely it was in the hope of living happily that they laid out such care and effort upon that study. And if virtue was discovered and brought to perfection by them, and if there is in virtue protection enough for living happily, who is there who would not judge that the work of philosophizing was nobly both undertaken by them and taken up by us? But if virtue, subject to changes various and uncertain, is the handmaid of fortune, and has not strength enough to keep itself secure, then I fear that we ought to lean toward the hope of living happily not so much by confidence in virtue as by the offering of prayers.

\ciceroSection{3} For my own part, when I consider with myself those reverses by which fortune has so violently exercised me, I begin at times to lose faith in this view, and to dread the weakness and frailty of the human race. For I fear that nature, having given us bodies that are feeble and joined to them diseases beyond cure and pains beyond endurance, may have given us minds as well — minds at once attuned to the pains of the body and, on their own account, entangled in their own anguish and distress.

\ciceroSection{4} But here I rebuke myself, because from the softness of others, and perhaps of my own, I form my estimate of virtue's strength, rather than from virtue itself. For virtue — if indeed there is any such thing as virtue, a doubt your uncle, Brutus, swept away — holds beneath itself all that can befall a man; looking down upon these, it scorns the accidents of mortal life, and, free of all fault, it judges that nothing concerns it but itself. We, on the other hand, magnify by fear every adversity as it comes, and by grief every present trouble, and choose to condemn the nature of things rather than our own error.

\ciceroSection{5} But the correction of this fault, and of the rest of our vices and sins, must all be sought from philosophy. Into her bosom, from the earliest season of our life, our own will and zeal drove us; and now, in these gravest reverses, tossed by a great tempest, we have fled for refuge to the same harbor from which we set out. O philosophy, guide of life! O searcher-out of virtue and expeller of vices! What, not only could we, but what could the whole life of man have been without you? You gave birth to cities; you called together into the fellowship of life men once scattered; you joined them to one another first by dwellings, then by marriages, then by the sharing of letters and of speech; you were the inventor of laws, you the teacher of morals and of order. To you we flee; from you we ask aid; to you we entrust ourselves — as before in great part, so now wholly and entirely. And one day well spent, by your precepts, is to be preferred to an eternity of wrongdoing.

\ciceroSection{6} Whose resources, then, should we use rather than yours — you who have lavished tranquillity of life upon us and taken away the terror of death? And yet philosophy, so far from being praised in proportion to what she has deserved of the life of man, is by most men neglected and by many even reviled. Does anyone dare to revile the parent of life, and to stain himself with this parricide, and to be so impiously ungrateful as to accuse her whom he ought to revere, even if he could not grasp her? But this error, I think, and this darkness spread over the minds of the unlearned, comes of their being unable to look so far back, and of their not believing that the men by whom the life of man was first set in order were philosophers.

\ciceroSection{7} Though we see the thing itself to be most ancient, we confess all the same that the name is recent. For wisdom itself — who can deny that it is ancient, not only in fact but even in name? It is the name that won, among the ancients, this most beautiful title by its knowledge of things divine and human, and of the beginnings and causes of each thing. And so those seven who by the Greeks were called \textit{[Greek: sophoi]}, and by our people were both held and named the wise; and many ages before them Lycurgus, in whose time Homer too is recorded to have lived, before this city was founded; and as far back as the age of heroes, Ulysses and Nestor — these, we are told, both were and were held to be wise.

\ciceroSection{8} Nor would Atlas have been said to bear up the sky, nor Prometheus to be fastened to the Caucasus, nor Cepheus set among the stars with his wife, his son-in-law, and his daughter, had not their divine knowledge of the heavens carried over their names into the error of fable. From these, in succession, all who set their zeal upon the contemplation of things were both held and named the wise; and that name of theirs ran on down to the age of Pythagoras. He, as Heraclides of Pontus writes — a pupil of Plato, and a man of the first rank in learning — is said to have come to Phlius, and there to have discoursed on certain matters with Leon, the chief man of the Phliasians, learnedly and at length. And when Leon, marvelling at his genius and eloquence, asked him in what art above all he placed his confidence, Pythagoras replied that he knew no art at all, but was a philosopher. Leon, struck by the novelty of the word, asked who these philosophers might be, and what was the difference between them and the rest of mankind;

\ciceroSection{9} and Pythagoras answered that the life of man seemed to him like that festival which was held with the most splendid display of games and the gathering of all Greece. For there, just as some men sought glory and the renown of a crown by training their bodies, while others were drawn by the profit and gain of buying and selling, but there was a certain kind of men — and this the most freeborn of all — who sought neither applause nor profit, but came for the sake of seeing, and looked on closely at what was being done, and in what manner: so too we, as if we had set out from some city to a kind of crowded festival, had come into this life from another life and nature; and some of us are slaves to glory, some to money, while there are a rare few who, holding all else as nothing, study closely the nature of things. These men, Pythagoras said, he called students of wisdom — that is, philosophers; and just as there it was most worthy of a free man to look on, acquiring nothing for himself, so in life the contemplation and knowledge of things far surpasses all other pursuits.

\ciceroSection{10} Nor indeed was Pythagoras only the inventor of the name, but an enlarger as well of the things themselves. After this conversation at Phlius he came into Italy, and adorned that Greece which has been called Great, both in private and in public, with the most excellent institutions and arts. Of his teaching there may, perhaps, be another time for speaking. But from the ancient philosophy down to Socrates — who had listened to Archelaus, the pupil of Anaxagoras — it was numbers and motions that were handled, and the question whence all things arise and into what they fall back; and those men inquired zealously into the magnitudes of the stars, their intervals and courses, and all things of the heavens. Socrates, however, was the first to call philosophy down from the sky, to set it in cities, to bring it even into homes, and to compel it to inquire about life and morals, about things good and evil.

\ciceroSection{11} His many-sided manner of arguing, the variety of his subjects, and the greatness of his genius, consecrated in the memory and writings of Plato, brought into being several kinds of dissenting philosophers; and of these we ourselves have followed above all the one which we judged Socrates to have used: that we should conceal our own opinion, free others from error, and in every discussion seek out what was most like the truth. This method, since Carneades held to it most acutely and most copiously, I have often elsewhere, and lately at Tusculum, made it my practice to argue after that fashion. And the discourse of the four days I have already sent you, written out in the earlier books; but on the fifth day, when we had taken our seats in the same place, this was set down as the matter we should argue:

\ciceroSection{12} — It does not seem to me that virtue can be sufficient for the happy life. — But it does, by Hercules, seem so to my friend Brutus, whose judgment, with your leave I shall say it, I rank far above yours. — I do not doubt it; nor is the question now how much you love him, but this, what the thing is which I have said seems true to me, and which I want you to argue. — You deny, then, that virtue can be sufficient for the happy life? — I deny it utterly. — What? For living rightly, honorably, laudably — in short, for living well — is there protection enough in virtue? — Certainly there is. — Can you, then, either not call wretched a man who lives badly, or deny that a man you admit lives well lives happily? — Why should I not? For even amid torments one can live rightly, honorably, laudably, and therefore well — provided only you understand what I now mean by "well." For I mean: steadfastly, with weight, wisely, bravely.

\ciceroSection{13} These things are flung even upon the rack, to which the happy life does not aspire. — What then? Is the happy life, I ask you, the one thing left outside the door and threshold of the prison, when steadfastness, weight, courage, wisdom, and the rest of the virtues are dragged off to the torturer and refuse neither punishment nor pain? — You, if you mean to do anything, must search out something new; these things do not move me in the least — not only because they are commonplace, but much more because, just as certain light wines have no strength in water, so these doctrines of the Stoics please more when tasted than when drunk. So this chorus of virtues, set upon the rack, sets images before our eyes of the most ample dignity, so that the happy life seems likely to press on toward them at a run, and not to suffer them to be deserted by it;

\ciceroSection{14} but once you have drawn your mind away from that picture and those images of the virtues to the thing itself and the truth, this is left bare: whether a man can be happy as long as he is being tortured. This, then, is what we must now ask; but do not fear that the virtues will protest and complain that they have been abandoned by the happy life. For if no virtue is without prudence, prudence itself sees this — that not all good men are also happy; and it recalls many things about Marcus Atilius, Quintus Caepio, Manius Aquilius, and the happy life, if it please us to use images rather than the things themselves, prudence itself holds back as it tries to mount the rack, and denies that it has anything in common with pain and torment.

\ciceroSection{15} — I am content to let you proceed in that way, though it is unfair of you to prescribe to me how you would have me argue. But I ask: are we to think that something was accomplished on the earlier days, or nothing? — Accomplished, surely, and a good deal. — And yet, if that is so, this question is already as good as won and brought almost to its conclusion. — How so, pray? — Because turbulent motions and agitations of the mind, roused and carried up by reckless impulse, repelling all reason, leave no part for the happy life. For who can fear death or pain — the one often present, the other always hanging over us — and not be wretched? What, moreover, if the same man, as commonly happens, fears poverty, disgrace, infamy; if he fears weakness, blindness; if, in the end, he fears that which has befallen not single men only but often whole peoples once mighty — slavery?

\ciceroSection{16} Can anyone be happy while he fears such things? And what of the man who not only fears them as things to come, but actually bears and endures them when present — add to these exile, mourning, the loss of children — the man who, broken by these blows, is crushed by grief: can he, after all, be anything but utterly wretched? And what of this? The man whom we see inflamed and raging with his lusts, reaching madly after everything with an appetite that nothing can fill, and the more abundantly he drinks in pleasures from every side, the more heavily and burningly he thirsts — would you not be right to call him utterly wretched? And what of the man lifted up by frivolity, exulting in empty joy and capering without cause — is he not all the more wretched the happier he thinks himself? Therefore, just as these men are wretched, so on the other side those are happy whom no fears terrify, whom no griefs eat away, whom no lusts goad, whom no idle exulting joys melt with their enervating pleasures. As the calm of the sea is recognized when not the slightest breath of air stirs the waves, so the soul's quiet and settled state is discerned when there is no disturbance by which it can be moved. And if there is a man who counts the force of fortune, and everything human that can befall anyone, as bearable — so that neither fear nor anguish can touch him — and if this same man craves nothing, is carried away by no empty pleasure of the soul, what reason is there why he should not be happy?

\ciceroSection{17} And if these things are brought about by virtue, what reason is there why virtue itself, of itself, should not make men happy? — Well, the one cannot be denied: that those who fear nothing, are troubled by nothing, crave nothing, are carried away by no ungoverned joy, are happy; and so I grant you that. But the other is no longer an open question. For it was established in our earlier discussions that the wise man is free of every disturbance of the soul. — Then surely the matter is settled;

\ciceroSection{18} for the inquiry seems to have reached its conclusion. — Nearly so, indeed. — And yet that is the way of the mathematicians, not of the philosophers. For when geometers wish to teach something, if anything bearing on the matter belongs among the things they have taught before, they take it as granted and proven, and explain only that about which nothing has been written before; whereas philosophers, whatever subject they have in hand, gather into it everything that bears upon it, even if it has been argued elsewhere. Were it not so, why would a Stoic, if the question were raised whether virtue can suffice for the happy life, say a great deal? It would be enough for him to answer that he had shown before that nothing is good except what is honorable, and that, this being proved, it follows that the happy life is content with virtue; and just as this follows from that, so the other follows from this — that, if the happy life is content with virtue, then, unless a thing is honorable, nothing else is good.

\ciceroSection{19} But still they do not proceed in that way; for there are separate books both on the honorable and on the highest good, and although it follows from the latter that there is power enough in virtue for living happily, none the less they treat this point separately. For each subject must be handled with its own proper arguments and reminders, a subject so great above all. For do not suppose that any voice more illustrious has been uttered in philosophy, or that any of philosophy's promises is richer or greater. For what does it profess? Good gods! That it will bring it about, for the man who has obeyed its laws, that he be forever armed against fortune, that he hold within himself every safeguard of living well and happily — that, in short, he be forever happy.

\ciceroSection{20} But I shall see what it can accomplish; in the meantime I value this very thing greatly, that it makes the promise. For Xerxes, indeed, glutted with all the rewards and gifts of fortune, not content with his cavalry, not with his foot soldiers, not with his multitude of ships, not with his boundless weight of gold, offered a prize to whoever should discover a new pleasure — with which pleasure itself he was not content either, for desire will never find an end; I could wish that we might draw out by a prize someone who would bring us something by which we might believe this more firmly.

\ciceroSection{21} — I could wish that too; but I have one small point I would press. For I agree that, of the propositions you have set down, the one follows from the other: that just as, if only what is honorable is good, it follows that the happy life is brought about by virtue, so, if the happy life lies in virtue, nothing is good but virtue. But your friend Brutus, on the authority of Aristo and Antiochus, does not hold this; for he thinks that there is some good besides virtue. — What then?

\ciceroSection{22} Do you suppose I shall speak against Brutus? — Just as you see fit; for to set limits beforehand is not my place. — What, then, is consistent with each view, we shall consider elsewhere. For on that matter I had a disagreement both with Antiochus often and with Aristo recently, when as commander I lodged with him at Athens. For it did not seem to me that anyone could be happy while he was in the midst of evils — and the wise man could be in the midst of evils, if there were any evils of the body or of fortune. These things were said — which Antiochus too wrote in several places — that virtue itself, of itself, can bring about the happy life, but not the happiest; further, that most things take their name from the greater part, even if some part should be missing — as strength, as health, as riches, as honor, as glory, which are reckoned by their kind, not by their number; and likewise that the happy life, even if it should be lame in some part, none the less holds its name from the much greater part.

\ciceroSection{23} It is not so necessary now to thresh these matters out, although they seem to me said not very consistently. For both, in the case of the man who is happy, I do not see what he should want in order to be happier — for if there is something he lacks, he is not even happy at all — and, as for their saying that each thing is named and regarded from the greater part, there is a place where this holds in that way; but when they say there are three kinds of evils, and a man is pressed by all the evils of two of those kinds — so that everything in his fortune is adverse, and his body is wholly overwhelmed and worn out with pains — shall we say that this man falls just a little short of the happy life, not to say the happiest?

\ciceroSection{24} This is the very thing that Theophrastus could not maintain. For when he had laid it down that lashings, tortures, torments, the overthrows of one's country, exiles, the loss of children have great force toward living badly and wretchedly, he did not dare to speak loftily and largely, since he felt humbly and meanly. How well he felt is not the question — consistently, at any rate, beyond doubt. And so it is not my habit to find fault with what follows, when you have granted the premises. But this man, the most elegant and most learned of all philosophers, is not greatly faulted when he says there are three kinds of goods; he is harried, however, by everyone, above all in that book which he wrote on the happy life, in which he argues at length why a man who is tortured, who is racked, cannot be happy. In it he is also thought to say that the happy life does not mount the wheel — that being a certain kind of torture among the Greeks. He nowhere says that at all, in fact, but what he does say comes to the same thing.

\ciceroSection{25} Can I, then, who have granted that bodily pains are among evils, that the shipwrecks of fortune are among evils, be angry with him for saying that not all good men are happy, when those things which he reckons among evils can fall upon all good men? This same Theophrastus is harried in both the books and the schools of all the philosophers, because in his \textit{Callisthenes} he praised that maxim: \textit{Fortune, not wisdom, rules our life.} They say that nothing more spineless was ever uttered by any philosopher. Rightly so, indeed; but I do not see that anything more consistent could have been said. For if there are so many goods in the body, so many outside the body in chance and fortune, is it not consistent that fortune — which is the mistress of all things both external and pertaining to the body — should have more power than counsel?

\ciceroSection{26} Or do we prefer to imitate Epicurus? Who often says many fine things; for how consistently and coherently he speaks with himself does not trouble him. He praises a frugal way of life. The remark of a philosopher, indeed — but only if a Socrates or an Antisthenes said it, not the man who declared pleasure to be the end of goods. He denies that anyone can live pleasantly unless he also lives honorably, wisely, and justly. Nothing weightier, nothing worthier of philosophy — were it not that he refers this very honorable, wise, and just living to pleasure. What is better than: that fortune intrudes but little upon the wise man? But does he say this — he who, after declaring pain to be not only the greatest evil but the only evil, can be overwhelmed in his whole body by the sharpest pains at the very moment when he most boasts against fortune? The same thing Metrodorus put in even better words:

\ciceroSection{27} I have forestalled you, he says, Fortune, and laid hold of you, and blocked up all your approaches, so that you cannot reach me. Splendid — if Aristo of Chios or the Stoic Zeno had said it, men who counted nothing evil except what was base; but you, Metrodorus, who lodged every good in the entrails and the marrow, and defined the highest good as bounded by a sound condition of the body and the assured hope of its continuance — you have blocked up the approaches of Fortune? How so? For of that very good you can be stripped at any moment. And yet by such talk the inexperienced are caught, and on account of maxims of this kind the crowd that follows these men is large;

\ciceroSection{28} but it is the mark of one who argues acutely to see, not what each man says, but what each man ought to say. As, for instance, in the very position which we have taken up in this discussion, we hold that all good men are always happy. Whom I call good men is plain; for those equipped and adorned with all the virtues we call now wise men, now good men. Let us see who are to be called happy.

\ciceroSection{29} For my part, I judge them to be those who are among goods with no evil joined to them; nor is any other notion attached to this word, when we say "happy," than the full assemblage of goods, with all evils set apart. This virtue cannot attain, if there is any good besides itself. For there will be present a certain throng of evils, if we count those things evils: poverty, obscurity, low standing, loneliness, the loss of one's own, grievous bodily pains, ruined health, infirmity, blindness, the destruction of one's country, exile, slavery at the last. Amid these so many and so great evils — and yet more can befall — the wise man can find himself; for these chance brings on, and chance can run upon the wise man. But if these are evils, who can guarantee that the wise man will always be happy, when he may be in the midst of them all at one and the same time?

\ciceroSection{30} I do not, then, easily grant — neither to my own Brutus, nor to our common masters, nor to those ancients, Aristotle, Speusippus, Xenocrates, Polemo — that, while they count among evils the things I enumerated above, they should yet say that the wise man is always happy. If this distinction delights them, splendid and beautiful and most worthy of Pythagoras, Socrates, and Plato, let them bring their minds to despise those things by whose splendor they are caught — strength, health, beauty, riches, honors, resources — and to count as nothing the things contrary to these: then they will be able to declare with the clearest voice that they are terrified neither by the onset of fortune, nor by the opinion of the crowd, nor by pain, nor by poverty, and that everything for them is lodged within themselves, and that there is nothing outside their own power which they reckon among goods.

\ciceroSection{31} To say these things now — which belong to some great and lofty man — and at the same time to reckon among evils and goods what the crowd reckons there, can in no way be allowed. Stirred by that glory Epicurus comes forward; to him too, if you please, the wise man seems always happy. He is captivated by the dignity of this opinion; but he would never say it, if he listened to himself. For what is there less consistent than that the very man who says that pain is either the highest or the only evil should also hold that the wise man will say, even while he is racked with pain, "How sweet this is!"? Philosophers, then, are not to be judged from single utterances, but from their continuity and consistency.

\ciceroSection{32} You bring me to agree with you. But see that your own consistency, too, is not found wanting. — How so? — Because I lately read your fourth book on the ends of good and evil; in it you seemed to me, in arguing against Cato, to wish to show this — which for my part I do approve — that between Zeno and the Peripatetics there is nothing at issue but a novelty of terms. And if that is so, what reason is there why, if it is consistent with Zeno's account that there should be force enough in virtue for living happily, the same should not be open to the Peripatetics to say? For the matter, I think, is what ought to be looked at, not the words.

\ciceroSection{33} You deal with me by sealed documents, calling to witness what I once said or wrote. Do that with others, who debate under imposed rules; we live from day to day. Whatever has struck our minds with probability, that we say; and so we alone are free. Yet since we spoke a little earlier about consistency, I do not think the question to be asked here is whether it is true — what was Zeno's view, and his pupil Aristo's, that the only good is the honorable — but whether, if it were so, it would then be consistent to rest this whole business of living happily in virtue alone.

\ciceroSection{34} So let us grant Brutus this much, by all means — that the wise man is always happy. How far this is consistent with himself, let him see to it; as for the glory of this opinion, who is worthier of it than that man? Yet let us hold to our point: that he is also supremely happy. And though Zeno of Citium, a kind of newcomer and obscure artificer of words, seems to have insinuated himself into the ancient philosophy, the weight of this opinion may be traced back to the authority of Plato, in whom this manner of speaking is often employed — that nothing is called good except virtue. So in the \textit{Gorgias} Socrates, when he had been asked whether he did not think Archelaus the son of Perdiccas, who was then held the most fortunate of men, a happy man, replied, "I do not know;

\ciceroSection{35} for I have never talked with him." — Is that so? Can you not know it any other way? — In no way. — Then you cannot say even of the great king of the Persians whether he is happy? — How could I, when I do not know how learned he is, how good a man? — What? Do you think the happy life consists in that? — Just so, entirely; I judge the good to be happy, the wicked wretched. — Archelaus, then, is wretched? — Certainly, if he is unjust.

\ciceroSection{36} Does he not seem here to rest the whole happy life in virtue alone? And what of this — how does the same man speak in the funeral oration? "For the man to whom all things that bear on living happily depend on himself, and are not left hanging on the good or ill chance of others, nor forced to drift with another's fortunes — for him the way of living best has been provided. He it is who is moderate, he who is brave, he who is wise; he it is who, when goods come and when they go — children above all, with the rest of his comforts — will yield and obey that old precept: for he will never rejoice too much nor grieve too much, because he sets all his hope of himself always in himself." From Plato, then, as from a kind of holy and august spring, all our discourse will flow. From where, then, can we more rightly begin than from our common parent, nature?

\ciceroSection{37} Nature willed that whatever she has brought forth — not only the animal, but also that which has so sprung from the earth as to be held up by its own roots — should be perfect, each thing in its own kind. And so trees and vines, and the lowlier things that cannot lift themselves higher from the ground — some are forever green, others, stripped bare in winter, put out leaves when the spring has warmed them; and there is nothing that does not so thrive, by a certain inner motion and by the seeds shut up in each, that it pours forth either flowers or fruits or berries; and all things, in all things, so far as lies in them, are perfect when no force stands in the way.

\ciceroSection{38} But the force of nature herself can be discerned even more easily in the beasts, because to them sense has been given by nature. For some beasts nature willed to be swimmers and dwellers in the waters; others, winged, to enjoy the free sky; some to be crawlers, some walkers; and of these themselves, some to range alone, some to gather in herds; some to be savage, others tame; some to be hidden away and covered by the earth. And each of them, keeping to its own office, since it cannot cross over into the life of an unlike creature, abides in the law of nature. And just as to the beasts something distinct was given to each by nature, which each retains as its own and does not depart from, so to man something far more excellent was given — though only those things ought to be called excellent that admit of some comparison. The human mind, plucked from the divine intelligence, can be compared with nothing else but God himself, if it is right to say so.

\ciceroSection{39} If, then, this mind has been cultivated, and if its keen edge has been so tended that it is not blinded by errors, it becomes perfect intelligence — that is, completed reason — which is the same as virtue. And if everything is happy that lacks nothing, and that is filled out and brought to fullness in its own kind, and if this is the property of virtue, then surely all who possess virtue are happy. And here I agree with Brutus — that is, with Aristotle, Xenocrates, Speusippus, Polemo. But to me they seem also supremely happy.

\ciceroSection{40} For what is lacking for living happily to the man who trusts in his own goods? Or how can one who distrusts them be happy? And distrust them he must, who divides goods into three kinds. For how will he be able to trust in the firmness of his body or the stability of fortune? Yet without a good that is stable and fixed and lasting, no one can be happy. What, then, of theirs is of that sort? So that the saying of the Laconian seems to me to fall upon these men — who, to a certain merchant boasting that he had sent out many ships to every coast of the sea, said, "That fortune, hung upon ropes, is not greatly to be desired." Is there any doubt that nothing is to be counted among the things by which the happy life is filled out, if it can be lost? For nothing of those things in which the happy life consists ought to wither, nothing to be quenched, nothing to fall away. For whoever fears to lose any of them cannot be happy.

\ciceroSection{41} For we want the man who is happy to be safe, unassailable, fenced about and fortified — not so as to be furnished with some small fear, but with none at all. For just as a man is called blameless not who does slight harm, but who does no harm, so he is to be held free from fear not who fears small things, but who is altogether empty of fear. For what else is courage but a disposition of the mind that endures in facing danger and in toil and pain, and at the same time stands far off from all fear? And these things surely would not be so, unless every good consisted in the honorable alone.

\ciceroSection{42} And how can anyone possess that most longed-for and sought-after security — by security I now mean freedom from distress, in which the happy life is set — when there is, or can be, a multitude of evils close upon him? How can anyone be lofty and upright, counting all that can befall a man as small — such as we want the wise man to be — unless he reckons that all his goods rest within himself? When Philip threatened the Spartans by letter that he would block everything they attempted, did they not ask whether he would block them even from dying? Will not the man we are seeking be found far more readily with such a spirit than a whole state was found? What more? When to this courage of which we speak temperance is joined — the moderator of all commotions — what can be lacking for living happily to the man whom courage delivers from distress and from fear, while temperance both calls him away from lust and does not let him exult in any insolent excitement? That virtue brings these things about I would show, were they not set out in the earlier days.

\ciceroSection{43} And since disturbances of the mind make for misery, while their calming makes life happy, and the account of disturbance is twofold — for distress and fear turn upon imagined evils, while exultant gladness and lust turn upon a mistaken view of goods, all of them at war with counsel and reason — when you have seen a man empty of these so weighty agitations, so much at variance and at odds among themselves, and so set loose and free, will you hesitate to call him happy? But the wise man is always so disposed; the wise man, therefore, is always happy. And further: every good is a thing to rejoice in; and what is to be rejoiced in is to be proclaimed and carried before one; and what is of that sort is also a thing of glory; and if of glory, then surely praiseworthy; and what is praiseworthy is assuredly also honorable.

\ciceroSection{44} What is good, therefore, is honorable. But the things these men reckon among goods, not even they themselves call honorable; the only good, then, is the honorable. From which it follows that the happy life is contained in the honorable alone. Those things, then, are not to be called goods, nor to be held such, in which a man may abound and yet be utterly wretched.

\ciceroSection{45} Or do you doubt that a man outstanding in health, strength, beauty, with the keenest and soundest senses — add even, if you like, agility and speed, throw in riches, honors, commands, power, glory — if the man who has these be unjust, intemperate, cowardly, of a dull and worthless wit, you would not hesitate to call him wretched? Of what sort, then, are those goods which the man who has them can be utterly wretched? Let us consider whether, just as a heap is made of grains of its own kind, so the happy life ought to be made of parts that resemble itself. And if that is so, it must be made of goods that are honorable alone; if these shall be mixed of unlike things, nothing honorable can be made of them; and once that is taken away, what can be understood as happy? For whatever is good is a thing to be desired; and what is to be desired is surely to be approved; and what you have approved is to be held welcome and acceptable; and so worth must be assigned to it as well. And if that is so, it must of necessity be praiseworthy; every good, then, is praiseworthy. From which it follows that what is honorable is the only good.

\ciceroSection{46} For unless we hold to this firmly, there will be many things that we shall have to call goods. I leave aside riches — which, since anyone, however unworthy, may possess them, I do not count among goods, for what is a good cannot be possessed by just anyone; I leave aside nobility and a fame stirred up by the agreement of fools and rogues: yet these, though they are the least of things, must still be called goods — bright little teeth, charming eyes, a pleasing complexion, and the things that Anticlea praises while washing the feet of Ulysses: \textit{the smoothness of his speech, the softness of his body.} If we shall reckon these as goods, what will there be in a philosopher's gravity that can be called either weightier or grander than what is in the opinion of the crowd and the throng of fools?

\ciceroSection{47} But, you say, the Stoics call these same things "preferred" or "advanced," which your school calls goods. They do call them so, to be sure; but they deny that the happy life is made complete by them — whereas these others think there is no happy life without them, or, if there is, certainly deny it is the happiest. We, however, will have it the happiest, and this is confirmed for us by that Socratic chain of reasoning. For thus did that prince of philosophy argue: as the disposition of each man's soul is, such is the man; and as the man himself is, such is his speech; and his deeds are like his speech, and his life like his deeds. But the disposition of a good man's soul is praiseworthy; and so the life of a good man is praiseworthy; and therefore honorable, since praiseworthy. From which it is concluded that the life of good men is happy.

\ciceroSection{48} For, by the faith of gods and men! — was it too little established in our earlier discussions, or did we speak merely for delight and to pass the leisure hours, that the wise man is always free from every commotion of soul, which I call a disturbance, and that there is always in his soul the most placid peace? A man, then, who is temperate, steadfast, without fear, without distress, without idle elation, without lust — is he not happy? But the wise man is always such; therefore he is always happy. And further, how can a good man fail to refer to what is praiseworthy everything that he does and everything that he feels? But he refers everything to living happily; therefore the happy life is praiseworthy; and nothing is praiseworthy without virtue: therefore the happy life is brought about by virtue.

\ciceroSection{49} And this too is concluded in the following way: there is nothing to be proclaimed or gloried in either in a wretched life or in one that is neither wretched nor happy. And yet in some life there is something to be proclaimed and gloried in and held up before the world, as in Epaminondas: \textit{by our counsels the glory of the Spartans was shorn;} or in Africanus: \textit{from the rising of the sun above the marshes of Maeotis, there is no man who can match my deeds.}

\ciceroSection{50} And if this is so, then the happy life is to be gloried in and proclaimed and held up before the world; for there is nothing else that ought to be proclaimed and held up before the world. With these things laid down, you understand what follows. And indeed, unless that life is happy which is also honorable, it follows of necessity that there is something better than the happy life; for whatever shall be honorable, they will certainly admit to be better. So the happy life will be something less than another thing; and what can be said more perverse than that? What of this? When they admit that there is force enough in vices to make a life wretched, must it not be admitted that there is the same force in virtue to make it happy? For of contraries the consequences are contrary.

\ciceroSection{51} At this point I ask what force there is in that balance of Critolaus, who, when he places the goods of the soul in one pan and the goods of the body and external goods in the other, thinks the pan of the soul's goods weighs down so heavily that it would sink the earth and the seas. What, then, prevents either this man, or that gravest of philosophers Xenocrates, who exalts virtue so greatly while making light of everything else and casting it away, from placing in virtue not merely a happy life but the happiest of lives?

\ciceroSection{52} And unless this is so, the destruction of the virtues will follow. For on whomever distress falls, fear must fall on the same man as well — for fear is the anxious expectation of distress to come; and on whomever fear falls, on the same man fall dread, timidity, panic, cowardice; and so the same man will at times be conquered, and will not think that precept of Atreus applies to him: \textit{let them so prepare themselves in life that they know not how to be conquered.} But this man will be conquered, as I have said, and not only conquered but enslaved as well; while we will have virtue always free, always unconquered; and unless these things hold, virtue is overthrown.

\ciceroSection{53} And if there is enough protection in virtue for living well, there is enough also for living happily; for there is certainly enough in virtue for us to live bravely; if bravely, then also with greatness of soul, and indeed so that we are never terrified by anything and are always unconquered. It follows that nothing is regretted, nothing lacking, nothing in the way; therefore all things flow on freely, perfectly, prosperously, and thus happily. But virtue has power enough for living bravely; therefore power enough also for living happily.

\ciceroSection{54} For just as folly, even when it has attained what it desired, never thinks it has gained enough, so wisdom is always content with what is at hand, and is never displeased with itself. Do you suppose the single consulship of Gaius Laelius was like another's — and that too coming with a defeat (if, when a man wise and good, such as he was, is passed over in the voting, it is not rather the people that has met a defeat at the hands of a good consul than he at the hands of a good people)? But still — which would you rather, if the choice were yours: to be consul once like Laelius, or four times like Cinna?

\ciceroSection{55} I have no doubt what you will answer; and so I see to whom I am committing the question. I would not put this same question to just anyone; for another might perhaps answer that he prefers not only four consulships to one, but a single day of Cinna's to whole lifetimes of many famous men. Laelius, had he so much as touched a man with his finger, would have paid the penalty; but Cinna ordered the head of his colleague, the consul Gnaeus Octavius, to be struck off, and the heads of Publius Crassus and Lucius Caesar, men of the highest nobility, whose virtue had been proved at home and in war, and of Marcus Antonius, the most eloquent of all men I ever heard, and of Gaius Caesar, in whom there seemed to me to be a model of refinement, of wit, of charm, of grace. Was he happy, then, who killed these men? To me, on the contrary, he seems wretched not only because he did these things, but also because he so conducted himself that he was free to do them — though in truth no one is free to do wrong; but we slip through an error of language, for we say a thing is permitted

\ciceroSection{56} when it is granted to a man to do it. Which, after all, was the happier — Gaius Marius, when he shared the glory of the Cimbric victory with his colleague Catulus, almost a second Laelius (for I count this man very like him), or when, victor in the civil war and enraged, he answered the kinsmen of Catulus who pleaded with him not once but again and again: "let him die"? In this the happier was the man who obeyed that wicked command than the one who issued so criminal an order. For as it is better to receive an injury than to do one, so it was better to go a little forward to meet death already drawing near — which is what Catulus did — than to do as Marius did: to crush, by the destruction of such a man, his own six consulships, and to defile the last span of his life.

\ciceroSection{57} For thirty-eight years Dionysius was tyrant of the Syracusans, having seized the lordship at the age of twenty-five. With what splendor, with what resources he held the city, and the state oppressed in slavery! And yet of this man we have received from good authorities this account: that he was in his manner of living of the utmost temperance, in the conduct of affairs a keen and industrious man, but at the same time malicious by nature and unjust; from which it must seem to all who rightly look upon the truth that he was the most wretched of men. For the very things he had set his heart on, not even then, when he believed he could do all things, did he attain.

\ciceroSection{58} For though he was born of good parents and of honorable station — though indeed this one point is handed down differently by different writers — and abounded both in the friendships of his peers and in the company of his kinsmen, and had besides, after the Greek fashion, certain young men bound to him by love, he trusted none of them, but committed the guarding of his person to men whom he had chosen as slaves out of the households of the rich, from whom he himself had stripped the name of slavery, and to certain immigrants and savage barbarians. Thus, through his unjust craving for lordship, he had in a manner shut himself up in prison. Indeed, that he might not entrust his throat to a barber, he taught his own daughters to cut hair. So in a sordid and servile craft these royal maidens, like little barber-girls, trimmed the beard and hair of their father. And yet from these very girls, when they were now grown, he took away the blade, and arranged that with glowing walnut shells they should singe off his beard and hair.

\ciceroSection{59} And though he had two wives, Aristomache, his fellow citizen, and Doris of Locri, he would come to them by night only after spying out and searching through everything beforehand. And since he had set a broad ditch around the chamber where his bed stood, and had joined the crossing of that ditch by a little wooden bridge, he himself would swing this very bridge aside once he had shut the door of the chamber. And since he did not dare to stand on the common platforms, he was accustomed to address the assembly from a high tower.

\ciceroSection{60} And when he wished to play ball — for he did this eagerly and often — and was laying aside his tunic, he is said to have handed his sword to a young man whom he loved. At this one of his companions said in jest: "To him, at any rate, you certainly entrust your life," and the young man smiled; whereupon he ordered both to be put to death, the one because he had pointed out the way of killing him, the other because he had approved the remark by his laughter. And by that deed he was so grieved that he bore nothing more heavily in all his life; for he had killed one he had loved intensely. So the desires of men without self-mastery are pulled apart in opposite directions: when you have indulged the one, the other must be resisted.

\ciceroSection{61} And yet this very tyrant passed judgment on his own happiness. For when one of his flatterers, Damocles, was recounting in conversation his forces, his power, the majesty of his rule, the abundance of his wealth, the splendor of his royal palace, and declaring that no one had ever been more fortunate, "Would you like, then, Damocles," he said, "since this life delights you, to taste it for yourself and try out my fortune?" When the man said he would gladly, Dionysius ordered him laid on a golden couch spread with a covering of the most beautiful weave, embroidered with magnificent designs, and had several sideboards decked with chased gold and silver. Then he ordered chosen boys of surpassing beauty to take their stand at the table and to wait upon him attentively, watching his every nod.

\ciceroSection{62} There were perfumes and garlands; incense was burning; the tables were piled high with the choicest delicacies. Damocles thought himself a fortunate man. In the midst of all this display Dionysius ordered a gleaming sword, fastened by a single horsehair, to be let down from the ceiling, so that it hung over the neck of this happy man. And so Damocles neither looked at those beautiful attendants, nor at the silver wrought with such art, nor reached out his hand to the table; the garlands by now were slipping from his head; at last he begged the tyrant to let him go, since he no longer wished to be happy. Does Dionysius not seem to have shown clearly enough that there is no happiness for a man over whom some terror forever hangs? Nor was it any longer open to him to return to justice, to give his citizens back their freedom and their rights; for in his youth, at an age that takes no thought, he had so entangled himself in errors and committed such crimes that he could not be safe if he once began to be sane.

\ciceroSection{63} How much he longed for friendships, whose faithlessness he dreaded, he showed in the case of those two Pythagoreans: when he had accepted one of them as surety for the other's death, and the other, to free his surety, had presented himself at the appointed hour of death, "Would," he said, "that I might be enrolled as a third friend among you!" How wretched it was for him to do without the company of friends, the fellowship of daily life, intimate conversation altogether — and this for a man learned from boyhood and schooled in the liberal arts, devoted besides to music; a writer of tragedy, too — how good a one is nothing to the point, for in this field, somehow more than in others, every man finds his own work beautiful. I have never yet known a poet (and I had a friendship with Aquinius) who did not think himself the best; that is how it stands: your work delights you, mine delights me. — But to return to Dionysius: he was deprived of all human refinement and society; he lived among runaways, among criminals, among barbarians; he counted no one his friend who was either worthy of freedom or wished to be free at all. With this man's life — than which I can conceive nothing more foul, more wretched, more detestable — I shall not now compare the life of Plato or Archytas, learned men and plainly wise:

\ciceroSection{64} from the same city I shall summon up a lowly little man from his dust and drawing-board, one who lived many years later, Archimedes. When I was quaestor I tracked down his grave, unknown to the Syracusans — they denied it existed at all — hemmed in on every side and overgrown with brambles and thickets. For I held in mind certain little verses that I had been told were inscribed upon his monument, which declared that a sphere with a cylinder had been set on top of the tomb.

\ciceroSection{65} Now as I was surveying everything with my eyes — for there is a great crowd of tombs at the Agrigentine Gate — I noticed a small column standing out a little from the thickets, on which there was the figure of a sphere and a cylinder. And at once I told the Syracusans — their leading men were with me — that I thought this was the very thing I was looking for. Men were sent in with sickles; they cleared and opened up the place.

\ciceroSection{66} When the way to it had been laid open, we approached the base facing us. There appeared the epigram, with the latter halves of the little verses eaten away — about half of it remaining. So the most renowned city of Greece, once indeed the most learned as well, would have known nothing of the monument of its one most brilliant citizen, had it not learned of it from a man of Arpinum. But let the discourse return to where it strayed: who is there of all men, provided he has any dealings at all with the Muses — that is, with culture and learning — who would not rather be this mathematician than that tyrant? If we ask after the manner and conduct of their lives, the one man's mind was nourished by the working out and investigation of reasonings, with the delight of his own ingenuity, which is the single sweetest food of the soul; the other's, by slaughter and injustice, with fear by day and by night. Come now, set beside him Democritus, Pythagoras, Anaxagoras: what kingdoms, what riches will you prefer to their pursuits and delights?

\ciceroSection{67} For that part which is best in a man — there it is that the best thing you seek must necessarily reside. And what is there in a man better than a keen and sound mind? It is the good of this mind, then, that we must enjoy, if we wish to be happy; but the good of the mind is virtue; therefore the happy life must necessarily be contained in virtue. From this come all things that are beautiful, honorable, and noble — as I said above, but the same point seems worth stating a little more fully — and they are full of joys. And since it is clear that the happy life arises from perpetual and abundant joys, it follows that it arises from the honorable.

\ciceroSection{68} But that we may not merely touch in words upon what we wish to demonstrate, certain considerations must be set out, so to speak, to move us, which may turn us more readily toward recognition and understanding. Let us, then, take a certain outstanding man, accomplished in the finest arts, and let us fashion him for a moment in our mind and thought. First, he must be of exceptional natural gift; for virtue does not readily keep company with sluggish minds. Next, he must have an eager zeal for tracking down the truth. From this there will arise that threefold yield of the mind: one part set in the knowledge of things and the unfolding of nature, a second in the marking out of what is to be sought and what shunned and in the principle of living well, a third in judging what is consequent upon each thing and what contradicts it — and in this lies all the subtlety of reasoning together with the truth of judgment.

\ciceroSection{69} With what joy, then, must the mind of the wise man be filled, dwelling and passing its nights among these concerns! When he has discerned the motions and revolutions of the whole universe, and seen the countless stars fixed in the sky moving in concert with its own motion, set in their appointed seats; and the seven others each holding their own courses, far apart from one another in their height or their lowness, whose wandering motions yet mark out the fixed and certain spans of their orbits — surely the sight of these things urged on those men of old and prompted them to seek further. From this was born the search into beginnings, and into seeds, as it were, from which all things rose, were generated, and took shape; and what is the origin of each kind, whether lifeless or living, whether mute or speaking; what its life, what its death, and what the succession and change from one thing into another; whence the earth, and by what weights it is balanced, in what hollows the seas are held, by what force of gravity all things are carried so as ever to seek the middle place of the world, which in a round body is also the lowest.

\ciceroSection{70} For the mind that handles these matters and reflects on them night and day, there arises that self-knowledge enjoined by the god at Delphi: that the mind should recognize itself and feel itself joined with the divine mind, from which it is filled with an insatiable joy. For the very contemplation of the power and nature of the gods kindles a zeal to imitate their eternity, and the mind does not think itself confined within the brevity of life, when it sees the causes of things linked one to another and bound by necessity — causes which, though they flow from everlasting time into everlasting, are yet governed by reason and mind.

\ciceroSection{71} Gazing upon these things and looking up at them — or rather, surveying every part and region of them — with what tranquility of mind, in turn, does he consider human affairs and things nearer at hand! From this comes that knowledge of virtue; the kinds and parts of the virtues come into flower; it is discovered what nature looks to as the ultimate end among goods and the last among evils, to what the duties are to be referred, what plan of living out one's life is to be chosen. And from the investigation of these and like matters there comes about, above all else, that very thing which is the business of this discussion: that virtue is self-sufficient for living happily.

\ciceroSection{72} There follows the third part, which runs and pours through every part of wisdom — which defines a thing, distinguishes its kinds, joins what follows, draws conclusions to their completion, and judges true from false: the method and science of reasoning. From it arises both the highest usefulness for weighing matters and, above all, a delight that is liberal and worthy of wisdom. But these are the works of leisure. Let this same wise man pass over to the defense of the commonwealth. What could be more excellent than he, when in his prudence he discerns what is to the advantage of the citizens, in his justice diverts nothing of it to his own house, and makes use of all his other so many and so various virtues? Add the harvest of friendships, in which the learned find both counsel for the whole of life, harmonious and almost of one breath, and the highest pleasure in their daily intercourse and shared life. What, in the end, does this life want, to be happier? Fortune herself must yield to a life crammed with so many and so great joys. But if to rejoice in such goods of the mind — that is, in the virtues — is to be happy, and if all the wise enjoy these joys to the full, then it must be admitted that they are all happy. — Even on the rack and under torture?

\ciceroSection{73} Or did you suppose I was speaking of him as lying on violets and roses? Or shall it be permitted to Epicurus — who merely put on the mask of a philosopher and inscribed the name upon himself — to say (a thing that, as the matter stands, he says to my applause) that there is no time for the wise man, even if he be burned, racked, cut, when he cannot cry out, "How utterly I scorn it!" — and this above all when he defines every evil by pain and every good by pleasure, mocks these honorable and base things of ours, and says that we are taken up with words and pour out empty sounds, that nothing concerns us except what is felt in the body as smooth or rough? To this man, then, differing little, as I said, from the judgment of beasts, it shall be permitted to forget himself, and at one moment to despise fortune — though his every good and evil lies in fortune's power — and at the next to call himself happy in the height of torment and torture, though he has laid it down that pain is not only the highest evil but the only one?

\ciceroSection{74} Nor indeed has he furnished himself with those remedies for enduring pain — firmness of mind, shame at what is base, the practice and habit of suffering, the precepts of fortitude, a manly hardness — but he says he finds rest in a single thing, the recollection of past pleasures; as if a man scorched by heat, when he cannot easily bear the violence of it, should wish to recall that once, in our own country of Arpinum, he was plunged about by cool rivers. For I do not see how past pleasures can soothe

\ciceroSection{75} present evils. — But since the man who says this declares the wise man always happy — a thing he has no right to say, if he wished to be consistent with himself — what then must they do who hold that nothing is to be sought, nothing to be counted among goods, that lacks honor? On my authority, then, let even the Peripatetics and the old Academics at last leave off their stammering and dare to say openly and in a clear voice that the happy life will go down into the bull of Phalaris.

\ciceroSection{76} For let there be three kinds of goods — to withdraw at last from the Stoic snares, which I see I have employed more than is my habit — let there be those kinds of goods, then, by all means, provided the goods of the body and the external ones lie low upon the ground and are called goods only because they are to be taken up, while those divine goods of the soul spread themselves far and wide and reach the very sky; so that the man who has attained them — why should I call him merely happy and not the happiest of all? But will the wise man dread pain? For it is pain above all that fights against this verdict. Against the death of ourselves and of those dear to us, against distress and the other disturbances of the soul, we seem armed and ready enough by the discussions of the earlier days; pain seems the fiercest adversary of virtue. It brandishes its blazing torches; it threatens to wear down courage, greatness of soul, and endurance.

\ciceroSection{77} Will virtue, then, give way to this? Will the happy life of the wise and steadfast man yield to it? How shameful, good gods! Spartan boys do not so much as groan when they are torn by the pain of the lash. We ourselves saw at Lacedaemon bands of young men contending with unbelievable fierceness — with fists, heels, nails, even teeth at the last — sooner letting themselves be beaten senseless than confess themselves beaten. What barbarian land is wilder or rougher than India? Yet even among that people, those first who are held to be wise pass their lives naked, and endure the snows of the Caucasus and the violence of winter without pain; and when they have laid themselves against the flame, they are scorched without a groan.

\ciceroSection{78} The women of India, moreover, when one of them has lost her husband, come into a contest and a trial: which of them he loved most — for several wives are usually married to a single man. She who is the victor, attended joyfully by her own people, is laid upon the pyre together with her husband; the one defeated departs in grief. Custom would never overcome nature, for nature is forever unconquered; but we have infected our souls with shadows, with luxuries, with idleness, sloth, and indolence; we have softened them, lulled by opinions and by bad habit. Who does not know the custom of the Egyptians? Their minds, steeped in the errors of perversity, would undergo any torture sooner than do violence to an ibis, an asp, a cat, a dog, or a crocodile; and even if they have done one of these creatures some harm unawares, they refuse no punishment for it.

\ciceroSection{79} I speak of men. What of beasts? Do they not endure cold, hunger, their roving courses and wanderings over mountains and through forests? Do they not fight for their young so fiercely that they take wounds upon themselves, flinching from no assault, no blow? I pass over what men endure and suffer who are greedy for office for ambition's sake, eager for praise for the sake of glory, kindled by love for the sake of desire. Life is full of such examples.

\ciceroSection{80} But let our discourse keep its measure and return to the point from which it turned aside. The happy life, I say, will give itself over to the torments, and having followed justice, temperance, and above all courage, greatness of soul, and endurance, when it has seen the face of the torturer it will halt — and with all the virtues setting out toward the torment without any terror of soul, it will stand its ground outside the doors, as I said before, and on the threshold of the prison. For what could be more loathsome than the happy life, what more disfigured, left alone, cut off from its fairest company? Yet that can in no way come to pass; for the virtues cannot hold together without the happy life, nor the happy life without the virtues.

\ciceroSection{81} And so they will not allow it to turn its back, but will carry it off with them to whatever pain and torment they themselves are led. For it is the wise man's own mark to do nothing he could ever repent of, nothing against his will, but all things splendidly, steadfastly, with weight and with honor; to expect nothing as though it were certain to come, to wonder at nothing when it has come, as though it had befallen unforeseen and strange; to refer all things to his own judgment, to stand by his own decisions. What could be happier than this, I for one cannot conceive.

\ciceroSection{82} The Stoics' conclusion, indeed, is easy. Since they have judged that the end of goods is to be in harmony with nature and to live in accord with it, and since this lies within the wise man's power not by duty alone but by capacity, it must follow that the man in whose power the highest good lies has the happy life in his power as well. Thus the wise man's life is always happy. — There you have what I think can be said most stoutly about the happy life, and, as matters now stand, unless you bring something better, most truly as well. — For my part I can bring nothing better; but I would gladly win this from you, if it is no trouble: since no bonds tie you to any one fixed school, and you draw from all of them whatever most moves you with the look of truth — what you seemed a little while ago to be urging upon the Peripatetics and the Old Academy, that without any holding back they should dare to say freely that the wise are always happiest, that is what I should like to hear, how you think it consistent for them to say it. For you said much against that verdict, and much that was deduced by the reasoning of the Stoics.

\ciceroSection{83} Let us use, then, the liberty that we alone in philosophy are allowed to use — we whose discourse pronounces no judgment of its own, but ranges over every side, so that it may be judged by others through itself, with no man's authority attached. And since you seem to wish this — that whatever the verdict of the disagreeing philosophers may be about the ends, virtue should nevertheless have safeguard enough for the happy life — which is what we are told Carneades used to argue; but he argued it against the Stoics, whom he was always most eager to refute, and against whose doctrine his genius had taken fire. I, however, will conduct it in peace — for if the Stoics have set down the end of goods rightly, the matter is settled: the wise man must always be happy —

\ciceroSection{84} but let us examine each of the remaining verdicts one by one, if it can be done, to see whether this splendid decree, so to speak, of the happy life can be made to agree with the verdicts and doctrines of all. Now these, as I suppose, are the verdicts about the ends that have been held and defended. First, four simple ones: nothing is good but the honorable, as the Stoics hold; nothing is good but pleasure, as Epicurus holds; nothing is good but freedom from pain, as Hieronymus holds; nothing is good but to enjoy the first goods of nature, either all of them or the greatest, as Carneades argued against the Stoics.

\ciceroSection{85} These, then, are the simple verdicts; the following are the mixed: three kinds of goods, the greatest of the soul, the second of the body, the third the external ones, as the Peripatetics hold, and not much otherwise the Old Academics. Pleasure with honor Dinomachus and Callipho coupled together; freedom from pain the Peripatetic Diodorus joined to the honorable. These are the verdicts that have some stability; for those of Aristo, of Pyrrho, of Erillus, and of certain others have faded away. What these can maintain, let us see — leaving the Stoics aside, whose verdict I think I have defended sufficiently. And the case of the Peripatetics has indeed been set out, except for Theophrastus and any who, following him, shrink from pain too feebly and dread it. The rest may do what they generally do — exalt the weight and worth of virtue. And when they have raised it to the sky, as eloquent men are wont to do with full abundance, it is easy to grind down and despise the rest by comparison. For it cannot be permitted to those who say that praise is to be sought even at the cost of pain to deny that those are happy who have attained it. For although they may be amid certain evils, still this name of "the happy" extends far and wide.

\ciceroSection{86} For just as trade is called profitable and plowing fruitful — not if the one is forever free of every loss and the other of every disaster of weather, but if good fortune stands out in much the greater part of each — so a life can rightly be called happy not only if it is crammed with goods on every side, but if the good things weigh down the scale by much the greater and heavier part.

\ciceroSection{87} By their reckoning, then, the happy life will follow virtue even to torture, and will go down with it into the bull, on the authority of Aristotle, Xenocrates, Speusippus, and Polemo, and corrupted by no threats or blandishments will not desert it. The verdict of Callipho and of Diodorus will be the same, since each of them embraces the honorable so as to judge that all things which exist apart from it are to be set far behind it. The rest seem to hold a narrower ground, yet they swim clear all the same — Epicurus, Hieronymus, and any there are who care to defend that deserted Carneadean end; for there is not one of them who does not hold the soul to be the judge of goods, and who does not train it so that it can despise whatever seems good or evil.

\ciceroSection{88} For the case that seems to you Epicurus's will be the same as Hieronymus's and Carneades's, and, by Hercules, that of all the rest. For who is too little prepared against death or pain? Let us begin, if you please, with the man whom we call soft, whom we call a pleasure-seeker. What? Does he seem to you to fear death or pain — he who calls the very day on which he dies a happy one, and who, struck by the greatest pains, confounds those very pains by the memory and recollection of his own discoveries? And he does not handle this in such a way as to seem to be merely babbling off the cuff. For about death he holds this view: that when the living being is dissolved, sensation is extinguished, and that what is without sensation he judges to be nothing to us. As for pain, likewise, he has fixed principles to follow: he consoles its severity by its brevity, and its long duration by its lightness.

\ciceroSection{89} How, after all, do those grandiloquent ones stand against these two things, which most torment us, better than Epicurus? And as for the other things that are reckoned evils — do not Epicurus and the rest of the philosophers seem prepared enough against them too? Who does not dread poverty? Yet not one of the philosophers does. And this very man himself — with how little is he content! No one has said more about a thin diet. For the things that bring on a craving for money — that there be means enough to supply love, ambition, and daily expenses — since he stands far off from all of these, why should he greatly long for money, or rather, why should he care for it at all?

\ciceroSection{90} Could the Scythian Anacharsis count money as nothing, and our own countrymen, the philosophers, not be able to do it? A letter of his is handed down in these words: \textit{Anacharsis to Hanno, greetings. My clothing is a Scythian hide, the soles of my feet their own hardened skin, my bed the earth, my relish hunger; I live on milk, cheese, and meat. So you may come to me as to a man at peace. But those gifts of yours, in which you took such delight, give them to your fellow citizens or to the immortal gods.} Nearly all the philosophers of every school — except those whom a flawed nature had twisted away from right reason — could be of this same mind.

\ciceroSection{91} Socrates, when a great mass of gold and silver was being carried past in a procession, said: "How many things there are that I do not want!" Xenocrates, when envoys from Alexander had brought him fifty talents — a sum that in those days, and at Athens above all, was very large — took the envoys off to dinner in the Academy; he set before them just so much as was enough, with no display. When on the next day they asked him to whom he wished the money paid out, he said: "What? Did you not understand from yesterday's modest little supper that I have no need of money?" And seeing that they looked rather downcast, he accepted thirty minas, so as not to seem to scorn the king's generosity.

\ciceroSection{92} But Diogenes spoke more freely, as a Cynic should: when Alexander asked him to say if he wanted anything, "Just now," he said, "move a little out of my sun." Alexander had, it seems, blocked the light as he was warming himself. And this man used to argue how far he surpassed the king of the Persians in his life and his fortune: that he himself lacked nothing, while for the king nothing would ever be enough; that he did not crave the king's pleasures, with which the king could never be satisfied, whereas the king could in no way attain his own. You see, I imagine, how Epicurus divided the kinds of cravings — not too subtly, perhaps, but usefully all the same:

\ciceroSection{93} some are natural and necessary, some natural and not necessary, some neither; the necessary ones can be satisfied at almost no cost — for nature's riches are easily come by; the second kind of cravings, he holds, are neither hard to gratify nor hard to go without; the third kind, since they were plainly empty and touched neither necessity nor even nature, he thought should be thrown out root and branch.

\ciceroSection{94} On this point the Epicureans say a great deal, and they belittle one by one those pleasures whose kinds they do not despise, yet whose abundance they do not seek. For even the pleasures of the flesh, on which they hold forth at length, they call easy, common, set out for all in plain reach; and they hold that, if nature requires them, they are to be measured not by family or place or rank, but by good looks, youth, and figure; that to abstain from them is by no means difficult, if health or duty or reputation demands it; and that this whole class of pleasures is, in short, desirable when it does no harm, but never beneficial.

\ciceroSection{95} And on this whole matter of pleasure his teaching is this: that pleasure itself, in and of itself, because it is pleasure, is always to be desired and sought after, and that by the same reasoning pain, for that very reason, because it is pain, is always to be fled; and so the wise man will employ a kind of balancing, fleeing a pleasure if it is going to produce a greater pain, and taking on a pain if it produces a greater pleasure; and that all delights, though they are judged by the senses of the body, are nonetheless referred back to the mind.

\ciceroSection{96} For this reason the body rejoices only so long as it feels the pleasure present to it, while the mind both perceives the present pleasure together with the body and looks ahead to the one that is coming, and does not let the past one slip away. Thus the wise man will always have unbroken, interwoven pleasures, since the expectation of pleasures hoped for is joined to the memory of pleasures already enjoyed.

\ciceroSection{97} And points like these are carried over to the matter of diet as well, and the magnificence and extravagance of banquets is belittled, on the ground that nature is content with little upkeep. For who does not see that all such things are seasoned by appetite? Darius, in his flight, when he had drunk water that was muddy and fouled with corpses, declared that he had never drunk anything more pleasant. He had simply never drunk while thirsty before. Nor had Ptolemy ever eaten while hungry: when, as he was traveling across Egypt and his companions had not kept up with him, he was given coarse bread in a hut, nothing seemed to him more pleasant than that bread. They say that Socrates, when he had walked rather briskly until evening and was asked why he was doing it, replied that he was foraging up an appetite by walking, so as to dine the better.

\ciceroSection{98} And again — do we not see the diet of the Spartans at their common messes? There the tyrant Dionysius, when he had dined, declared that he took no delight in that black broth which was the centerpiece of the meal. Then the man who had cooked it said: "No wonder; the seasonings were missing." "What seasonings, pray?" he asked. "Toil in the hunt, sweat, a run to the Eurotas, hunger, thirst. For it is by these that the meals of the Spartans are seasoned." And this can be understood not from the way of men alone, but from the beasts as well, which, whenever something is set before them that is not foreign to their nature, are content with it and seek nothing more.

\ciceroSection{99} Whole communities, schooled by custom, take delight in thrift, as we said a little while ago of the Spartans. The diet of the Persians is set out by Xenophon, who says that to their bread they add nothing but cress. And yet, if nature should desire certain things more agreeable as well, how many such are produced from the earth and the trees, both in easy abundance and of surpassing sweetness! Add the leanness that follows upon this restraint in diet, add the soundness of health; then compare the men who sweat and belch, stuffed with banquets like fattened cattle:

\ciceroSection{100} you will understand then that those who pursue pleasure most attain it least, and that the pleasantness of a diet lies in appetite, not in being sated. They say that Timotheus, a famous man at Athens and a leader of the state, when he had dined at Plato's and had been quite delighted with the meal, said on seeing him the next day: "Your dinners, to be sure, are pleasant not only at the moment but the day after as well." And what of the fact that we cannot even use our minds properly when we are filled with too much food and drink? There is a celebrated letter of Plato's to the kinsmen of Dion, in which it is written more or less in these words: "When I had come there, that life called happy, full of Italian and Syracusan tables, in no way pleased me — to be stuffed twice a day and never to spend a night alone, and all the rest that accompanies a life in which no one will ever be made wise, and a man of moderation far less. For what nature can be so marvelously tempered?"

\ciceroSection{101} How, then, can a life be pleasant from which good sense is absent, from which moderation is absent? From this the folly of Sardanapalus, the richest king of Syria, is laid bare, who ordered to be cut on his tomb: \textit{What I ate I have, and what my glutted lust drank up; but those many splendid things lie left behind.} "What else," says Aristotle, "would you inscribe on the tomb of an ox rather than of a king?" He says that, dead, he has those things which, even alive, he had no longer than he was enjoying them.

\ciceroSection{102} Why, then, should riches be longed for, or where does poverty not permit men to be happy? You are keen, I suppose, on statues and paintings. If there is anyone who delights in these, do not men of slender means enjoy them better than those who have them in abundance? For of all such things there is in our city the greatest supply in public places. Those who own them privately neither see so many of them, and see them rarely, only when they have come to their country estates; and even so something pricks them, when they recall from where they have them. The day would fail me, were I to undertake to plead the case of poverty. The matter is plain, and nature herself reminds us daily how few, how small, how cheap are the things she needs. Will obscurity, then, or low birth, or even the disfavor of the crowd, keep the wise man from being happy?

\ciceroSection{103} Consider whether favor with the crowd and this glory that men strive for do not bring more vexation than pleasure. Our Demosthenes was decidedly vain, who said that he was delighted by the whisper of a little woman carrying water — as the custom is in Greece — murmuring to her neighbor: "That is the famous Demosthenes." What could be more trifling than this? And yet what an orator! But he had clearly learned to speak before others, not much with himself.

\ciceroSection{104} We must understand, then, that popular glory is not in itself to be sought, nor obscurity to be dreaded. "I came to Athens," says Democritus, "and no one there recognized me" — a steady and serious man, who could glory in having been free of glory! Or shall pipers and those who play the lyre tune their melodies and rhythms by their own judgment, not the multitude's, while the wise man, endowed with a far greater art, will seek out not what is truest but what the crowd wants? Or is there anything more foolish than to suppose that men who, taken one by one, you would despise as laborers and barbarians, amount to something taken all together? He, in truth, will despise our ambitions and our frivolities, and will reject the honors of the people even when they are offered unasked; but we do not know how to despise them until we have begun to regret them. There is a passage in Heraclitus the natural philosopher about Hermodorus, the foremost man of the Ephesians;

\ciceroSection{105} he says that all the Ephesians deserve to be put to death, because, when they drove Hermodorus from the city, they spoke thus: "Let no single one of us stand out; if anyone does, let him be in another place and among other men." Does this not happen so among every people? Do they not hate all preeminence in virtue? What of Aristides — for I prefer to bring forward examples from the Greeks rather than our own — was he not banished from his country for this very reason, that he was just beyond measure? How free, then, from vexations are those who have no dealings at all with the people! For what is sweeter than a leisure given to letters? — letters, I mean, by which we come to know the boundlessness of things and of nature, and, within this very world of ours, the heaven, the lands, the seas. With honor scorned, then, and money scorned as well, what is left that should be dreaded?

\ciceroSection{106} Exile, I suppose — which is reckoned among the greatest evils. But if it is an evil on account of the hostility and ill will of the people toward another man, how worthy of contempt that is has just been said a moment ago. And if it is wretched to be away from one's country, the provinces are full of wretched men, of whom very few ever return home.

\ciceroSection{107} "But exiles are stripped of their goods." What of it? Has not enough been said about enduring poverty? And as for exile itself, if we look to the nature of the thing and not to the disgrace of the name, how much, after all, does it differ from a perpetual residence abroad? In such a life the most distinguished philosophers spent their days — Xenocrates, Crantor, Arcesilas, Lacydes, Aristotle, Theophrastus, Zeno, Cleanthes, Chrysippus, Antipater, Carneades, Clitomachus, Philo, Antiochus, Panaetius, Posidonius, and countless others, who once they had set out never returned home. But exile carries disgrace, you say. Can exile touch the wise man with disgrace? For this whole discourse is about the wise man, on whom it cannot justly fall; and there is no need to console one who is justly exiled.

\ciceroSection{108} Lastly, those who refer to pleasure everything that comes about in life have the readiest reckoning for every circumstance, since wherever pleasure is supplied, there they can live happily. And so the saying of Teucer can be adapted to every situation: \textit{My country is wherever it goes well with me.} Socrates indeed, when he was asked what country he claimed as his own, said, "The world"; for he held himself to be an inhabitant and a citizen of the whole world. And what of Titus Albucius? Did he not philosophize at Athens, an exile, with the most untroubled spirit? Yet even that would not have befallen him, had he kept quiet in public life and obeyed the precepts of Epicurus.

\ciceroSection{109} For how was Epicurus the happier because he lived in his own country, than Metrodorus because he lived at Athens? Or did Plato surpass Xenocrates, or Polemo Arcesilas, in being any happier for it? And how much is that state to be valued from which good and wise men are driven out? Damaratus, indeed, the father of our king Tarquin, because he could not endure the tyrant Cypselus, fled from Corinth to Tarquinii, and there established his fortunes and begot children. Did he choose foolishly when he preferred the freedom of exile to servitude at home?

\ciceroSection{110} Furthermore, the agitations of the mind, its anxieties and its griefs, are soothed by forgetfulness when minds are drawn off toward pleasure. Not without reason, then, did Epicurus venture to say that the wise man is always in the midst of more goods, because he is always in the midst of pleasures. From which he thinks it follows — the very thing we are after — that the wise man is always happy.

\ciceroSection{111} "Even if he is robbed of the use of his eyes, or his ears?" Even then; for those very things he despises. For, first of all, of what pleasures, after all, does that dreaded blindness deprive him? Some even argue that the other pleasures reside in the senses themselves, but that the things perceived by sight have no dealings with any delight of the eyes; so that what we taste, smell, handle, and hear are dealt with in that very part where we perceive them. With the eyes nothing of the sort happens: the mind takes in what we see. But the mind can be delighted in many and various ways, even when sight is not brought to bear. For I am speaking of a learned and cultivated man, for whom to live is to think. And the wise man's thought does not, as a rule, call in the eyes as helpers for its inquiries.

\ciceroSection{112} For indeed, if night does not take away the happy life, why should a day like night take it away? For that saying of Antipater of Cyrene is a little too coarse, but the thought is not absurd: when some women were lamenting his blindness, he said, "What are you about? Does no pleasure seem to you to belong to the night?" That old Appius, too, who was blind for many years — we understand from his offices and his achievements that in that misfortune of his he failed in no duty either private or public. We have it on record that the house of Gaius Drusus used to be thronged with men seeking his counsel; when those whose business it was could not see their own affairs, they took a blind man for their guide. When we were boys, Gnaeus Aufidius, a former praetor, used to deliver his opinion in the Senate, was never wanting to friends seeking his advice, wrote a history in Greek, and saw with his letters.

\ciceroSection{113} Diodotus the Stoic lived in our house for many years, blind. He, indeed — a thing scarcely credible — applied himself to philosophy far more assiduously even than before, played the lyre after the manner of the Pythagoreans, had books read to him night and day, and for these studies had no need of eyes; and, what seems scarcely possible to be done without eyes, he kept up the work of geometry, instructing his pupils in words from where to where they should draw each line. They tell of Asclepiades, no obscure philosopher of Eretria, that when someone asked what blindness had brought him, he answered: that he went about with one boy more in his train. For just as even the extreme of poverty would be bearable, if a man were allowed what some Greeks earn day by day, so blindness could be borne easily, if the supports for one's infirmities were not lacking.

\ciceroSection{114} Democritus, having lost his sight, could not, to be sure, tell white from black; but good from bad, just from unjust, honorable from base, useful from useless, great from small, that he could; and without the variety of colors it was possible to live happily, but without the knowledge of things it was not. And this man held that the mind's keen edge is even hindered by the sight of the eyes; and while others often could not see what lay before their feet, he ranged abroad through all infinity, halting at no farthest boundary. It is handed down too that Homer was blind; yet we look upon his painting, not his poetry. What region, what coast, what spot of Greece, what shape and form of battle, what battle-line, what fleet, what movements of men and of beasts has he not so vividly depicted that he made us see the very things he himself did not see? What then? Are we to suppose that Homer, or any learned man, was ever wanting in delight and pleasure of the mind?

\ciceroSection{115} Or, were the matter not so, would Anaxagoras, or this same Democritus, have abandoned their fields and their patrimonies, and given themselves wholeheartedly to this divine delight of learning and inquiry? And so the augur Tiresias, whom the poets imagine as a wise man, they never bring on bewailing his blindness; whereas Homer, having made Polyphemus monstrous and savage, has him even converse with his ram and praise its good fortune, in that it could go where it pleased and reach what it pleased. And rightly so; for the Cyclops himself was no wiser than that ram.

\ciceroSection{116} In deafness, again, what evil is there, after all? Marcus Crassus was rather hard of hearing, but there was something more vexing to him — that he heard ill spoken of, even if, as it seemed to me, unjustly. Our Epicureans, for the most part, know no Greek, nor the Greeks Latin. These men, then, are deaf to the speech of those, and those to the speech of these; and likewise all of us are surely deaf to those tongues we do not understand, which are beyond counting. But the deaf do not hear the voice of the lyre-player. Neither do they hear the screech of the saw when it is being sharpened, nor the grunting of a pig when its throat is cut, nor, when they wish to rest, the roar of the murmuring sea; and if songs perhaps delight them, they ought first to reflect that, before these were invented, many wise men lived happily, and then, that far greater pleasure can be drawn from reading such things than from hearing them.

\ciceroSection{117} Then, just as a little while ago we drew the blind across to the pleasure of the ears, so we may draw the deaf across to that of the eyes. For indeed, the man who can converse with himself will not require another's talk. But let everything be heaped together at once, so that the same man is bereft of both eyes and ears, and let him be crushed besides by the sharpest pains of the body. These, in the first place, generally make an end of a man by themselves; but if perhaps, drawn out by their long duration, they torment him more violently than there is any reason why he should bear them — what is there, after all, good gods, that we should trouble ourselves over? For the harbor is at hand, since death is the same haven there, the eternal refuge of feeling nothing.

\ciceroSection{118} Theodorus, when Lysimachus threatened him with death, said, "A great thing indeed you have accomplished, if you have come by the power of a blister-beetle." Paulus, when Perseus begged not to be led in the triumph, said, "That, at least, is in your own power." Much was said on the first day, when we were inquiring into death itself, and not a little on the day after as well, when pain was the subject, about death; and whoever calls these things to mind can hardly be in any danger of not judging death either something to be wished for or, at the very least, not to be feared. To me, indeed, that rule which holds at the banquets of the Greeks seems one to keep in life: "Let him drink," it says, "or let him go." And rightly. For either let a man enjoy the pleasure of drinking equally with the rest, or let him withdraw beforehand, lest, while sober, he fall foul of the violence of the drunk. In the same way you may, by flight, leave behind the wrongs of fortune which you cannot bear. These same things that Epicurus says, Hieronymus says in just so many words.

\ciceroSection{119} But if those philosophers whose doctrine it is that virtue of itself has no power, and that everything we call honorable and praiseworthy they call empty and decked out with a hollow sound of words — if even they nonetheless judge the wise man always happy, what then ought to be thought right for the philosophers who set out from Socrates and Plato? Of these, some say there is so great a preeminence in the goods of the mind that the goods of the body and the external things are overwhelmed by them; while others do not even count these as goods at all, but lay everything up in the mind.

\ciceroSection{120} It was the controversy between these that Carneades, like an honorary arbiter, was wont to decide. For since whatever goods the Peripatetics held seemed to the Stoics mere advantages, and yet the Peripatetics granted no more to riches, to good health, and to the rest of that kind than the Stoics did — when these things were weighed by their substance and not by words, he denied that there was any ground for the dispute. As for this point, how the philosophers of the other schools may maintain it, let them see to it themselves; for my part, I am glad that, concerning the wise man's unbroken capacity for living well, they profess something worthy of a philosopher's voice.

\ciceroSection{121} But since we must set out in the morning, let us gather up in memory these discussions of five days. For my part, I think I shall even commit them to writing — for where could we better use this leisure, of whatever sort it is? — and we shall send these second five books to our friend Brutus, by whom I was not only urged on to the writing of philosophy, but even challenged to it. In this work, how much good we shall do the rest of the world I could not easily say; but for my own most bitter griefs, and the manifold troubles besetting me on every side, no other relief could be found.
